\section{Surjectivity II: geometric pages and exponent bounds}\label{sec:pf:onto2}

\providecommand{\ontoIIix}{\operatorname{ix}}
\providecommand{\ontoIITw}{\mathrm{Tw}}
\providecommand{\ontoIIcut}{\operatorname{cut}}
\providecommand{\ontoIIeig}{\operatorname{eig}}
\providecommand{\ontoIIexp}{\operatorname{ex}}
\providecommand{\ontoIIJA}{J^{\mathrm{A}}}
\providecommand{\ontoIIJr}{J^{\mathrm{rec}}}
\providecommand{\ontoIIJu}{J^{\mathrm{u}}}
\providecommand{\ontoIIloc}[2]{{#2}_{(#1)}}
\providecommand{\ontoIIRB}{\mathrm{RB}}
\providecommand{\ontoIIPN}{\mathrm{PN}}
\providecommand{\ontoIICl}{\mathrm{Cl}}
\providecommand{\ontoIILI}{\mathrm{LI}}

This section proves two of the statements about \emph{record objects} from which the surjectivity estimate
$\sfname{SurjectiveOntoX}$ of \cref{prop:families:points} is assembled
(\lref{GlobalTheory.BigOn.ontoX_of_twist_shares_exp_alt}):
\begin{enumerate}[label=(\Alph*)]
\item the \emph{page statement} implies the \emph{geometric statement} (\cref{pf:onto2:thmA});
\item eight published results (Clozel's purity theorem, the $\ell$-adic characters attached to algebraic Hecke
  characters and their behaviour above $\ell$, Chebotarev's theorem with degree-one primes, Galois representations of
  polarised cuspidal representations, base change for definite unitary groups, the Bhatt--Scholze theorem on
  crystalline lattices, and the existence of local Langlands parameters) imply the \emph{exponent statement}
  (\cref{pf:onto2:thmB}); the first of them is not used (\cref{pf:onto2:rem:unused}).
\end{enumerate}
All three statements are \emph{closed statements over record objects} (\cref{pf:onto2:forall}): for every family of
the class $\sfname{FamHyp}$ which is large in the sense of the quantifier $\sfname{Closed}$ (\cref{def:families:closed}),
for all large constants and all large scales $x$, they assert a list of explicit conditions (``leaves'') for every
record object $O$ attached to a point $p$ of the family and to a pair of runs. The geometric statement consists of
lattice leaves (the lattices attached to $O$ are lattices, stable under the relevant finite group, and suitably
nested), a local index bound, and four \emph{depth leaves} bounding the normalised exponents $e_0(w,j)$ of $O$ at
twisted places $w$ by the content $F_0(w,j)$ of the point plus explicit error terms, separately at good primes, at
the ramified primes $q$ and $d$ of the host, at the small primes, and at the two level primes. The exponent statement
asserts that the exponents $J_0$ of the numbers of the two-sided reduction are non-negative at twisted places and are
dominated by a function $J_p(j,\ell)$ of the residue characteristic whose weighted sums
$\sum_{\ell\in Q}J_p(j,\ell)\log\ell$ are bounded independently of the finite set of primes $Q$.

\paragraph{The idea.} The geometric statement is a reduction to three inputs proved elsewhere: the index-bound
hypotheses along runs (\lref{GlobalTheory.FinalJoin.runHolds_of_parOK}), the identification of the lattices of
$O$ with lattices of modular forms (the \emph{form ties} contained in the record relation), and local depth bounds for
integral points of the relevant moduli spaces. The two pieces of lattice theory proved here in full are a
local--global passage (\cref{pf:onto2:l5}), which converts a bound for the index of a theta lattice at the places of a
number field $K_0$ into local index bounds at the places of an extension $K$, and the averaging argument for the
eigen-part of a lattice (\cref{pf:onto2:avg}). The local depth bounds at good primes rest on \emph{depth packages} extracted
from the jets of integral charts (\cref{pf:onto2:packages,pf:onto2:nugood}). The exponent statement compares, at a twisted place $w$ of residue
characteristic $\ell$, the normalised $\ell$-adic exponent $\ontoIIJu$ of the inclusion of a ``Weyl-type'' block lattice
into a ``dual-Weyl-type'' one with an explicit integer $\ontoIIJr(\ell)$ built from exponents of Weyl modules. Outside an
explicit range of $\ell$ the comparison follows from an inclusion whose exponent is the sum of the Weyl exponents of the weights
(\cref{pf:onto2:sumexp}); in the range, the coarse bound defining $\ontoIIJr(\ell)$ rests
on \emph{local models}: Hida-type settings over $\ell$-adic integers covering the block, in which a lattice inclusion with the coarse
exponent holds (\cref{pf:onto2:lincl}); their inputs are obtained from the Galois representations attached to the automorphic
eigensystems of the block via base change, together with Chebotarev's theorem and the Bhatt--Scholze theorem
(\cref{pf:onto2:latlink,pf:onto2:winputs,pf:onto2:model}). Summing over $\ell$ uses the cut inequality for Weyl
modules and the prime number bound $\pi(y)<1.26\,y/\log y$ of Rosser and Schoenfeld.

\paragraph{Organisation.} \Cref{sec:pf:onto2:setup} recalls record objects, the leaves and the three statements.
\Cref{sec:pf:onto2:geo} proves (A): runs (\cref{sec:pf:onto2:runs}), lattice leaves (\cref{sec:pf:onto2:lat}), the
index bound (\cref{sec:pf:onto2:index}), the depth leaves (\cref{sec:pf:onto2:depth}) and the assembly
(\cref{sec:pf:onto2:geoproof}). \Cref{sec:pf:onto2:exp} proves (B): weight dictionaries (\cref{sec:pf:onto2:dict}), the
reduction to four statements (\cref{sec:pf:onto2:red}), the numerical and exponent-sum statements
(\cref{sec:pf:onto2:sum}), and the smallness, frame, Galois and model statements (\cref{sec:pf:onto2:four}).
\Cref{sec:pf:onto2:remarks} collects remarks on the formal statements. Lemmas of the formal development whose proofs
are not reproduced here are cited by their labels together with their statements.

%======================================================================
\subsection{Record objects, leaves and the three statements}\label{sec:pf:onto2:setup}

\paragraph{Conventions.} Throughout, $K$ is a number field with an involution $a\mapsto\bar a$ (in the applications
$K$ is the CM coefficient field $\mathrm{LL}$ of a datum $D$ and the involution is complex conjugation). For a finite
place $w$ of $K$ (a non-zero prime of $\OO_K$) we write $p_w$ for its residue characteristic, $v_w$ for the normalised
valuation (with the convention $v_w(0)=0$ of the formal development), $e_w=v_w(p_w)$ for the absolute ramification
index, $\ell_w=\log\Nm w/[K:\Q]$, and $\bar w$ for the conjugate place. For an $\OO_K$-submodule $A$ of a $K$-vector
space $Z$, $\ontoIIloc{w}{A}$ denotes the localisation of $A$ at $w$ (the $\OO_{K,(w)}$-span of $A$), and for two such
submodules
\begin{equation}\label{eq:pf:onto2:ix}
\ontoIIix_w(A,B)=\len\bigl(\ontoIIloc{w}{A}/\ontoIIloc{w}{(A\cap B)}\bigr)-\len\bigl(\ontoIIloc{w}{B}/\ontoIIloc{w}{(A\cap B)}\bigr)\in\Z ,
\end{equation}
lengths being lengths of $\OO_K$-modules, with the convention that a module of infinite length contributes $0$
(\cref{pf:onto2:rem:len}). For full lattices this is the $w$-adic valuation of the generalised index $[A:B]$. For a
$K$-subspace $W\subseteq Z$, $A$ is a \emph{lattice in $W$} if $A$ is finitely generated and spans $W$; we write
$\ontoIIcut_W(A)=A\cap W$ and $aA$ for the image of $A$ under multiplication by $a\in K$. For a representation $\rho$ of a
group $P$ on $Z$, a subgroup $U\le P$ and a function $\psi:P\to\OO_K$, we put
$\ontoIIeig(\rho,U,\psi)=\{z\in Z:\rho(u)z=\psi(u)z\ \text{for all }u\in U\}$.

An \emph{indefinite datum} $S$ (Part~I, \cref{sec:families:nonempty}) carries in particular a finite set $\ontoIITw(S)$ of
finite places of $K$ (the \emph{twisted places}), a type $\mu(w)$ at every finite place, an integer $D(S)$ and an index
set $J$; its \emph{numbers} $N$ contain real functions $F_0(w,j)$ (the \emph{content}), $J_0(t,w,j)$ (the
\emph{exponents}) and $E_1(j)$.

\begin{definition}[Record objects]\label{pf:onto2:onobj}
Let $S$ be an indefinite datum over $K$ with index set $J$, let $P$ be a finite group and $c=(C,C_{b1},C_{b2})$ a
triple of reals. A \emph{record object} $O\in\mathcal O(S,P,c)$ (\lref{Formalism.L2.Indef.Tree.OnObj}) consists
of the following data, of which we list those used in this section. For each $j\in J$: an integer $n_Z(j)$, a
representation $\rho_j$ of $P$ on $K^{n_Z(j)}$, an $\OO_K$-submodule $M_s(j)\subseteq K^{n_Z(j)}$, and an element
$\mathrm{nn}_K(j)\in K$; a subgroup $U\le P$ and a function $\psi:P\to\OO_K$; for each finite place $v$ and $j\in J$:
$\OO_K$-submodules $T_{h1}(v,j)$, $T_h(v,j)$, $T_{hK}(v,j)$, $T_{h0}(v,j)$ of $K^{n_Z(j)}$, elements
$\sigma(v,j),a(v,j)\in K$, an integer $E(v,j)\ge0$ and reals $r_1(v,j),r_2(v,j)$; a prime $p_{\mathrm{de}}$; real
functions $J_p(j,\ell)$, $\theta(\ell)$, $\alpha(\ell)$ ($\ell\in\N$) and $\mathrm{wt}(j)$; reals
$k,m,f,b',S_c,\eta,T,D_F,X_2,x,\varepsilon,v_d,r_r,\delta_{e2},N_c$; and natural numbers $T_2,g_1,g_2,q_1,d$ and a finite
set $s_m$ of natural numbers. Derived quantities are
\begin{gather*}
e_0(v,j)=\frac{E(v,j)-v_v(\sigma(v,j))}{e_v},\qquad \mathrm{nn}(v,j)=\frac{v_v(\mathrm{nn}_K(j))}{e_v},\\
x_e=x^{\varepsilon},\quad b=e^{x_e},\quad L=\log b,\quad Y=C D_F,\qquad
r(v,j)=\begin{cases} r_2(v,j)&p_v=p_{\mathrm{de}},\\ r_1(v,j)&\text{otherwise},\end{cases}\\
S_e=S_c+\eta,\qquad s_B(j)=\mathrm{wt}(j)+fS_e,\qquad s_C(j)=Cm(2+v_d)\,s_B(j),\\
A_1(j)=C\bigl(\log k\cdot m\,\mathrm{wt}(j)+fkS_e\bigr),\qquad A_{kw}(j)=Ck\,\mathrm{wt}(j),\\
a_E(j,\ell)=\frac{A_1(j)}{\ell}+A_{kw}(j)\theta(\ell)+T\alpha(\ell)+[X_2<\ell]\frac{A_{kw}(j)}{\ell^2},\\
E_J=30\,T_2\,m\,b\,x_e^2+2mb,\qquad E_{1u}=x_e(fmbL+N_c),
\end{gather*}
where $[\cdot]$ is $1$ if the condition holds and $0$ otherwise.
\end{definition}

\begin{definition}[Leaves]\label{pf:onto2:leaves}
Let $O\in\mathcal O(S,P,c)$ and let $N$ be numbers for $S$. Write ``at pairs'' for ``for every $w\in\ontoIITw(S)$ and both
$v\in\{w,\bar w\}$''. The \emph{leaves} used here are the following conditions (\lref{Formalism.L2.Indef.Tree.OnObj}).
\begin{enumerate}[label=(L\arabic*)]
\item\label{pf:onto2:L1} \emph{Stable lattices}: for all $j$, $M_s(j)$ is a lattice in $K^{n_Z(j)}$ and $\rho_j(g)M_s(j)\subseteq
  M_s(j)$ for all $g\in P$.
\item\label{pf:onto2:L2} \emph{Theta lattices}: at pairs, $T_h(v,j)$ and $T_{hK}(v,j)$ are lattices in $K^{n_Z(j)}$ for all $j$.
\item\label{pf:onto2:L3} \emph{Nesting}: at pairs, for all $j$: $T_h(v,j)\subseteq M_s(j)$, $T_h(v,j)$ is $\rho_j(P)$-stable,
  $\sigma(v,j)\ne0$ and $T_h(v,j)\subseteq\sigma(v,j)\,T_{hK}(v,j)$.
\item\label{pf:onto2:L4} \emph{Eigen-part}: at pairs, for all $j$: $T_{h0}(v,j)$ is finitely generated,
  $T_{h0}(v,j)\subseteq\ontoIIeig(\rho_j,U,\psi)$, and $|U|\cdot\bigl(T_{hK}(v,j)\cap\ontoIIeig(\rho_j,U,\psi)\bigr)\subseteq
  T_{h0}(v,j)$.
\item\label{pf:onto2:L5} \emph{Index bound}: there are $T_{h1},a,r_1,r_2,\delta_{e2}$ (replacing the corresponding data of $O$)
  such that: at pairs, for all $j$, $T_{h1}(v,j)$ is a lattice contained in $T_h(v,j)$, $a(v,j)\ne0$, $v_v(a(v,j))=E(v,j)$, and
  for every $K$-subspace $W\subseteq K^{n_Z(j)}$
  \begin{equation}\label{eq:pf:onto2:LIB}
  \ontoIIix_v\bigl(\ontoIIcut_W M_s(j),\ontoIIcut_W T_{h1}(v,j)\bigr)\le E(v,j)\dim W+r(v,j);
  \end{equation}
  moreover $0\le\delta_{e2}\le1$, $r_1,r_2\ge0$, and for every $j$ and every finite set $S'$ of finite places,
  $\sum_{v\in S'}\ell_v r_i(v,j)\le n_Z(j)\,fmb\,\delta_{e2}$ for $i=1,2$.
\item\label{pf:onto2:L6} \emph{Depth at good primes} (for $N$): $\theta(\ell),\alpha(\ell)\ge0$ for all $\ell$;
  $\sum_{p\in Q}\theta(p)\log p\le C/k$ and $T\sum_{p\in Q}\alpha(p)\log p\le Cmb$ for every finite set $Q$ of primes; and for
  every $w\in\ontoIITw(S)$ with $p_w\notin s_m\cup\{q_1,d,g_1,g_2\}$ and every $j$: if $p_w\le Y$ then $e_0(v,j)\le
  F_0(w,j)+a_E(j,p_w)+1$ for $v\in\{w,\bar w\}$, and if $p_w>Y$ then $e_0(w,j)=e_0(\bar w,j)=0\le F_0(w,j)$.
\item\label{pf:onto2:L7} \emph{Depth at $q_1$ and $d$}: for $w\in\ontoIITw(S)$ with $p_w\in\{q_1,d\}$ and all $j$:
  $e_0(v,j)\le s_B(j)+T+1$ for $v\in\{w,\bar w\}$; $F_0(w,j)=0$ if $p_w=q_1$; and $-2fkb'\le F_0(w,j)\le0$ if $p_w=d$.
\item\label{pf:onto2:L8} \emph{Depth at small primes}: for $w\in\ontoIITw(S)$ with $p_w\in s_m$ and all $j$: $e_0(v,j)\le
  s_C(j)+T+1$ for $v\in\{w,\bar w\}$, and $F_0(w,j)=0$.
\item\label{pf:onto2:L9} \emph{Depth at level primes}: for $w\in\ontoIITw(S)$ with $p_w\in\{g_1,g_2\}$ and all $j$:
  $e_0(v,j)\le 2T-\mathrm{nn}(v,j)$ for $v\in\{w,\bar w\}$, and $\mathrm{nn}(w,j)+\mathrm{nn}(\bar w,j)=-2F_0(w,j)$.
\item\label{pf:onto2:L10} \emph{Exponent bound}: $J_0(\mu(w),w,j)\ge0$ for $w\in\ontoIITw(S)$ and all $j$, and there is a
  function $J_p(j,\ell)\ge0$ (replacing that of $O$) such that $\sum_{p\in Q}J_p(j,p)\log p\le E_J$ for every finite set $Q$
  of primes and every $j$, and $J_0(\mu(w),w,j)\le J_p(j,p_w)$ for $w\in\ontoIITw(S)$ and all $j$.
\end{enumerate}
The \emph{page leaves} are:
\begin{enumerate}[label=(P\arabic*)]
\item\label{pf:onto2:P1} $C,k,m\ge1$, $f\ge8$, $T,S_c,\eta,b',v_d,N_c\ge0$, $x\ge1$; $g_1,g_2,q_1,d$ are primes with $g_1\ne
  g_2$, $q_1\ne d$, $\{q_1,d\}\cap\{g_1,g_2\}=\emptyset$, and none of them lies in $s_m$;
\item\label{pf:onto2:P2} $0\le\mathrm{wt}(j)\le2fb$, $1<Y$, $\log Y\le 2L$, $Y\le mb$, $kS_eL\le mb$, $k\le X_2$, $\log k\le L$,
  $1\le L$;
\item\label{pf:onto2:P3} $\log x\le x_e$ and $T_2\le f$;
\item\label{pf:onto2:P4} three inequalities bounding $\log(q_1d)$, $\sum_{p\in s_m}\log p$ and $\log(g_1g_2)$ against
  $C_{b1}fmb$, $C_{b2}fmbL$ and $4b$, and $C_{b1},C_{b2}\ge0$;
\item\label{pf:onto2:P5} a condition on the residue fields of the places of $R.P_g$ (each has $g_1$ or $g_2$ elements, $f$ of each);
\item\label{pf:onto2:P6} $2f/(\min(g_1,g_2)-1)\le 1/(4(fk)^3)$.
\end{enumerate}
\end{definition}

\begin{definition}[Closed statements over record objects]\label{pf:onto2:forall}
For a point $p$ (\cref{sec:families:nonempty}) let $\mathcal M(p)$ denote the \emph{extra point conditions}: every place at
which the complex of $p$ has a single index is twisted, the core of $p$ is tied to the two-sided reduction $R_2(p)$ (the
\emph{core ties}), and the record ties of $p$ hold (\lref{GlobalTheory.BigOn.More}). For a pair of runs $r$ for the
core of $p$ (\lref{GlobalTheory.BigOn.TwoRuns}), a \emph{record object of $p$ along $r$} is a record object $O$ for
the false side $S_2(p)$ of $R_2(p)$, the group $\mathrm{Mir}(R_2(p))$ and constants $c$, which satisfies the record relation
(\cref{pf:onto2:isrec}) for the parameters of the family and the own-prime pin of $r$. Let $Q$ be a property of a family
$\mathcal F$, a scale $x$, an index $i$, a point $p\in\mathrm{Pt}(\mathcal F,x,i)$, constants $c$ and a record object $O$. Then
$\mathrm{Rec}(Q)$ (\lref{GlobalTheory.BigOn.ForallRecordObjects}) is the statement
\begin{multline*}
\sfname{Closed}\Bigl(\sfname{FamHyp},\ \exists c_0\ \forall c:\ \bigl(c_0\le C\wedge 14\le C_{b1}\wedge
66C+4\le C_{b2}\bigr)\Rightarrow{}\\ \mathrm{Ev}\bigl(\mathcal F,\ p\mapsto \mathcal M(p)\Rightarrow\forall r\ \forall O\
\text{record object of $p$ along $r$}:\ Q(\mathcal F,x,i,p,c,O)\bigr)\Bigr),
\end{multline*}
where $\mathrm{Ev}(\mathcal F,\cdot)$ means ``there is $x_0$ such that for all $x>x_1(\mathcal F)$ with $x\ge x_0$, all $i$ and all
points $p$ over $(x,i)$ satisfying the point conditions''. Similarly $\mathrm{AtPts}(Q')$
(\lref{GlobalTheory.BigOn.AtPtsAlt}) is $\sfname{Closed}(\sfname{FamHyp},\mathrm{Ev}(\mathcal F,p\mapsto\mathcal
M(p)\Rightarrow Q'(p)))$. For a point $p$ and reals $C,C_H$, $N_p(C,C_H)$ denotes the numbers of the two-sided reduction
(\lref{GlobalTheory.BigOn.R2}).
\begin{itemize}
\item The \emph{page statement} is $\mathrm{Rec}$ of: the page leaves, $O.b=b(L)$ (where $b(L)$ is the scale $b$ of the size data
  $L=L(x,i)$), $p.Z.E_{1u}\ge0$, and $E_1(N_p(C,C_H))(j)=C\cdot O.E_{1u}$ for all $C,C_H,j$
  (\lref{GlobalTheory.BigOn.PageStatement}).
\item The \emph{geometric statement} is $\mathrm{Rec}$ of: \ref{pf:onto2:L1}--\ref{pf:onto2:L5}, and
  \ref{pf:onto2:L6}--\ref{pf:onto2:L9} for $N=N_p(C,C_H)$ for all $C,C_H$ (\lref{GlobalTheory.BigOn.GeometricStatement'}).
\item The \emph{exponent statement} is $\mathrm{Rec}$ of: \ref{pf:onto2:L10} for $N=N_p(C,C_H)$ for all $C,C_H$
  (\lref{GlobalTheory.BigOn.ExponentStatement}).
\end{itemize}
\end{definition}

\begin{remark}[The record relation]\label{pf:onto2:isrec}
The record relation (\lref{GlobalTheory.BigOn.IsRec}) ties a record object to the runs $r$ and to the parameters of the
family. Its components are: (i) the \emph{run conditions} (\lref{GlobalTheory.BigOn.RunsOk}): the top structures of both
runs have all their laws, one-cover conditions hold, the twist character is trivial on determinants of the rational subgroup,
run, level and action pins hold at the twisted places, and the level ideal is the product of the places in $R.P_g$; (ii) the
weight families $\kappa_i$ of the face are the run weights $\kappa R(i+1)$; (iii) a \emph{level compatibility}: the core has
its laws, the map $\OO\to\mathrm{LL}$ is injective, and for every $j$ there are an isomorphism $z_j:K^{n_Z(j)}\cong
K\otimes\mathrm{Cusp}_j$ with a space of cusp forms (of weight $\kappa R(j+1)$) and an action $\mathrm{act}_j$ of $P$ on these
cusp forms, and for every place $v$ a further identification, such that the \emph{trace tie} holds at $(v,j)$ for the run
$r_v$ \emph{used at $v$} (the second run of $r$ if $p_v=p_{\mathrm{de}}$, the first otherwise) and the prime $\mathrm{own}(v)$
($p'_{\mathrm{de}}$ if $p_v=p_{\mathrm{de}}$, $p_{\mathrm{de}}$ otherwise); the trace tie (\lref{GlobalTheory.OnTies.OnTie})
contains the \emph{form tie} of \cref{pf:onto2:formtie} and identifies $U$, $\psi$, $T_{h0}$ and the remaining lattices of $O$
with objects defined from theta series; (iv) the \emph{depth tie} of \cref{pf:onto2:runtie}; (v) the identities
$D_F=\max_j\deg(lR(j+1))$, $r_r=D.R.r$, $N_c=L.N_c$, $T_2=|D.\mathrm{host}.T_2|$ and $p_{\mathrm{de}}=D.\mathrm{host}.p_{\mathrm{de}}$.
The \emph{own-prime pin} of $r$ (\lref{GlobalTheory.BigOn.OwnPin}) says that all auxiliary prime data of the run used at
$v$ are located at $\mathrm{own}(v)$.
\end{remark}
%======================================================================
\subsection{The geometric statement}\label{sec:pf:onto2:geo}

\begin{theorem}[Page statement implies geometric statement]\label{pf:onto2:thmA}
The page statement implies the geometric statement. In words: if, closedly over $\sfname{FamHyp}$, for large constants and
eventually at points, every record object satisfies the page leaves, has $O.b=b(L)$ and $E_1=C\cdot O.E_{1u}$, then
closedly over $\sfname{FamHyp}$, for large constants and eventually at points, every record object satisfies the lattice leaves
\ref{pf:onto2:L1}--\ref{pf:onto2:L4}, the index bound \ref{pf:onto2:L5} and the depth leaves \ref{pf:onto2:L6}--\ref{pf:onto2:L9}
for the numbers $N_p(C,C_H)$, for all $C,C_H$.
\formallabel{GlobalTheory.BigOn.geoSt_of_pageSt_alt}
\end{theorem}

The proof (\cref{sec:pf:onto2:geoproof}) combines four results about a single record object, proved in
\cref{sec:pf:onto2:runs,sec:pf:onto2:lat,sec:pf:onto2:index,sec:pf:onto2:depth}, with numerical conditions along families. In
these four results we fix a datum $D$, a finite group $G$, a family $A$ of forms on $G$, a core $J$ with a real datum $R$, a
face $F$, weight families $\kappa$ with vectors $v_0$ and a tie $\mathrm{fw}$, size data $L$, parameter choices with parameters
$T_p$, the weight definition $\kappa_{\mathrm{def}}$ of $D$, an error datum $\mathrm{EX}$, a pair of runs $r$ for $J$ and
constants $c$; $O$ is a record object for the false side of the complex of $R$ at $F$, the group $\mathrm{Mir}(R)$ and $c$; and
$K=\mathrm{LL}$ is the coefficient field of $D$.

\subsubsection{Index-bound hypotheses along runs}\label{sec:pf:onto2:runs}

\begin{definition}[Run bounds and page-run conditions]\label{pf:onto2:runbounds}
For a pair of runs $r$ and $C_K\in\R$, the \emph{run bounds} $\ontoIIRB(C_K,r)$ say: for every place $v$ such that the top
structure of the run $r_v$ used at $v$ satisfies its laws, if $\OO\to\mathrm{LL}$ is injective, then for every
admissible index $j$ the \emph{index-bound hypotheses} and the \emph{run conclusion}
$\ontoIICl(X,S,\kappa,w_L,w_F,T,\nu_s,\beta)$ (\cref{pf:onto2:concl}) hold for the record $X$ of the run at $j$
(\cref{pf:onto2:formtie}), the scalar datum $S$ of the run at the prime $\mathrm{own}(v)$, the run weights
$(\kappa,w_L,w_F)=(\kappa R(j),lR(j),fR(j))$, the integer $T=T_O$ of the parameters, the local depths $\nu_s$ of the run, and
\begin{equation}\label{eq:pf:onto2:beta}
\beta=\rank_{K_0}\bigl(K_0\cdot\mathrm{Mlat}(\kappa)\bigr)\cdot f\,m\,b\,\delta_{e2},
\end{equation}
where $K_0$ is the coefficient field of the record, $f,m$ are ranks of $D$, $b$ is the scale of $L$ and $\delta_{e2}$ the thick
density of $L$. The \emph{page-run conditions} $\ontoIIPN(x_1)$ (\lref{GlobalTheory.BigOn.PageNum}) are: the size conditions
on $L$, $C_{nW}=10$, $\log g_i\le 9L+1$ for the two level primes, $1000\le C_P$, $C_P^{20}\le z_{\mathrm{Base}}(L)$ and $x_1\le x$.
\formallabel{GlobalTheory.BigOn.RunsHold}
\end{definition}

\begin{lemma}[Index-bound hypotheses at a prime]\label{pf:onto2:hypconcl}
There are constants $C_K$ and $x_1$ such that for every top structure $T$ satisfying its laws with $\OO\to\mathrm{LL}$ injective,
every error datum $\mathrm{EX}$ and every prime $p$ for which the parameter conditions of $(T,\mathrm{EX},p)$ hold at all
admissible indices: if
$\ontoIIPN(x_1)$ holds and the field smallness of $(D,L)$ holds, then the index-bound hypotheses and the run conclusion (with
the data of \cref{pf:onto2:runbounds} for the prime $p$) hold at every admissible index $j$.
\formallabel{GlobalTheory.FinalJoin2.hyp_concl_at_run}
\formallabel{GlobalTheory.T.exists_constants_indexBound_hypotheses_liftWeights}
\end{lemma}

\begin{proof}
By \lref{GlobalTheory.FinalJoin.runHolds_of_parOK} there is a threshold function $x_0(\cdot,\cdot)$ such that for every top
structure $T$ with its laws, every error datum and prime $p$ for which the parameter conditions of $(T,\mathrm{EX},p)$ hold at all
admissible indices, and every record run satisfying the page-run hypotheses at a scale $x\ge x_0(20,C_K)$, both the
index-bound hypotheses and the run conclusion hold for the record at $j$. By
\lref{GlobalTheory.FinalJoin.RecordRun.exists_pageRun_record_sub_four} there is a constant $C_K$ such that the page-run
hypotheses hold whenever the size conditions, $C_{nW}=10$, $\log g_i\le9L+1$, $x\ge10^{60}$, $C_P\ge1000$, $C_P^{20}\le
z_{\mathrm{Base}}(L)$ and a smallness condition on the coefficient field hold; by
\lref{GlobalTheory.FinalJoin2.fieldSmall_iff_opo} this smallness condition is equivalent to the field smallness of $(D,L)$.
Put $x_1=\max(10^{60},x_0(20,C_K))$; then $x\ge x_1$ gives both scale conditions.
\end{proof}

\begin{proposition}[Runs satisfy the index-bound hypotheses]\label{pf:onto2:runshold}
There are constants $C_K$ and $x_1$ such that for all data as above: if the core ties of $(L,\mathrm{EX},J,R)$, the run
conditions, the own-prime pin and the page-run conditions $\ontoIIPN(x_1)$ hold, then the run bounds $\ontoIIRB(C_K,r)$ hold.
\formallabel{GlobalTheory.BigOn.runsHold_of_ties}
\end{proposition}

\begin{proof}
Let $C_K,x_1$ be the constants of \cref{pf:onto2:hypconcl}. Let $v$ be a place and apply \cref{pf:onto2:hypconcl} to the run
$T=r_v$ and $p=\mathrm{own}(v)$. Field smallness is one of the core
ties, and the laws of $T$ are part of the run conditions. It remains to check the parameter conditions for
$(T,\mathrm{EX},\mathrm{own}(v))$. By \lref{GlobalTheory.FinalJoin2.liftIndexParameterConditions_of_coreAndShortLaws} they
hold at all $j$ as soon as: $T$ has the laws of its core, the laws of the top structure and the short laws (all contained in
the run conditions); the core of $T$ is $J$ and is canonical (core ties); the field conditions at $\mathrm{own}(v)$ hold (one of
the core ties, for $p'_{\mathrm{de}}$ if $p_v=p_{\mathrm{de}}$ and for $p_{\mathrm{de}}$ otherwise); the error datum is
half-compatible with the core (core ties); and every auxiliary prime datum $d_\delta(j,n)$ of $T$ has level $1$ and is located at
$\mathrm{own}(v)$. The level statement is one of the laws in the run conditions, and the location statement is the own-prime pin.
\end{proof}

\subsubsection{Lattice leaves}\label{sec:pf:onto2:lat}

\begin{definition}[Records of a run and form ties]\label{pf:onto2:formtie}
The \emph{record} $X$ of a run at an admissible index $j$ (\lref{Formalism.MainNotions.WeightAlgebra.PJ}) has a coefficient
field $K_0$ (a number field; $K$ is a $K_0$-algebra), a commutative $K_0$-algebra $\mathrm{Forms}$, for each weight $w$ a
$K_0$-space $\mathrm{Cusp}(w)$ of cusp forms with an injective $K_0$-linear map $\iota_w:\mathrm{Cusp}(w)\to\mathrm{Forms}$, and
$\OO_{K_0}$-submodules $\mathrm{Mlat}(w)$ (integral forms), $\mathrm{thAll}(w)$ (theta series) and $\mathrm{shortTh}(w)$ (short theta
series) of $\mathrm{Cusp}(w)$. For $s\in\OO_{K_0}$ write $\mathrm{Thc}(s,w)=s\,(g_1g_2)^{frn(w)}\,\mathrm{shortTh}(w)$, where $n(w)$ is
the size of $w$. For a scalar datum $S=(s_W,s_F,s_0,s_K)$, weights $\kappa,w_L,w_F$ and $T\in\N$, the \emph{theta lattice} is the
$\OO_{K_0}$-submodule of $\mathrm{Forms}$
\[
\Theta=\iota(s_W\mathrm{thAll}(w_L))^{T}\iota(\mathrm{Thc}(s_F,w_F))+\iota(s_W\mathrm{thAll}(w_L))^{T-\dim Y}
\iota(\mathrm{Thc}(s_0,w_{F0}))+\iota(\mathrm{Thc}(s_K,\kappa)),
\]
with $w_{F0}=w_F+(\dim Y)w_L$, and $\Theta^{\mathrm{pl}}$ is the same expression with all scalars removed and $\mathrm{Thc}(\cdot,w)$
replaced by $\mathrm{shortTh}(w)$. For an isomorphism $z:K^n\cong
K\otimes_{K_0}\mathrm{Cusp}(\kappa)$ and a subset $Y'\subseteq\mathrm{Cusp}(\kappa)$ let $\langle Y'\rangle_z$ be the $\OO_K$-span of
$\{z^{-1}(1\otimes F):F\in Y'\}$. The \emph{form tie} at $(v,j)$ (\lref{GlobalTheory.FormTies.FormTieK}) says:
\begin{equation}\label{eq:pf:onto2:formtie}
\begin{gathered}
M_s(j)=\langle\mathrm{Mlat}(\kappa)\rangle_z,\quad \rho_j(g)z^{-1}(1\otimes F)=z^{-1}(1\otimes\mathrm{act}_gF),
\quad T_h(v,j)=\langle\{F:\iota_\kappa F\in\Theta\}\rangle_z,\\
\sigma(v,j)=\sigma_{\mathrm{rec}},\qquad T_{hK}(v,j)=\langle\mathrm{shortTh}(\kappa)\rangle_z ,
\end{gathered}
\end{equation}
for the run weights $(\kappa,w_L,w_F)=(\kappa R(j),lR(j),fR(j))$, $T=T_O$, the scalar datum of the run at $\mathrm{own}(v)$ and an
explicit non-zero scalar $\sigma_{\mathrm{rec}}\in\OO_{K_0}$ (\lref{GlobalTheory.FormTies.sigRec}). The \emph{weak form tie}
replaces the last identity by: $T_{hK}(v,j)$ is finitely generated and contains
$\langle\{F:\iota_\kappa F\in\Theta^{\mathrm{pl}}\}\rangle_z$.
\formallabel{GlobalTheory.FormTies.thetaPlain}
\formallabel{GlobalTheory.FormTies.FormTie}
\end{definition}

We use the following facts about records, proved elsewhere or among the deferred results of this section. Let $(\kappa,w_L,w_F,T)$ be
run weights at an admissible $j$; by \lref{GlobalTheory.RunDepth.run_weights_sum} they satisfy the \emph{weight relations}
(a compatibility of the three weights, $n(\kappa)=n(w_F)+Tn(w_L)$, $\dim Y\le T$ and $r\le n(w_F)$), and $n(\kappa)=n_V$
(\lref{Formalism.MainNotions.DIV.Cons.kaWt_n}).
\begin{enumerate}[label=(F\arabic*)]
\item\label{pf:onto2:F1} \emph{Integrality.} Under the integrality hypothesis of the record, $\Theta\subseteq
  \iota_\kappa(\mathrm{Mlat}(\kappa))$ (\lref{Formalism.MainNotions.Gaps.T0.theta_le_msup}).
\item\label{pf:onto2:F2} \emph{Scaling.} $\sigma_{\mathrm{rec}}\ne0$ and every element of $\Theta$ is $\sigma_{\mathrm{rec}}$ times an
  element of $\Theta^{\mathrm{pl}}$ (\lref{GlobalTheory.T.theta_scaled_span_le_sigma_smul_unscaled_span}).
\item\label{pf:onto2:F3} \emph{Products of short forms.} If the top structure has its laws and short laws, the one-cover condition
  holds at $j$ and $d_\delta(j,0)$ has level $1$, then $\Theta^{\mathrm{pl}}\subseteq\iota_\kappa(\mathrm{shortTh}(\kappa))$; this is
  \cref{pf:onto2:prodshort} below.
\item\label{pf:onto2:F4} \emph{Level action.} The level action $\mathrm{act}$ attached to the level maps of the run preserves
  $\mathrm{Mlat}(\kappa)$ (\lref{GlobalTheory.FormTies.levAct_mem_Mlat}) and, if the run has all laws and its twist character
  is trivial on determinants of the rational subgroup, it preserves $\Theta$ (\lref{GlobalTheory.FormTies.levAct_theta}).
\item\label{pf:onto2:F5} \emph{Finiteness and spanning.} If $\mathrm{Mlat}(\kappa)$ is a lattice and the integrality hypothesis
  holds, $\mathrm{shortTh}(\kappa)$ is finitely generated (\lref{GlobalTheory.FormTies.shortTh_fg}). The index-bound
  hypotheses of a run imply that $\mathrm{Mlat}(\kappa)$ is a lattice, that $\mathrm{shortTh}(\kappa)$ spans $\mathrm{Cusp}(\kappa)$
  over $K_0$, and the integrality hypothesis (\lref{GlobalTheory.BigOn.lat_span_of_hypC}, from
  \lref{GlobalTheory.latAt_record} and \lref{GlobalTheory.RunLat.spanAt_of_shortG}).
\end{enumerate}

\begin{lemma}[Lattice algebra under a form tie]\label{pf:onto2:latalg}
Assume the weak form tie at $(v,j)$ for run weights as above and a prime $p>0$ (in the applications $p=\mathrm{own}(v)$), that
$\mathrm{Mlat}(\kappa)$ is a lattice preserved by $\mathrm{act}$, that $\mathrm{shortTh}(\kappa)$ spans $\mathrm{Cusp}(\kappa)$, that
the integrality hypothesis holds and that $\mathrm{act}$ preserves $\Theta$. Then: $M_s(j)$ is a lattice in $K^{n_Z(j)}$ stable
under $\rho_j(P)$; $T_h(v,j)$ and $T_{hK}(v,j)$ are lattices in $K^{n_Z(j)}$; $T_h(v,j)\subseteq M_s(j)$, $T_h(v,j)$ is
$\rho_j(P)$-stable, $\sigma(v,j)\ne0$ and $T_h(v,j)\subseteq\sigma(v,j)T_{hK}(v,j)$.
\formallabel{GlobalTheory.FormTies.supLattice_isLatIn_stable}
\formallabel{GlobalTheory.FormTies.theta_spans_isLatIn}
\formallabel{GlobalTheory.scaledSpan_clauses}
\end{lemma}

\begin{proof}
For a finitely generated $\OO_{K_0}$-module $Y'\subseteq\mathrm{Cusp}(\kappa)$ spanning it over $K_0$, $\langle Y'\rangle_z$ is a
lattice in $K^{n}$: it is generated by the images of finitely many generators, and its $K$-span contains
$z^{-1}(1\otimes\mathrm{Cusp}(\kappa))$, which spans $K^n$ since $z$ is a $K$-linear isomorphism. Moreover $\langle\cdot\rangle_z$ is
monotone and $\langle aY'\rangle_z=a\langle Y'\rangle_z$ for $a\in K_0$.

\emph{$M_s(j)$.} By \eqref{eq:pf:onto2:formtie} $M_s(j)=\langle\mathrm{Mlat}(\kappa)\rangle_z$, a lattice. For $g\in P$ and
$F\in\mathrm{Mlat}(\kappa)$, $\rho_j(g)z^{-1}(1\otimes F)=z^{-1}(1\otimes\mathrm{act}_gF)\in M_s(j)$ since $\mathrm{act}_g$ preserves
$\mathrm{Mlat}(\kappa)$; as $\rho_j(g)$ is $\OO_K$-linear, $\rho_j(g)M_s(j)\subseteq M_s(j)$.

\emph{$T_h\subseteq M_s$ and stability.} If $\iota_\kappa F\in\Theta$ then by \ref{pf:onto2:F1} $\iota_\kappa F=\iota_\kappa G$ with
$G\in\mathrm{Mlat}(\kappa)$, so $F=G$ by injectivity of $\iota_\kappa$. Hence $\{F:\iota_\kappa F\in\Theta\}\subseteq\mathrm{Mlat}(\kappa)$
and $T_h(v,j)\subseteq M_s(j)$; in particular $T_h(v,j)$ is finitely generated ($\OO_K$ is noetherian). The same computation as
for $M_s$, using that $\mathrm{act}$ preserves $\Theta$, shows that $T_h(v,j)$ is $\rho_j(P)$-stable.

\emph{$T_h$ spans.} Let $c=s_K(g_1g_2)^{frn(\kappa)}\in\OO_{K_0}$; $c\ne0$ because $s_K\ne0$ for $p>0$
(\lref{GlobalTheory.FormTies.scalOfP_sK_ne_zero}) and $g_1,g_2$ are primes. For $F\in\mathrm{shortTh}(\kappa)$ we have
$cF\in\mathrm{Thc}(s_K,\kappa)$, so $\iota_\kappa(cF)$ lies in the third summand of $\Theta$ and $z^{-1}(1\otimes cF)\in T_h(v,j)$;
hence $z^{-1}(1\otimes F)=c^{-1}z^{-1}(1\otimes cF)$ lies in the $K$-span of $T_h(v,j)$. Since $\mathrm{shortTh}(\kappa)$ spans
$\mathrm{Cusp}(\kappa)$, $T_h(v,j)$ spans $K^{n_Z(j)}$.

\emph{$T_{hK}$.} It is finitely generated by the weak form tie and contains
$\langle\{F:\iota_\kappa F\in\Theta^{\mathrm{pl}}\}\rangle_z\supseteq\langle\mathrm{shortTh}(\kappa)\rangle_z$ (the third summand of
$\Theta^{\mathrm{pl}}$ is $\iota(\mathrm{shortTh}(\kappa))$); hence $T_{hK}(v,j)$ spans.

\emph{Scaling.} By \ref{pf:onto2:F2}, $\sigma(v,j)=\sigma_{\mathrm{rec}}\ne0$. Let $F$ with $\iota_\kappa F\in\Theta$ and write
$\iota_\kappa F=\sigma_{\mathrm{rec}}G$ with $G\in\Theta^{\mathrm{pl}}$. Then $\iota_\kappa(\sigma_{\mathrm{rec}}^{-1}F)=G\in\Theta^{\mathrm{pl}}$,
so $z^{-1}(1\otimes\sigma_{\mathrm{rec}}^{-1}F)\in T_{hK}(v,j)$ and $z^{-1}(1\otimes F)=\sigma(v,j)\,z^{-1}(1\otimes\sigma_{\mathrm{rec}}^{-1}F)\in
\sigma(v,j)T_{hK}(v,j)$. Since the generators of $T_h(v,j)$ lie in the $\OO_K$-module $\sigma(v,j)T_{hK}(v,j)$, so does $T_h(v,j)$.
\end{proof}

\begin{lemma}[Averaging over $U$]\label{pf:onto2:avg}
Let $\rho$ be a representation of a group $P$ on a $K$-vector space $Z$, $U\le P$ a finite subgroup and $\psi:P\to\OO_K$ with
$\psi(1)=1$ and $\psi(ab)=\psi(a)\psi(b)$ for $a,b\in U$. Put $\mathrm{tr}_{U,\psi}=\sum_{u\in U}\psi(u)\rho(u)^{-1}$. Then
$\mathrm{tr}_{U,\psi}(Z)\subseteq\ontoIIeig(\rho,U,\psi)$, and $\mathrm{tr}_{U,\psi}(y)=|U|\,y$ for $y\in\ontoIIeig(\rho,U,\psi)$.
\formallabel{GlobalTheory.TraceLattice.trace_mem_eig}
\formallabel{GlobalTheory.TraceLattice.trace_of_mem_eig}
\end{lemma}

\begin{proof}
For $a\in U$, $\rho(a)\mathrm{tr}_{U,\psi}(y)=\sum_u\psi(u)\rho(au^{-1})y$; substituting $u=u'a$ gives
$\sum_{u'}\psi(u'a)\rho(u'^{-1})y=\psi(a)\mathrm{tr}_{U,\psi}(y)$. If $y\in\ontoIIeig(\rho,U,\psi)$ then $\rho(u)^{-1}y=\psi(u^{-1})y$
and $\psi(u)\psi(u^{-1})=\psi(1)=1$, so every term equals $y$.
\end{proof}

\paragraph{Products of short theta series.} For a top structure $T$ and an admissible index $j$, the \emph{theta series of
weight $w_L$} are the sections of the moduli space $X_j$ attached to column vectors $x$ of rank $m=n(w_L)$ and harmonic
polynomials $P$ of weight $w_L$ for the standard lattice of rank $m$ (\lref{Formalism.L2.Pj.thGen}: the section whose values
at the points of $X_j$ are given by $P$ evaluated on the theta coordinates $x$), and the \emph{short theta series of weight $w$}
are the sections attached to a unimodular lattice $\Lambda$ of rank $n=n(w)$ which is trivial at the bad primes, a level $N$
compatible with the auxiliary datum $d_\delta(j,n)$, an admissible column vector $x_0$ and a harmonic polynomial $P$ of weight $w$
for $\Lambda$, whose values on the cover of $d_\delta(j,n)$ are the \emph{short values} of $P$ normalised by the constant
$\beta(\Lambda)$ (\lref{Formalism.L2.Pj.shortGen}, \lref{Formalism.L2.Pj.betaH}: $\beta(\Lambda)$ is a non-zero element of
$K_0$ generating the quotient of the determinant ideals of the two lattices of $\Lambda$, of absolute value fixed by their index).
The sets $\mathrm{thAll}(w)$ and $\mathrm{shortTh}(w)$ of the record are the $\OO_{K_0}$-modules generated by these sections. The
\emph{chart-independence statement} says that the value of a harmonic polynomial on the germs of sections at a frame point does
not depend on the chart; it is proved in the development (\lref{GlobalTheory.chartIndependence_unconditional}).

\begin{proposition}[Short theta series times theta series]\label{pf:onto2:shortmul}
Assume the chart-independence statement. Let $T$ be a top structure with its laws and its short laws, and $j$ an admissible index
at which the one-cover condition holds and $d_\delta(j,0)$ has level $1$. Let $w,w_L$ be weights at $j$ with $r\le n(w)$, and such
that $n(w_L)=0$ if $n(w)=n_V$. Then for every short theta series $g$ of weight $w$ and every theta series $f$ of weight $w_L$ the
product $gf$ lies in the $\OO_{K_0}$-span of the short theta series of weight $w+w_L$.
\formallabel{GlobalTheory.ShortTimesTheta.short_mul_theta_mem_span}
\end{proposition}

\begin{proof}
Theta series and short theta series are sections of the relevant line bundles (\lref{GlobalTheory.thetaRule_mem_sec}), and
two sections which agree at all points of $X_j$ over algebraically closed fields are equal
(\lref{GlobalTheory.ClassSumUnique.sec_eq_of_geometric}). If $X_j$ has no such point, then $gf=0$ and there is nothing to prove.
Assume from now on that there is one.

\emph{Coefficients of the auxiliary data.} The pins of the laws and of the short laws give the same star structure on $\OO_E$, and
the signature $(p,q)$ at $j$ satisfies $p+q=r\ge1$. Put $n=n(w)$, $m=n(w_L)$. The one-cover condition gives a coefficient function
$\varphi'$ on $(\mathrm{La}/M\mathrm{La})^{n+m}$ such that $d_\delta(j,n+m)$ is $d_\delta(j,n)$ with coefficients $\varphi'$. The short
laws provide, at the sizes $n$ and $n+m$, cover characters $\psi,\psi'$ with their defining identities; since a cover character is
determined by these identities at a geometric point (\lref{GlobalTheory.CoverCharacter.psi_unique}), $\psi=\psi'$. Comparing the
two coefficient identities (\lref{GlobalTheory.MoreColumns.coeff_more_columns}) gives a root of unity $\zeta\in K_0$ with
\begin{equation}\label{eq:pf:onto2:phiprime}
\varphi'(q)=\zeta\,\varphi_\delta\bigl(q|_{[1,n]}\bigr)\cdot\bigl[q_b\in I\cdot\mathrm{La}\ \text{for}\ b>n\bigr]\cdot|\mathrm{La}/I\,\mathrm{La}|^{m},
\end{equation}
where $\varphi_\delta$ are the coefficients of $d_\delta(j,n)$ and $I$ its ideal.

\emph{The product formula} (\lref{GlobalTheory.ShortTimesTheta.mul_mem_span}). Let $g$ be attached to $(\Lambda,N,x_0,P')$ and $f$
to $(x,P)$. The block lattice $\Lambda\oplus\mathrm{std}_m$ is again unimodular, trivial at the bad primes and of compatible level,
and $P'P$ is harmonic of weight $w+w_L$ for it; moreover the short value of $P'P$ at a point of the cover is the product of the
short value of $P'$ and the value of $P$ (\lref{GlobalTheory.shortSeries_blockProduct}). The constants satisfy
$\beta(\Lambda\oplus\mathrm{std}_m)=\zeta'\beta(\Lambda)\beta(\mathrm{std}_m)$ with $\beta(\mathrm{std}_m)=1$ and a root of unity
$\zeta'$ (\lref{GlobalTheory.BetaOfSum.blockBeta}). For every vector $x'$ of $m$ further columns, $x_0$ extended by $x'$ is
admissible (\lref{GlobalTheory.ShortTimesTheta.admissibleColumns_append}), and the total short-value condition of the short laws
provides a short theta series $F(x')$ of weight $w+w_L$ attached to $(\Lambda\oplus\mathrm{std}_m,N,(x_0,x'),P'P)$. We claim
\begin{equation}\label{eq:pf:onto2:gfsum}
gf=\sum_{x'}\zeta'\zeta^{-1}c(x')\,F(x'),\qquad c(x')=\zeta_M^{-\langle x',\,p_\delta x\rangle},
\end{equation}
where $\zeta_M$ is the primitive $M$-th root of unity of the coefficient field, $\langle\cdot,\cdot\rangle$ the level pairing and
$p_\delta$ the integer of $d_\delta$. Both sides are sections, so it suffices to compare their values at a point $y$ of $X_j$ over an
algebraically closed field. Lift $y$ to a point $y_2$ of the cover of $d_\delta(j,n)$. The value of $g$ at $y$ is a value of the
short series of $P'$ at $y_2$, the value of $f$ at $y$ equals the short value at $y_2$ of the standard theta polynomial of $(x,P)$
(\lref{GlobalTheory.T.standardCover_shortValue_eq_thetaValue}), and by \lref{GlobalTheory.shortSeries_blockProduct} their
product is the short value of $P'P$ for the block lattice with the column-coefficient function of \eqref{eq:pf:onto2:phiprime}.
Expanding the indicator in \eqref{eq:pf:onto2:phiprime} as a character sum over $x'$
(\lref{GlobalTheory.ShortTimesTheta.sum_data}) and using the linearity of short values in the coefficient function
(\lref{GlobalTheory.ShortTimesTheta.shortVal_sum}), this value equals the value of the right-hand side of
\eqref{eq:pf:onto2:gfsum} multiplied by $\zeta\beta(\Lambda)$; cancelling the non-zero factor $\zeta\beta(\Lambda)$ proves the claim.
Finally $\zeta$, $\zeta'$ and $c(x')$ are roots of unity, so the coefficients $\zeta'\zeta^{-1}c(x')$ are integral over $\Z$ and lie in
$\OO_{K_0}$; hence \eqref{eq:pf:onto2:gfsum} exhibits $gf$ as an element of the $\OO_{K_0}$-span of short theta series of weight
$w+w_L$.
\end{proof}

\begin{proposition}[Products of short forms]\label{pf:onto2:prodshort}
Assume the chart-independence statement. Let $T$ be a top structure with its laws and short laws and $j$ an admissible index at
which the one-cover condition holds and $d_\delta(j,0)$ has level $1$. Let $\kappa,w_L,w_F$ be weights and $T'\in\N$ satisfying the weight relations (in particular
$\kappa=w_F+T'w_L$ entrywise), with $n(\kappa)=n(w_F)+T'n(w_L)=n_V$, $\dim Y\le T'$ and $r\le n(w_F)$. Then the plain theta lattice
$\Theta^{\mathrm{pl}}$ of the record of $T$ at $j$ (formed with $T'$) satisfies $\Theta^{\mathrm{pl}}\subseteq\iota_\kappa(\mathrm{shortTh}(\kappa))$.
\formallabel{GlobalTheory.FormTies.productsShort_rec}
\end{proposition}

\begin{proof}
\emph{One factor.} Let $w$ be a weight with $r\le n(w)$ and such that $n(w_L)=0$ if $n(w)=n_V$. Then
\begin{equation}\label{eq:pf:onto2:onefactor}
\iota(\mathrm{thAll}(w_L))\cdot\iota(\mathrm{shortTh}(w))\subseteq\iota(\mathrm{shortTh}(w+w_L)).
\end{equation}
Indeed, the left-hand side is generated by the products $\iota(F)\iota(G)$ of generators $F$ of $\mathrm{thAll}(w_L)$ and $G$ of
$\mathrm{shortTh}(w)$; the product $GF$ lies in the forms of weight $w+w_L$ by the multiplicativity law of the core, in the span of
the short theta series of weight $w+w_L$ by \cref{pf:onto2:shortmul}, and $\iota(G)\iota(F)=\iota(GF)$ because the embeddings are
compatible with products.

\emph{Iteration.} Let $s\in\N$ and $w$ a weight with $r\le n(w)$ and $w+sw_L=\kappa$ (in degree). Then
$\iota(\mathrm{thAll}(w_L))^s\iota(\mathrm{shortTh}(w))\subseteq\iota(\mathrm{shortTh}(\kappa))$. This follows by induction on $s$
from \eqref{eq:pf:onto2:onefactor} applied to $w,w+w_L,\dots,w+(s-1)w_L$: these weights have size at least $n(w)\ge r$, and if
$n(w+iw_L)=n_V$ for some $i<s$, then $n_V=n(\kappa)=n(w+iw_L)+(s-i)n(w_L)$ forces $n(w_L)=0$, so the side condition of
\eqref{eq:pf:onto2:onefactor} holds at every step.

\emph{Conclusion.} By definition $\Theta^{\mathrm{pl}}$ is the sum of
$\iota(\mathrm{thAll}(w_L))^{T'}\iota(\mathrm{shortTh}(w_F))$, of $\iota(\mathrm{thAll}(w_L))^{T'-\dim Y}\iota(\mathrm{shortTh}(w_{F0}))$ with
$w_{F0}=w_F+(\dim Y)w_L$, and of $\iota(\mathrm{shortTh}(\kappa))$. The iteration applies to the first with $(s,w)=(T',w_F)$, since
$w_F+T'w_L=\kappa$, and to the second with $(s,w)=(T'-\dim Y,w_{F0})$, since $\dim Y\le T'$ gives $w_{F0}+(T'-\dim Y)w_L=\kappa$ and
$n(w_{F0})\ge n(w_F)\ge r$. The third summand is trivially contained in $\iota(\mathrm{shortTh}(\kappa))$.
\end{proof}

\begin{lemma}[Lattice leaves from form ties]\label{pf:onto2:lemmaS}
Assume the chart-independence statement. Suppose that for every place $v$ and every index $i$ there are a full top structure
$T$ with its laws, an admissible index $j$, a prime $p>0$, identifications $z$, $\mathrm{act}$ and level maps such that: the form tie
holds at $(v,i)$ for the record of $T$ at $j$, the run weights at $j$ and $p$; $\mathrm{act}$ is the level action of the level maps;
$T$ has all its laws; the twist character is trivial on determinants of the rational subgroup; the one-cover condition holds at
$j$; $\mathrm{Mlat}(\kappa)$ is a lattice; $\mathrm{shortTh}(\kappa)$ spans $\mathrm{Cusp}(\kappa)$; and the integrality hypothesis
holds. Then $O$ satisfies \ref{pf:onto2:L1}, \ref{pf:onto2:L2} and \ref{pf:onto2:L3}.
\formallabel{GlobalTheory.FormTies.leaves_sup_span_lemmaS_of_formTieK}
\end{lemma}

\begin{proof}
Fix $(v,i)$ and the data of the hypothesis. The run weights at $j$ satisfy the weight relations
(\lref{GlobalTheory.RunDepth.run_weights_sum}) and $n(\kappa)=n_V$ (\lref{Formalism.MainNotions.DIV.Cons.kaWt_n}). The form
tie implies the weak form tie: $\mathrm{shortTh}(\kappa)$ is finitely generated by \ref{pf:onto2:F5}, and $T_{hK}(v,i)=\langle
\mathrm{shortTh}(\kappa)\rangle_z$ contains $\langle\{F:\iota_\kappa F\in\Theta^{\mathrm{pl}}\}\rangle_z$ by \cref{pf:onto2:prodshort}
(its hypotheses are the laws, the one-cover condition, the level law of all laws, and the weight relations). By \ref{pf:onto2:F4}
the level action preserves $\mathrm{Mlat}(\kappa)$ and $\Theta$. Hence \cref{pf:onto2:latalg} applies at $(v,i)$: $M_s(i)$ is a
$\rho_i(P)$-stable lattice, $T_h(v,i)$ and $T_{hK}(v,i)$ are lattices, and $T_h(v,i)\subseteq M_s(i)$ is $\rho_i(P)$-stable with
$T_h(v,i)\subseteq\sigma(v,i)T_{hK}(v,i)$, $\sigma(v,i)\ne0$. Taking for $v$ any place (one exists,
\lref{Formalism.L2.Indef.Tree.exists_place3}) gives \ref{pf:onto2:L1}; taking $v\in\{w,\bar w\}$ for $w\in\ontoIITw$ gives
\ref{pf:onto2:L2} and \ref{pf:onto2:L3}.
\end{proof}

\begin{proposition}[The lattice leaves]\label{pf:onto2:latleaves}
For every $C_K$ there is $x_1$ such that for all data as above: if the record relation, the core ties, the own-prime pin, the
page-run conditions $\ontoIIPN(x_1)$, the run bounds $\ontoIIRB(C_K,r)$, the group condition of $(R,F)$ and the three tuple
conditions of $R$ hold, then $O$ satisfies \ref{pf:onto2:L1}--\ref{pf:onto2:L4}.
\formallabel{GlobalTheory.BigOn.leaves_sup_span_data}
\end{proposition}

\begin{proof}
Take $x_1=0$. By component (ii) of the record relation the weights of the face are the run weights. Let $v$ be a place and $j$
an index, let $T=r_v$ (a full top structure with its laws, by the run conditions) and let $z_j$, $\mathrm{act}_j$
be the identifications of component (iii); the trace tie at $(v,j)$ contains the form tie and the statement that
$\mathrm{act}_j$ is the level action of the level maps of $r$ composed with the inclusion $\mathrm{Mir}(R)\to G$. By the run bounds
and \ref{pf:onto2:F5}, $\mathrm{Mlat}(\kappa)$ is a lattice, $\mathrm{shortTh}(\kappa)$ spans, and the integrality hypothesis holds;
the rationality of the twist character and the one-cover condition are run conditions. Hence the hypotheses of
\cref{pf:onto2:lemmaS} hold at every $(v,j)$, and the chart-independence statement holds
(\lref{GlobalTheory.chartIndependence_unconditional}); \cref{pf:onto2:lemmaS} gives \ref{pf:onto2:L1}--\ref{pf:onto2:L3}.

For \ref{pf:onto2:L4} let $w\in\ontoIITw$, $v\in\{w,\bar w\}$ and $j$. The trace tie identifies $U$ with the preimage of a subgroup
$U'\le G$, $\psi$ with a character of $U'$ (so $\psi(1)=1$ and $\psi$ is multiplicative on $U$), and $T_{h0}(v,j)$ with a lattice
$\Theta^{\mathrm{eig}}$ spanned by eigen-combinations of theta series (\lref{GlobalTheory.OnTies.thetaEig}). The group
condition identifies the subgroup $F.P_0$ with the elements of $G$ whose images at all places of $R.P_g$ are unipotent, and with the
level pins of the run conditions this gives that the level maps of $U'$ have determinant one
(\lref{GlobalTheory.TraceLattice.detOneOn_of_pins}); moreover $\kappa$ has size $n_V$ and positive last entry
(\lref{GlobalTheory.TraceLattice.kaR_lastPos}). Under these conditions
(\lref{GlobalTheory.TraceLattice.clauses_of_onTie}): (a) $T_{h0}(v,j)\subseteq\ontoIIeig(\rho_j,U,\psi)$, since each generator of
$\Theta^{\mathrm{eig}}$ is the image of a vector under $\mathrm{tr}_{U,\psi}$ and \cref{pf:onto2:avg} applies
(\lref{GlobalTheory.TraceLattice.thetaEig_le_eig}); (b) $\mathrm{tr}_{U,\psi}$ maps $\langle\mathrm{shortTh}(\kappa)\rangle_z=T_{hK}(v,j)$
into $T_{h0}(v,j)$ (\lref{GlobalTheory.TraceLattice.trace_mem_thetaEig}); (c) $T_{h0}(v,j)\subseteq T_{hK}(v,j)$, so $T_{h0}(v,j)$
is finitely generated since $T_{hK}(v,j)$ is. Finally, for $y\in T_{hK}(v,j)\cap\ontoIIeig(\rho_j,U,\psi)$, \cref{pf:onto2:avg} and (b)
give $|U|\,y=\mathrm{tr}_{U,\psi}(y)\in T_{h0}(v,j)$. This is \ref{pf:onto2:L4}
(\lref{GlobalTheory.TraceLattice.leafData_of_onTie_of_leafSpan}).
\end{proof}

\begin{remark}\label{pf:onto2:rem:latunused}
In \cref{pf:onto2:latleaves} the core ties, the own-prime pin, the page-run conditions and the three tuple conditions are not used,
and $x_1=0$. The constant $C_K$ enters only through the run bounds $\ontoIIRB(C_K,r)$, of which only the index-bound hypotheses are
used.
\end{remark}

\subsubsection{The local index bound}\label{sec:pf:onto2:index}

The index bound \ref{pf:onto2:L5} is obtained from the run conclusion, which is a statement over the coefficient field $K_0$ of
the record, by a local--global passage to the places of $K$.

\begin{definition}[The run conclusion]\label{pf:onto2:concl}
Let $X$ be a record with coefficient field $K_0$, $(S,\kappa,w_L,w_F,T)$ as in \cref{pf:onto2:formtie}, $\nu_s$ a real function on
the finite places of $K_0$ and $\beta\in\R$. For a place $\lambda$ of $K_0$ let $K_{0,\lambda}$, $\OO_\lambda$ be the completions,
$V_\lambda=K_{0,\lambda}\otimes_{K_0}\mathrm{Forms}$, and for an $\OO_{K_0}$-submodule $A\subseteq\mathrm{Forms}$ let $\widehat A_\lambda$ be
the $\OO_\lambda$-span of $1\otimes A$. Then $\ontoIICl(X,S,\kappa,w_L,w_F,T,\nu_s,\beta)$
means: there is $r_0:\{\lambda\}\to\R_{\ge0}$ such that
$\sum_\lambda r_0(\lambda)\ell_\lambda$ converges with sum $\le\beta$ ($\ell_\lambda=\log\Nm\lambda/[K_0:\Q]$), and for every $\lambda$,
every integer $e\ge T\nu_s(\lambda)$ and every $K_{0,\lambda}$-subspace $V'$ of the $K_{0,\lambda}$-span of $\widehat{M}_\lambda$, where
$M=\iota_\kappa(\mathrm{Mlat}(\kappa))$, the $\OO_\lambda$-module $(\widehat M_\lambda\cap V')/(\widehat\Theta_\lambda\cap V')$ has finite
length, at most $e\dim V'+r_0(\lambda)$.
\formallabel{Formalism.MainNotions.WeightAlgebra.Concl}
\end{definition}

We use three standard facts about completion, for a Dedekind domain $R$ with fraction field $K_0$, a maximal ideal $\lambda$ and a
$K_0$-vector space $F$: (i) for $R$-submodules $A,T\subseteq F$ and a subspace $W$, if $(\widehat A_\lambda\cap\widehat
W_\lambda)/(\widehat T_\lambda\cap\widehat W_\lambda)$ has finite length, then $\len\bigl(\ontoIIloc{\lambda}{(A\cap
W)}/\ontoIIloc{\lambda}{(A\cap W\cap T)}\bigr)$ is at most this length (\lref{Formalism.MainNotions.Gaps.LocalGlobalPassage.l1});
(ii) $\dim_{K_{0,\lambda}}\widehat W_\lambda=\dim W$ for finite-dimensional $W$
(\lref{Formalism.MainNotions.Gaps.LocalGlobalPassage.l1b}); (iii) if $W$ lies in the $K_0$-span of $A$ then $\widehat W_\lambda$ lies
in the $K_{0,\lambda}$-span of $\widehat A_\lambda$ (\lref{Formalism.MainNotions.Gaps.LocalGlobalPassage.l1c}). These follow from the
flatness of $\OO_\lambda$ over $R_{(\lambda)}$ and $\OO_\lambda/\lambda^n\OO_\lambda=R_{(\lambda)}/\lambda^nR_{(\lambda)}$
\cite[Ch.~II]{Serre1979}, \cite[Thms.~7.2, 8.8]{Matsumura1989}. We also use the following facts on local indices of lattices in a
finite-dimensional $K$-vector space $Z$: for lattices $A,B\subseteq Z$ with $B\subseteq A$, $a\in K^\times$ with $v_w(a)=E\in\N$ and
$\rho\in\R$,
\begin{equation}\label{eq:pf:onto2:threeforms}
\Bigl(\forall W\subseteq Z:\ \ontoIIix_w(\ontoIIcut_WA,\ontoIIcut_WB)\le E\dim W+\rho\Bigr)\iff
\ontoIIix_w(A,A\cap a^{-1}B)\le\rho
\end{equation}
(\lref{Formalism.EigenvalueLevelCeiling.three_forms}); $\ontoIIix_w(A,B')\le\ontoIIix_w(A,B)$ for lattices $B\subseteq B'$ in a
common subspace with $A$ (\lref{Formalism.EigenvalueLevelCeiling.index_monotone}); and $A\cap aB$ is a lattice for lattices $A,B$
and $a\ne0$ (\lref{Formalism.EigenvalueLevelCeiling.lattice_intersection}).

\begin{lemma}[Base change of indices]\label{pf:onto2:l3}
Let $K/K_0$ be an extension of number fields, $Z_0$ a finite-dimensional $K_0$-vector space and $\iota:Z_0\to K^n$ a $K_0$-linear
map which sends $K_0$-linearly independent families to $K$-linearly independent families and whose image spans $K^n$ over $K$. For
an $\OO_{K_0}$-submodule $A_0\subseteq Z_0$ let $\mathrm{up}(A_0)$ be the $\OO_K$-span of $\iota(A_0)$. Then for full lattices
$A_0,C_0$ in $Z_0$ and every place $v$ of $K$ over the place $\lambda$ of $K_0$,
\[
\ontoIIix_v\bigl(\mathrm{up}(A_0),\mathrm{up}(C_0)\bigr)=e(v|\lambda)\,\ontoIIix_\lambda(A_0,C_0).
\]
Moreover $v(a)=e(v|\lambda)\lambda(a)$ for $a\in K_0^\times$, and $e(v|\lambda)\ne0$.
\formallabel{Formalism.MainNotions.Gaps.LocalGlobalPassage.l3}
\formallabel{Formalism.MainNotions.Gaps.LocalGlobalPassage.l3a}
\formallabel{Formalism.MainNotions.Gaps.LocalGlobalPassage.l3e}
\end{lemma}

\begin{proof}
The valuation statement is the definition of the ramification index: $\lambda\OO_{K,(v)}=v^{e(v|\lambda)}\OO_{K,(v)}$, and
$e(v|\lambda)\ge1$. By hypothesis $\iota$ extends to a $K$-linear isomorphism $K\otimes_{K_0}Z_0\cong K^n$. The local ring
$R=\OO_{K_0,(\lambda)}$ is a discrete valuation ring with uniformiser $\pi$; the localisations $\ontoIIloc{\lambda}{A_0}$,
$\ontoIIloc{\lambda}{C_0}$ are free $R$-modules of rank $d=\dim Z_0$, and by the elementary divisor theorem there are a basis
$b_1,\dots,b_d$ of $\ontoIIloc\lambda{A_0}$ and integers $a_1,\dots,a_d$ such that $\pi^{a_1}b_1,\dots,\pi^{a_d}b_d$ is a basis of
$\ontoIIloc\lambda{C_0}$; then $\ontoIIix_\lambda(A_0,C_0)=\sum_ia_i$ (both sides of \eqref{eq:pf:onto2:ix} are additive along the
decomposition, and $\ontoIIix(R,\pi^aR)=a$). Localising $\mathrm{up}(A_0)$ at $v$ gives the $\OO_{K,(v)}$-span of
$\iota(\ontoIIloc\lambda{A_0})$, which is free with basis $\iota(b_i)$ (these are $K$-linearly independent), and similarly
$\ontoIIloc{v}{\mathrm{up}(C_0)}$ has basis $\pi^{a_i}\iota(b_i)$. Since $v(\pi)=e(v|\lambda)$, the same computation gives
$\ontoIIix_v(\mathrm{up}A_0,\mathrm{up}C_0)=\sum_ie(v|\lambda)a_i$.
\end{proof}

\begin{lemma}[Sums over places]\label{pf:onto2:l4}
Let $K/K_0$ be number fields and $r_0\ge0$ a function on the places of $K_0$ with $\sum_\lambda r_0(\lambda)\ell_\lambda<\infty$. Then
for every finite set $S'$ of places of $K$,
\[
\sum_{v\in S'}\ell_v\,e(v|\lambda_v)\,r_0(\lambda_v)\le\sum_\lambda r_0(\lambda)\ell_\lambda ,
\]
where $\lambda_v$ is the place of $K_0$ below $v$.
\formallabel{Formalism.MainNotions.Gaps.LocalGlobalPassage.l4}
\end{lemma}

\begin{proof}
For $v|\lambda$ with residue degree $f(v|\lambda)$ we have $\log\Nm v=f(v|\lambda)\log\Nm\lambda$, hence
$\ell_v=f(v|\lambda)\,\ell_\lambda\,[K_0:\Q]/[K:\Q]$. Group the sum over $S'$ according to $\lambda_v$. For fixed $\lambda$,
$\sum_{v\in S',v|\lambda}e(v|\lambda)f(v|\lambda)\le[K:K_0]$ \cite[Prop.~I.8.2]{Neukirch1999}, so the contribution of $\lambda$ is at
most $r_0(\lambda)\ell_\lambda$. Summing over the finitely many $\lambda$ below $S'$ and using $r_0\ge0$ gives the claim.
\end{proof}

\begin{lemma}[Local--global passage]\label{pf:onto2:l5}
Let $\iota:Z_0\to K^n$ be as in \cref{pf:onto2:l3}, let $M_0,N_0$ be $\OO_{K_0}$-submodules of $Z_0$ spanning $Z_0$ with $M_0$
finitely generated and $N_0\subseteq M_0$, let $v$ be a place of $K$ above $\lambda$, $e_0\in\Z$, $\rho\in\R$ and $E\in\N$ with
$E=e(v|\lambda)e_0$. Assume that for every $K_0$-subspace $W_0\subseteq Z_0$
\[
\ontoIIix_\lambda(M_0\cap W_0,N_0\cap W_0)\le e_0\dim W_0+\rho .
\]
Then for every $\OO_K$-submodule $T_h\supseteq\mathrm{up}(N_0)$ and every $K$-subspace $W\subseteq K^n$,
\[
\ontoIIix_v\bigl(\ontoIIcut_W\mathrm{up}(M_0),\ontoIIcut_WT_h\bigr)\le E\dim W+e(v|\lambda)\rho .
\]
\formallabel{Formalism.MainNotions.Gaps.LocalGlobalPassage.l5}
\end{lemma}

\begin{proof}
Write $e=e(v|\lambda)\ge1$. Since $E\ge0$ and $e>0$, $e_0\ge0$. Choose $a_0\in K_0^\times$ with $\lambda(a_0)=e_0$
(a power of an element of $\lambda\setminus\lambda^2$); $N_0$ is finitely generated as a submodule of the
noetherian module $M_0$, so $M_0$ and $N_0$ are full lattices. By \eqref{eq:pf:onto2:threeforms} over $K_0$ (with $a=a_0$, $E=e_0$),
the hypothesis gives $\ontoIIix_\lambda(M_0,C_0)\le\rho$ for the lattice $C_0=M_0\cap a_0^{-1}N_0$. By \cref{pf:onto2:l3},
$\ontoIIix_v(\mathrm{up}M_0,\mathrm{up}C_0)=e\,\ontoIIix_\lambda(M_0,C_0)\le e\rho$. Let $a$ be the image of $a_0$ in $K$; then
$v(a)=ee_0=E$. If $x\in C_0$ then $a_0x\in N_0$, so $\iota(x)\in a^{-1}\mathrm{up}(N_0)$; hence $\mathrm{up}(C_0)\subseteq\mathrm{up}(M_0)\cap
a^{-1}\mathrm{up}(N_0)$, a lattice, and by monotonicity
\[
\ontoIIix_v\bigl(\mathrm{up}M_0,\mathrm{up}M_0\cap a^{-1}\mathrm{up}N_0\bigr)\le\ontoIIix_v(\mathrm{up}M_0,\mathrm{up}C_0)\le e\rho .
\]
By \eqref{eq:pf:onto2:threeforms} over $K$ (with $a$, $E$) this gives, for every $W$,
$\ontoIIix_v(\ontoIIcut_W\mathrm{up}M_0,\ontoIIcut_W\mathrm{up}N_0)\le E\dim W+e\rho$. Finally put $A=\ontoIIcut_W\mathrm{up}M_0$,
$B'=\ontoIIcut_W\mathrm{up}N_0$ and $B=\ontoIIcut_WT_h$; $A$ and $B'$ are lattices in $W$, $B'\subseteq A$ and $B'\subseteq B$. Then
$\ontoIIix_v(A,B)\le\len(\ontoIIloc vA/\ontoIIloc v{(A\cap B)})\le\len(\ontoIIloc vA/\ontoIIloc v{B'})=\ontoIIix_v(A,B')$
(the first inequality because lengths are non-negative, the second because $B'\subseteq A\cap B$, so that the middle quotient is a
quotient of $\ontoIIloc vA/\ontoIIloc v{B'}$, which has finite length as $A$ and $B'$ are lattices in $W$), which gives the claim.
\end{proof}

\begin{lemma}[Localising the run conclusion]\label{pf:onto2:l2}
Let $X$ be a record, $M_0=\mathrm{Mlat}(\kappa)$ a lattice in $\mathrm{Cusp}(\kappa)$, $\Theta_0=\iota_\kappa^{-1}(\Theta)$, and assume
$\ontoIICl(X,S,\kappa,w_L,w_F,T,\nu_s,\beta)$. Then there is $r_0\ge0$ on the places of $K_0$ with $\sum_\lambda
r_0(\lambda)\ell_\lambda\le\beta$ such that for every $\lambda$: $M_0\cap\Theta_0$ spans $\mathrm{Cusp}(\kappa)$, and for every integer
$e_0\ge T\nu_s(\lambda)$ and every $K_0$-subspace $W_0$,
\[
\ontoIIix_\lambda\bigl(M_0\cap W_0,M_0\cap\Theta_0\cap W_0\bigr)\le e_0\dim W_0+r_0(\lambda).
\]
\formallabel{Formalism.MainNotions.Gaps.LocalGlobalPassage.l2}
\end{lemma}

\begin{proof}
Take $r_0$ from the run conclusion. Fix $\lambda$, $e_0\ge T\nu_s(\lambda)$ and $W_0$, and let $V'=\widehat{\iota W_0}$. Since $M_0$
spans, $\iota W_0$ lies in the $K_0$-span of $M=\iota M_0$, so $V'$ lies in the span of $\widehat M_\lambda$ by fact (iii). By the run
conclusion, $(\widehat M_\lambda\cap V')/(\widehat\Theta_\lambda\cap V')$ has finite length $\le e_0\dim V'+r_0(\lambda)$, and $\dim
V'=\dim W_0$ by fact (ii). By fact (i) (with $A=M$, $T=\Theta$, $W=\iota W_0$) and the injectivity of $\iota$, which identifies
$M\cap\iota W_0$ with $M_0\cap W_0$ and $M\cap\iota W_0\cap\Theta$ with $M_0\cap\Theta_0\cap W_0$, the length of
$\ontoIIloc\lambda{(M_0\cap W_0)}/\ontoIIloc\lambda{(M_0\cap\Theta_0\cap W_0)}$ is finite and at most $e_0\dim W_0+r_0(\lambda)$; as
$M_0\cap\Theta_0\cap W_0\subseteq M_0\cap W_0$, this length is the local index. Taking $W_0=\mathrm{Cusp}(\kappa)$ and an integer
$e_0\ge T\nu_s(\lambda)$, the finiteness of $\ontoIIloc\lambda{M_0}/\ontoIIloc\lambda{(M_0\cap\Theta_0)}$ shows that $M_0$ lies in the
$K_0$-span of $M_0\cap\Theta_0$ (a module of finite length over the discrete valuation ring $\OO_{K_0,(\lambda)}$ is torsion, so
every element of $M_0$ has a non-zero multiple in $\ontoIIloc\lambda{(M_0\cap\Theta_0)}$), so $M_0\cap\Theta_0$ spans
$\mathrm{Cusp}(\kappa)$.
\end{proof}

\begin{lemma}[Remainders from the run conclusion]\label{pf:onto2:rem}
Let $X$ be a record, $z:K^n\cong K\otimes_{K_0}\mathrm{Cusp}(\kappa)$ a $K$-linear isomorphism such that
$\langle\mathrm{Mlat}(\kappa)\rangle_z$ is a lattice in $K^n$, and assume $\ontoIICl(X,S,\kappa,w_L,w_F,T,\nu_s,\rho_X abcd)$ where
$\rho_X$ is the $K_0$-rank of the span of $\mathrm{Mlat}(\kappa)$ and $a,b,c,d\in\R$. Then there is $r:\{v\}\to\R_{\ge0}$ on the
places of $K$ with $\sum_{v\in S'}\ell_vr(v)\le n\,abcd$ for every finite $S'$, such that for every place $v$ of $K$, every $e\in\N$
with $T\nu_s(\lambda_v)\le e$ and every $K$-subspace $W\subseteq K^n$,
\[
\ontoIIix_v\Bigl(\ontoIIcut_W\langle\mathrm{Mlat}(\kappa)\rangle_z,\ \ontoIIcut_W\langle\{F:\iota_\kappa F\in\Theta\}\rangle_z\Bigr)
\le e(v|\lambda_v)\,e\,\dim W+r(v).
\]
\formallabel{GlobalTheory.BigOn.exists_rem_of_concl}
\end{lemma}

\begin{proof}
Since $\langle\mathrm{Mlat}(\kappa)\rangle_z$ is a lattice, $M_0=\mathrm{Mlat}(\kappa)$ is a lattice in $\mathrm{Cusp}(\kappa)$
(\lref{GlobalTheory.BigOn.isLattice_of_spanK}), so $\rho_X=\dim_{K_0}\mathrm{Cusp}(\kappa)=n$ because $z$ is an isomorphism. Let
$r_0$ be given by \cref{pf:onto2:l2} and put $r(v)=e(v|\lambda_v)r_0(\lambda_v)\ge0$. By \cref{pf:onto2:l4}, $\sum_{v\in
S'}\ell_vr(v)\le\sum_\lambda r_0(\lambda)\ell_\lambda\le n\,abcd$. Let $\iota:\mathrm{Cusp}(\kappa)\to K^n$, $F\mapsto z^{-1}(1\otimes
F)$; it sends independent families to independent families and its image spans $K^n$
(because $F\mapsto1\otimes F$ maps a $K_0$-basis of $\mathrm{Cusp}(\kappa)$ to a $K$-basis of $K\otimes_{K_0}\mathrm{Cusp}(\kappa)$ and $z$ is
an isomorphism), and
$\mathrm{up}(A_0)=\langle A_0\rangle_z$. Apply \cref{pf:onto2:l5} with $M_0$, $N_0=M_0\cap\Theta_0$ (which spans by \cref{pf:onto2:l2}),
$e_0=e$, $\rho=r_0(\lambda_v)$, $E=e(v|\lambda_v)e$, and $T_h=\langle\{F:\iota_\kappa F\in\Theta\}\rangle_z\supseteq\mathrm{up}(N_0)$;
the hypothesis of \cref{pf:onto2:l5} is the conclusion of \cref{pf:onto2:l2}.
\end{proof}

\begin{proposition}[The index bound]\label{pf:onto2:index}
For every $C_K$ and all data as above: if the record relation, the core ties, the own-prime pin, the run bounds $\ontoIIRB(C_K,r)$,
\ref{pf:onto2:L1}, \ref{pf:onto2:L2} and $O.b=b(L)$ hold, then $O$ satisfies \ref{pf:onto2:L5}.
\formallabel{GlobalTheory.BigOn.hasIndexBound_of_runs}
\end{proposition}

\begin{proof}
By the record relation the face weights are the run weights, and the depth tie gives $E(v,i)=e(v|\lambda_v)\,E^{*}(\lambda_v,i)$
(\cref{pf:onto2:depdat,pf:onto2:runtie}). Let $q$ be a truth value and $i$ an index. We claim that there is $\rho_{q,i}\ge0$ on the
places of $K$ with $\sum_{v\in S'}\ell_v\rho_{q,i}(v)\le n_Z(i)\,fmb\,\delta_{e2}$ for all finite $S'$, such that
\eqref{eq:pf:onto2:LIB} holds with $T_{h1}=T_h$ and remainder $\rho_{q,i}(v)$ at every place $v$ with
$(p_v=p_{\mathrm{de}})\Leftrightarrow q$. If there is no such place, $\rho_{q,i}=0$ works ($\delta_{e2}>0$, $b>0$). Otherwise fix such a
place $v_1$. All places $v$ with $(p_v=p_{\mathrm{de}})\Leftrightarrow q$ have the same run $r_v=r_{v_1}$ and the
same prime $\mathrm{own}(v)=\mathrm{own}(v_1)$. By the run bounds at $v_1$, the run conclusion holds for the record of this run with total
\eqref{eq:pf:onto2:beta}. The form tie at $(v_1,i)$ identifies $M_s(i)$ with $\langle\mathrm{Mlat}\rangle_z$, which is a lattice by
\ref{pf:onto2:L1}. \Cref{pf:onto2:rem} (with $(a,b,c,d)=(f,m,b,\delta_{e2})$, transported through the form tie) gives $\rho_{q,i}$ with the sum bound and, at every place $v$ where the form tie holds for
this run and prime (in particular every $v$ with $(p_v=p_{\mathrm{de}})\Leftrightarrow q$), the bound
$\ontoIIix_v(\ontoIIcut_WM_s(i),\ontoIIcut_WT_h(v,i))\le e(v|\lambda_v)e\dim W+\rho_{q,i}(v)$ for every integer $e\ge
T_O\,\nu_s(\lambda_v)$, where $\nu_s$ are the local depths of the run. By
\lref{GlobalTheory.RunDepthScale.nusOf_eq_ram_mul_nuStarAt}, $\nu_s(\lambda)=e_\lambda\nu^*(\lambda)$ with $\nu^*$ the normalised
depth of \cref{pf:onto2:depdat} and $e_\lambda$ the absolute ramification index of $\lambda$, so $T_O\nu_s(\lambda)\le\lceil e_\lambda
T_O(\nu^*(\lambda)+[p_\lambda\ \text{bad}])\rceil=E^{*}(\lambda)$ (\lref{GlobalTheory.BigOn.nusR_le_eLam}). Taking
$e=E^{*}(\lambda_v,i)$ proves the claim.

Now take $T_{h1}=T_h$, $r_1(v,i)=\rho_{\mathrm{false},i}(v)$, $r_2(v,i)=\rho_{\mathrm{true},i}(v)$, $\delta_{e2}$ the thick density of $L$
(which lies in $[0,1]$), and $a(v,i)\in K^\times$ with $v(a(v,i))=E(v,i)$
(a power of an element of $v\setminus v^2$). Then $r(v,i)=r_2(v,i)$ if $p_v=p_{\mathrm{de}}$ (recall that the
record relation gives $O.p_{\mathrm{de}}=p_{\mathrm{de}}$) and $r_1(v,i)$ otherwise, so \eqref{eq:pf:onto2:LIB} holds at all places;
$T_h(v,i)$ is a lattice at pairs by \ref{pf:onto2:L2}; and the sum bounds hold because $f$, $m$ and $b$ in the claim are those of $O$
($O.f=D.R.f$, $O.m=D.R.m$ by the depth tie and $O.b=b(L)$ by hypothesis). This is \ref{pf:onto2:L5}.
\end{proof}

\subsubsection{The depth leaves}\label{sec:pf:onto2:depth}

The depth leaves bound the normalised exponents $e_0(v,j)$ of a record object, whose numerators $E(v,j)$ and
$v_v(\sigma(v,j))$ are fixed by the \emph{depth tie} in terms of local depths of integral points of the moduli spaces of the
core.

\begin{definition}[Depth data]\label{pf:onto2:depdat}
Let $T$ be a core of $D$, with host $H$ (a CM field $E$ with complex conjugation $a\mapsto\bar a$ and primes $d,q$), level $M$,
coefficient field $c.F$ and, for each index $j$, a moduli space $X_j$ with datum $D_{w,j}$ (\lref{Formalism.L2.Pj.Par}). Its
\emph{depth datum} (\lref{GlobalTheory.RunDepth.depOf}, \lref{Formalism.L2.Pj.depRun}) consists of the integer $T=T_O$,
the primes $g_1,g_2$ (the level primes of the host), $q_1=q$, $d$, the finite set $s_m$ of small primes of the host, the ranks
$f,r$, a function $v_d$ on primes, weight vectors $a_L(j)$ (those of the run weight $lR(j)$) and $a_F(j)$ with degrees
$\mathrm{DF}_L(j)$, $\mathrm{DF}_F(j)$, integers $n_W(j)$, $D_\pm(j)$, and elements $\pi_1,\pi_2\in\OO_E$. A prime is
\emph{ramified} if it lies in $s_m\cup\{q_1,d\}$ and \emph{bad} if it is ramified or a level prime. For a place $\lambda$ of $c.F$
with residue characteristic $p_\lambda$ and absolute ramification index $e_\lambda$:
\begin{itemize}
\item the \emph{local depth} $\nu_W(\lambda,j)\in[0,\infty]$ (\lref{Formalism.L2.Pj.nuW}) is $\max(0,\cdot)$ of the supremum,
  over complete discrete valuation rings $T'$ over $\lambda$ and integral algebraic frame points $y$ of $X_j$ over $T'$, of the
  infimum over integral charts at $y$ and over admissible test polynomials and integral power series of the normalised
  valuations of the associated jet values of weight $a_L(j)$;
\item the \emph{normalised depth} is $\nu^*(\lambda,j)=\mathrm{DF}_L(j)/(g-1)$ if $p_\lambda=g\in\{g_1,g_2\}$, and
  $\nu^*(\lambda,j)=\nu_W(\lambda,j)/e_\lambda+c_P(p_\lambda)$ otherwise, where $c_P(p)=\lceil v_d(p)\,\mathrm{DF}_L(j)\rceil$ for
  ramified $p$ and $c_P(p)=0$ otherwise (\lref{Formalism.L2.Pj.nuStarAt});
\item $E^{*}(\lambda,j)=\bigl\lceil e_\lambda T\bigl(\nu^*(\lambda,j)+[p_\lambda\ \text{bad}]\bigr)\bigr\rceil$.
\end{itemize}
For a place $v$ of $K$ above the place $\lambda_v$ of $c.F$ put $E(v,j)=e(v|\lambda_v)\,E^{*}(\lambda_v,j)$
(\lref{Formalism.L2.Pj.EAt}), $\mathrm{nn}_K(j)=\prod_{i=1,2}\bar\pi_i^{D_+(j)}\pi_i^{D_-(j)}$ and
$\sigma(j)=\mathrm{nn}_K(j)\cdot N_\sigma(j)$ with the positive integer $N_\sigma(j)=(g_1g_2)^{fr(8fr+4)}\prod_{p\
\text{ramified}}p^{Tc_P(\mathrm{DF}_L(j),p)+c_P(\mathrm{DF}_F(j),p)}$ (\lref{Formalism.L2.Pj.sigAt}).
\end{definition}

\begin{definition}[Depth tie and content tie]\label{pf:onto2:runtie}
The \emph{depth tie} of a record object $O$ with the core $T$ and an index map $\iota$ ($\iota(j)=j+1$ in the applications)
(\lref{GlobalTheory.RunDepth.RunTie}) says: the indices $\iota(j)$ are admissible; the involution of $\OO_E$ is the host
conjugation (also after mapping to $K$); $E(v,j)$, $\sigma(v,j)$, $\mathrm{nn}_K(j)$ of $O$ are those of \cref{pf:onto2:depdat} at
$\iota(j)$; $T=T_O$; $(g_1,g_2,q_1,d,s_m)$ are the host primes; $k,m,f,x,\varepsilon,S_c,\eta,b',X_2,v_d$ are the corresponding
parameters of $D$ and $L$; $\mathrm{wt}(j)=2fb$; $\theta(\ell)=\mathrm{th}_k(\ell)$ and $\alpha(\ell)=0$; and $D_F\ge\deg lR(\iota(j))$.
The numbers $N$ satisfy the \emph{content tie} if $F_0(w,j)=F_{\iota(j)}(p_w)$ for all $w,j$, where for the numerical datum
$\mathrm{num}(D,L)$ of $(D,L)$ (\lref{GlobalTheory.RunDepth.numOf}) the content of a prime $p$ is
(\lref{Formalism.L2.DIV.NumDatum.FjP})
\[
F_j(p)=\begin{cases}
-\tfrac12\bigl(D_+(j)+D_-(j)\bigr)&p\in\{g_1,g_2\},\\
-\tfrac{k-1}{2}\sum_{t}\min\bigl(w_j(0,t),w_j(1,t)\bigr)&p=d,\\
F^{\mathrm{in}}_j(p)&p\ \text{inert or signal and}\ p\le X_2,\\
0&\text{otherwise},
\end{cases}
\]
with the weights $w_j$ of the datum and an explicit orbit average $F^{\mathrm{in}}_j(p)$.
\end{definition}

\begin{lemma}[Rounding inequalities]\label{pf:onto2:round}
Let $e_1,e_2\ge1$ be integers; all other letters denote real numbers.
\begin{enumerate}[label=(\roman*)]
\item If $E\in\N$ and $E\le e_1\lceil e_2x\rceil$, then $E/(e_1e_2)<x+1/e_2$.
\item If $g\ge3$, $T\ge0$, $\nu\le D/(g-1)$, $TD\le D'$ and $D'<(g-2)/2$, then $T\nu<1/2$.
\item If $e_2\ge2$, $T\ge1$, $T\nu<1/2$, $s\ge0$ and $X<T(\nu+1)+1/e_2$, then $X-(n+s)\le2T-n$.
\item If $l>1$, $T\ge0$, $\nu\le D/(l-1)$, $s'\ge0$, $TD\le D'\le CkB$, $Ck\le l-1$, $B\ge0$ and $X<T(\nu+s+1)+1$, then
  $X-(Ts+s')\le B+T+1$.
\end{enumerate}
\formallabel{Formalism.L2.Indef.Tree.DepthCond.rounding_condition_holds}
\formallabel{Formalism.L2.Indef.Tree.DepthCond.nu_level_condition_holds}
\formallabel{Formalism.L2.Indef.Tree.DepthCond.level_condition_holds}
\formallabel{Formalism.L2.Indef.Tree.DepthCond.ramification_condition_holds}
\end{lemma}

\begin{proof}
(i) $\lceil e_2x\rceil<e_2x+1$, so $E<e_1(e_2x+1)$; divide by $e_1e_2$. (ii) $T\nu\le TD/(g-1)\le D'/(g-1)<1/2$ because
$2D'<g-2<g-1$. (iii) $1/e_2\le1/2$, so $X<T\nu+T+1/2<T+1\le2T$, and $s\ge0$. (iv) $T\nu\le TD/(l-1)\le D'/(l-1)\le B$ since
$D'\le CkB\le(l-1)B$; hence $X-(Ts+s')<T\nu+T+1\le B+T+1$.
\end{proof}

\begin{lemma}[Ramification at level primes]\label{pf:onto2:ramlevel}
Let $\lambda$ be a place of the coefficient field $c.F$ of a core whose residue characteristic $p$ divides the level $M$. Then
$e_\lambda\ge p-1$.
\formallabel{GlobalTheory.RunDepth.resChar_sub_one_le_ramOverZ_of_dvd_level}
\end{lemma}

\begin{proof}
The coefficient field contains a primitive $M'$-th root of unity for $M'=\mathrm{lcm}(M,A)$ (part of the definition of coefficient
fields), hence a primitive $p$-th root of unity $\zeta$ (\lref{Formalism.L2.Pj.zeta_mem_of_dvd}). The prime $p$ is totally
ramified in $\Q(\zeta)$ with ramification index $p-1$ \cite[Lemma~1.4]{Washington1997}, and ramification indices are
multiplicative in towers, so $e_\lambda=e(\lambda|\mathfrak p)(p-1)\ge p-1$ for the place $\mathfrak p$ of $\Q(\zeta)$ below $\lambda$
(\lref{GlobalTheory.RunDepth.resChar_sub_one_le_ramOverZ_of_primitiveRoot}).
\end{proof}

In the following four propositions we fix a core $T$ with its laws, an error datum, parameters, a number field $K$ over $c.F$ and
$E$, an indefinite datum $S$ over $K$, a finite group $P$, constants $c$, a record object $O$, an index map $\iota$ and numbers
$N$, and assume the depth tie, \ref{pf:onto2:P1}, \ref{pf:onto2:P2}, $3\le x$, $kb'\le b$ and the content tie for
$\mathrm{num}(D,L)$. For $v$ above $\lambda_v$ we write $e_1(v)=e(v|\lambda_v)$ and $e_2(v)=e_{\lambda_v}$, so that $e_v=e_1(v)e_2(v)$
(\lref{Formalism.L2.Indef.Tree.DepthTie.ramw_eq_below}). By \cref{pf:onto2:depdat}, for every $v$ and $j$,
\begin{equation}\label{eq:pf:onto2:tieE}
E(v,j)=e_1(v)\bigl\lceil e_2(v)\,T\bigl(\nu^*(\lambda_v,j)+[p_v\ \text{bad}]\bigr)\bigr\rceil
\end{equation}
(\lref{Formalism.L2.Indef.Tree.DepthTie.tieE}); here $p_v=p_{\lambda_v}$, as $v$ lies above $\lambda_v$.

\begin{proposition}[Depth at the level primes]\label{pf:onto2:dlevel}
Under these assumptions $O$ satisfies \ref{pf:onto2:L9} for $N$.
\formallabel{GlobalTheory.RunDepth.leafDepthLevel_run}
\end{proposition}

\begin{proof}
Let $w\in\ontoIITw(S)$ with $p_w=g\in\{g_1,g_2\}$, $v\in\{w,\bar w\}$ and $j$. Since $2\in s_m$ and $g\notin s_m$, $g\ge3$. Put
$\mathrm{DF}=\mathrm{DF}_L(\iota j)$ and $D_k=\deg\kappa R(\iota j)$. (a) $\nu^*(\lambda_v,j)=\mathrm{DF}/(g-1)$
(\lref{Formalism.L2.Indef.Tree.DepthTie.nuLevel}). (b) $T\,\mathrm{DF}\le D_k$ and $D_k<(g-2)/2$: the first by the degree identity
for weight exponents (\lref{GlobalTheory.WeightExponents.weight_exponents}: the degree of a weight of the datum equals the degree
of the associated weight on the embeddings) and the row bound $T_O\deg lR\le\deg\kappa R$
(\lref{GlobalTheory.RunDepth.mul_rows_lR_le_kaR}), the second by the page bound $\deg\kappa R<(g_i-2)/2$
(\lref{GlobalTheory.RunDepth.DF_kaR_le}, valid for $3\le x$ and $kb'\le b$). (c) By \cref{pf:onto2:ramlevel}, $e_2(v)\ge g-1\ge2$, since $g$ divides the
level (\lref{Formalism.L2.Pj.level_M}). (d) $v_v(\sigma(v,j))/e_v=\mathrm{nn}(v,j)+s$ with $s=v_v(N_\sigma)/e_v\ge0$, because
$\sigma=\mathrm{nn}_K\cdot N_\sigma$ with $N_\sigma$ a non-zero integer and $\mathrm{nn}_K\ne0$
(\lref{Formalism.L2.Indef.Tree.DepthTie.sigmaLevel}). (e) $T\ge1$ (a law of the depth datum). By \cref{pf:onto2:round}(ii) with
$D=\mathrm{DF}$, $D'=D_k$, $T\nu^*<1/2$; by \eqref{eq:pf:onto2:tieE} ($g$ is bad) and \cref{pf:onto2:round}(i), $X:=E(v,j)/e_v<T(\nu^*+1)+
1/e_2(v)$; and \cref{pf:onto2:round}(iii) gives $e_0(v,j)=X-(\mathrm{nn}(v,j)+s)\le2T-\mathrm{nn}(v,j)$.

For the content identity: the involution of $K$ restricts to the host conjugation, so $e_{\bar w}=e_w$ and
$v_{\bar w}(\mathrm{nn}_K)=v_w(\overline{\mathrm{nn}_K})$, whence $\mathrm{nn}(w,j)+\mathrm{nn}(\bar w,j)=v_w(\mathrm{nn}_K
\overline{\mathrm{nn}_K})/e_w$. By the splitting law of the depth datum, $\pi_i\bar\pi_i=g_i$, so
$\mathrm{nn}_K\overline{\mathrm{nn}_K}=(g_1g_2)^{D_+(\iota j)+D_-(\iota j)}$, whose valuation at $w$ is $e_w(D_+(\iota j)+D_-(\iota j))$
because $p_w$ is exactly one of the distinct primes $g_1,g_2$. Hence $\mathrm{nn}(w,j)+\mathrm{nn}(\bar w,j)=D_+(\iota j)+D_-(\iota j)$
(\lref{Formalism.L2.Indef.Tree.DepthTie.nnLevel}, using that $D_\pm$ of the depth datum are those of the numerical datum,
\lref{GlobalTheory.T.depthParams_plusMinusCounts_eq_kappa}). Since $F_0(w,j)=F_{\iota j}(g)=-\tfrac12(D_++D_-)$ by the content tie
(\lref{Formalism.L2.Indef.Tree.DepthCond.content_prime_condition_holds}), $\mathrm{nn}(w,j)+\mathrm{nn}(\bar w,j)=-2F_0(w,j)$
(\lref{Formalism.L2.Indef.Tree.DepthCond.leafDepthLevel_of}).
\end{proof}

\begin{proposition}[Depth at the ramified primes $q$ and $d$]\label{pf:onto2:dram}
Under these assumptions $O$ satisfies \ref{pf:onto2:L7} for $N$.
\formallabel{GlobalTheory.RunDepth.leafDepthRam_run}
\end{proposition}

\begin{proof}
Let $w\in\ontoIITw(S)$ with $p_w=l\in\{q_1,d\}$, $v\in\{w,\bar w\}$, $j$, and $\mathrm{DF}=\mathrm{DF}_L(\iota j)$. Write $\nu^*=\nu_W/e_2+s_W$
with $s_W=c_P(l)$ (\cref{pf:onto2:depdat}; $l$ is not a level prime) and $v_v(\sigma(v,j))/e_v=Ts_W+s'$. (a) $s'\ge0$, because the
factor $l^{Tc_P(\mathrm{DF}_L,l)+c_P(\mathrm{DF}_F,l)}$ of $N_\sigma$ contributes at least $Ts_W$ and the remaining factors have
non-negative valuation (\lref{Formalism.L2.Indef.Tree.DepthTie.sprimeRam}). (b) The crude depth bound: integral charts exist and
are saturated (\lref{GlobalTheory.RunDepth.laws_and_nuFinAt}, \lref{Formalism.L2.Pj.SaturatedCharts.saturation_charts_holds}),
and the weights are dominant with enough variables, so $\nu_W(\lambda_v,j)/e_2\le\mathrm{DF}/(l-1)$
(\lref{Formalism.L2.Pj.nuW_le_crude}, \lref{Formalism.L2.Indef.Tree.DepthTie.crudeRam}). (c) $T\,\mathrm{DF}\le D_k\le k\,s_B(j)$
with $D_k=\deg\kappa R(\iota j)$ (\lref{GlobalTheory.WeightExponents.weight_exponents}, \lref{GlobalTheory.RunDepth.mul_rows_lR_le_kaR},
\lref{GlobalTheory.RunDepth.DF_kaR_le}). (d) $k\le l-1$ for $l\in\{q,d\}$, because $q$ and $d$ exceed $k$
(\lref{GlobalTheory.RunDepth.k_le_sub_one_of_lt}). (e) $s_B(j)\ge0$ by \ref{pf:onto2:P1}, \ref{pf:onto2:P2}. As $l$ is bad,
\eqref{eq:pf:onto2:tieE} and \cref{pf:onto2:round}(i) give $X:=E(v,j)/e_v<T(\nu_W/e_2+s_W+1)+1$, and \cref{pf:onto2:round}(iv) (with
$\nu=\nu_W/e_2$, $s=s_W$, $C=1$, $B=s_B(j)$) gives $e_0(v,j)=X-(Ts_W+s')\le s_B(j)+T+1$.

Content: by the content tie, $F_0(w,j)=F_{\iota j}(l)$. If $l=q$, then $l$ is not a level prime and is ramified in the host, so
$F_{\iota j}(q)=0$; if $l=d$ then $F_{\iota j}(d)=-\frac{k-1}2\sum_t\min(w(0,t),w(1,t))$
(\lref{Formalism.L2.Indef.Tree.DepthCond.content_prime_condition_holds}). The weights satisfy $w(\cdot,t)\ge0$ and
$\min(w(0,t),w(1,t))\le2b'$ (\lref{GlobalTheory.RunDepth.sizes_le_two_mul}); summing over the $n=f/2$ values of $t$ gives
$-n(k-1)b'\le F_{\iota j}(d)\le0$ (\lref{Formalism.L2.Indef.Tree.DepthCond.ramified_discriminant_condition_holds}), and
$n(k-1)b'\le2fkb'$.
\end{proof}

\begin{proposition}[Depth at the small primes]\label{pf:onto2:d2}
Assume moreover the formal-group statements $(\mathrm A_p)$ and $(\mathrm H_p)$ below for all primes $p$. Under these assumptions $O$
satisfies \ref{pf:onto2:L8} for $N$.
\formallabel{GlobalTheory.RunDepth.leafDepth2_run}
\end{proposition}

Here, for a complete noetherian local domain $T'$ of characteristic $0$ in which $p$ is not a unit, a commutative group object over
$T'$ with a degeneration datum and a uniformisation, and a jet chart with formal group law $F$ of dimension $d$:
$(\mathrm H_p)$ (\lref{GlobalTheory.ChF.hodgeTate_inclusion_chart}) says that over the valuation ring $\OO_C$ of every
algebraically closed valued field $C$ of characteristic $0$ above $T'$ there are finitely many characters $u_k$ of $F$ such that
$c\cdot f$ lies in the span of the logarithmic differentials $du_k$ for all $f\in\OO_C^d$ and all $c$ with $|c|^{p-1}\le|p|$; and
$(\mathrm A_p)$ (\lref{GlobalTheory.ChF.colength_bound_split_place_chart}) says that, in the presence of a compatible action of
the integers of a CM field and eigenbases $W_j$ of this action at a place over $p$, suitable invariants $x_j$ of the eigenbases
satisfy $\prod_jx_j\ge|p|^{n r_m(1/(p-1)+v_d/2)}$. Both are proved in the development for every $p$.

\begin{proof}
Let $\lambda$ be a place of $c.F$ with $p=p_\lambda\in s_m$ and $j$ an index. By
\lref{GlobalTheory.CharDepth.twoPlaces_host}, $d$ splits at $p$, the places of the host over $p$ satisfy the
splitting condition with parameter $v_{dH}(p)$, and $0\le v_{dH}(p)\le v_{d2}$. By \lref{GlobalTheory.T.unscaledDepth_le_placeOverTwo}
(a consequence of $(\mathrm A_p)$ and $(\mathrm H_p)$), for dominant weights $a$ with $a(\tau,0)\le A_p$ on the CM type $\Psi$,
\[
\nu_W(\lambda,j)\le e_\lambda\Bigl(f\,q_D\,(1+v_{dH}(p)/2)A_p+q_D\sum_{\tau\notin\Psi}a(\tau,0)\Bigr)=:e_\lambda B ,
\]
with $q_D$ the second signature entry of the datum; and by \lref{GlobalTheory.T.sizes_bound_at_two_first_rows_weights} there is
$A_p$ with $T\,B\le Cm(2+v_{d2})(2fb+f(S_c+\eta))=s_C(j)$ (recall $\mathrm{wt}(j)=2fb$ and $v_d=v_{d2}$ by the depth tie, and $C\ge1$).
Now let $w\in\ontoIITw(S)$ with $p_w\in s_m$ and $v\in\{w,\bar w\}$. As $p_w$ is bad, \eqref{eq:pf:onto2:tieE} and
\cref{pf:onto2:round}(i) give $E(v,j)/e_v<T(\nu_W/e_2+c_P+1)+1/e_2$, while $v_v(\sigma)/e_v\ge Tc_P$ as in \cref{pf:onto2:dram}(a)
(\lref{GlobalTheory.RunDepth.sprime_nonneg_two}); hence $e_0(v,j)<TB+T+1\le s_C(j)+T+1$. Finally $F_0(w,j)=F_{\iota j}(p_w)=0$ since
$p_w\in s_m$ is neither a level prime nor $d$ (\lref{GlobalTheory.RunDepth.leafDepthTwo_of_unscaled_bound}).
\end{proof}

\paragraph{Depth packages.} The bound at good primes is proved in three steps: at an integral point over a good place, the jets
of an integral chart provide \emph{depth packages} (\cref{pf:onto2:packages}); every package bounds the depth by a content-type
quantity plus an error (\cref{pf:onto2:nugood}); and the resulting bound for the local depth is converted into the leaf
\ref{pf:onto2:L6} (\cref{pf:onto2:leafgood,pf:onto2:leafgoodrun}). The input from the theory of formal groups is the statement
$(\mathrm U_p)$ (\lref{GlobalTheory.ChF.unramified_companion_chart_characters}), the analogue of $(\mathrm A_p)$ at places of the
CM field unramified over $p\ne2$, $p$ prime to the discriminant; it is proved in the development for every $p$.

Throughout this paragraph, $E$ is a CM field with a host structure $H$, $D_w$ a datum for $H$ with signature $(p,q)$, $M$ a level,
$X$ a moduli space with coefficient field $c.F$, and the involution of $\OO_E$ is the host conjugation. A \emph{good situation}
$(N,j,\mathrm{pl},a,\lambda,T',y,C,n_W,\nu)$ consists of a numerical datum $N$ (with ranks $k,m,r=k+m-1,f$, weights $w_j$ and
content $F_j$), an index $1\le j\le m$ with $(p,q)=\sigma_W(j)$, an injective map $\mathrm{pl}$ from the places of $N$ to the CM type
of $H$ satisfying the place tie (\lref{GlobalTheory.GoodDepth.PlaceTie}: $E/\Q$ is abelian, the primes outside $s_m\cup\{q,d\}$
are unramified in $E$, and the Frobenius congruences hold at inert primes), a weight $a$ dominant for $D_w$, a place $\lambda$ of
$c.F$ whose residue characteristic $\ell$ is not in $s_m\cup\{q,d\}$, does not divide $M$, and is at least $5$ if it is inert, a
complete discrete valuation ring $T'$ over $\lambda$ with ramification index $e$, an integral algebraic frame point $y$ of $X$
over $T'$, an integral chart $C$ at $y$, an integer $n_W$, and a real $\nu$ with $e_\lambda\nu\le\nu_{\mathrm{loc}}(C)$, where
$\nu_{\mathrm{loc}}(C)=\nu_{\mathrm{loc}}(C;e,M,a,n_W)$ is the depth of the chart (the infimum in the definition of the local depth,
\cref{pf:onto2:depdat}). The class of $\ell$ for $N$ is \emph{inert}, \emph{signal} or \emph{split}; it is \emph{of conjugate type}
if it is inert and the element $h_N(\ell)$ attached to it has odd additive order (then the orbit length $\mathrm{ord}_N(\ell)$ is
odd), and of \emph{type $\Phi$} if it is inert and $h_N(\ell)=0$ (a special case of conjugate type). If $\ell$ is inert or signal
and not of conjugate type, $\mathrm{ord}_N(\ell)$ is even.

\begin{definition}[Depth packages]\label{pf:onto2:pkg}
Let $P=(\mathrm{Pl},k,m,r,f,d,\mathrm{pos},w,S_E)$ be a \emph{row parameter}: a finite set $\mathrm{Pl}$ of places, integers, a
permutation-like map $\mathrm{pos}:\mathrm{Pl}\times\N\to\mathrm{Pl}$ whose orbits have length $d$, positive weights $w$ and $S_E\ge0$
(for $(N,j,\ell)$ one takes $P(N,j,\ell)$ with $d=\mathrm{ord}_N(\ell)$, $w=w_j$ and $S_E=S_c+\eta$). A set $Z\subseteq\mathrm{Pl}$ of
\emph{representatives} meets every orbit of $\mathrm{pos}$ exactly once. A \emph{row family} for $(P,\ell,d)$ consists of a field $C$ with a valuation $V$ of residue characteristic $\ell$
and, for every $z\in\mathrm{Pl}$, \emph{row data} $X_z$ over $(C,V)$ with $d$ \emph{types} and two \emph{groups} $g\in\{0,1\}$ of rows,
of size $k-1$ and second size $m$, with a flag $F_z$ enumerating the rows of each type; $\mathrm{rk}_z(g,s)$ denotes the number of
rows of group $g$ and type $s$. A \emph{gauge} is a function $\mathrm{ga}(z,g,s,k')\in\N$. For $n_W\in\N$ and $\nu\in\R$, the
\emph{valuation-block property} of $(Z,n_W,\mathrm{ga},\nu)$ says: for every choice, for $z\in Z$, of finite lists of \emph{blocks}
(polynomials in the rows of $X_z$ of three kinds: bi-blocks built from powers of maximal minors of two groups, character blocks
built from powers of one maximal minor, and monomial blocks), each attaining a valuation bound $c$ on the integral germs of
$X_z$, such that the total number of letters is at most $n_W$ and the total weights of the blocks equal $\mathrm{ga}$, the depth
$\nu$ is at most the sum of the bounds $c$. The \emph{degree hypothesis} $\mathrm{DegHyp}(P,W,T,\theta_0,\rho,C_1)$
for $W:\mathrm{Pl}\to\N$ is: $T\ge1$, $w>0$,
$(1-\theta_0)w_z\le TW_z\le w_z+C_1S_E$, $0\le\theta_0\le1$, $(k-1)\theta_0\le5m$, $0\le\rho$ and $T\rho\le C_1fkS_E$.

A \emph{package of inert type} with time $T$ and depth $\nu$ consists of a
row family, representatives $Z$ whose rows are alternating, a gauge $\mathrm{ga}$, weights $W$, $n_W$, $\theta_0$ and $\rho$ such that:
$\mathrm{ga}(z,g,s,k')\ne0$ implies $1\le k'\le\mathrm{rk}_z(g,s)$ for $z\in Z$; the main term $\mathrm{ga}(z,\mathrm{bg}(s),s,k-1)$ equals
$W$ along the orbit of $z$ (where $\mathrm{bg}(s)=0$ for even $s$ and $1$ otherwise); $\rho$ is the sum over $z\in Z$ of
$\sum_{g,s,k'}k'\mathrm{ga}(z,g,s,k')$ with the main terms omitted; $f\,i_0(k,\ell)(k-1)+2rf\le n_W$, where $i_0(k,\ell)$ is the least
odd $i\ge3$ with $k^{3/2}<\ell^i$; $\mathrm{DegHyp}(P,W,T,\theta_0,\rho,C_1)$; and the valuation-block property of
$(Z,n_W,\mathrm{ga},\nu)$. The other kinds are variants:
packages of \emph{type $\Phi$} have $d=2$, all rows alternating and $8rf\le n_W$; \emph{split} packages have all rows of group $0$ big
and $4rf\le n_W$; \emph{zero} packages (and their $d=2$ variant) have no main term and no degree hypothesis, a gauge supported in
$1\le k'\le\min(\mathrm{rk},K)$ with $\sum_{k'}\mathrm{ga}(z,g,s,k')\le s_0(\ell-1)$, and $2Ks_0|Z|\le n_W$ (resp.\ $4Ks_0f\le n_W$);
packages of \emph{conjugate type} (and their zero variant) are indexed by $Z\times\Z/2d$ and have
$f\bigl((i_0(k,\ell)+1)(k-1)+4r\bigr)\le n_W$.

The \emph{package statement} for $(X,c)$ says that in every good situation, for all $T,\theta_0,C_1$ with
$\mathrm{DegHyp}(P(N,j,\ell),W_a,T,\theta_0,\rho_a,C_1)$, where $W_a(z)=a(\mathrm{pl}\,z,k-2)-a(\mathrm{pl}\,z,k-1)$ and
$\rho_a=\deg(a)-(k-1)\sum_zW_a(z)$, and the respective size bound for $n_W$: (i) if $\ell$ is inert or signal and not of conjugate
type, there is a package of inert type for $P(N,j,\ell)$ with time $T$ and depth $\nu$; (ii) if $\ell$ is of type $\Phi$, a package
of type $\Phi$; (iii) if $\ell$ is split, a split package for some $(d,\mathrm{pos})$; (iv) if $\ell$ is not of conjugate type,
$w_j\ge0$ and $a(\tau,0)\le s_0(\ell-1)$ for all $\tau$ with positive multiplicity, a zero package with depth $\nu$; (v) if $\ell$ is of type $\Phi$, $w_j\ge0$ and $a(\tau,0)\le s_0(\ell-1)$ for all $\tau$ with positive
multiplicity, a zero package of the $d=2$ variant; (vi) if $\ell$ is inert and $\mathrm{ord}_N(\ell)$ is odd, a package of conjugate
type; (vii) if $\ell$ is inert, $\mathrm{ord}_N(\ell)$ is odd, $w_j\ge0$ and $a(\tau,0)\le s_0(\ell-1)$ for all $\tau$ with positive
multiplicity, a zero package of conjugate type. In (iv), (v) and (vii) the degree hypothesis is not assumed
(and $T$ plays no role); in (iii) it is assumed for $P(N,j,1)$.
\formallabel{Formalism.L2.FormalSeriesBounds.RF.RowFam}
\formallabel{Formalism.L2.FormalSeriesBounds.RF.DegHyp2}
\formallabel{Formalism.L2.FormalSeriesBounds.RF.PkgInert}
\formallabel{Formalism.L2.FormalSeriesBounds.RF.PkgPhiC}
\formallabel{Formalism.L2.FormalSeriesBounds.RF.PkgSplit}
\formallabel{Formalism.L2.FormalSeriesBounds.RF.PkgZero}
\formallabel{Formalism.L2.FormalSeriesBounds.RF.PkgZeroC}
\formallabel{UpperTheory.G2.PkgCs}
\formallabel{UpperTheory.G2.PkgZeroCs}
\formallabel{GlobalTheory.GoodDepth.AllPts}
\formallabel{GlobalTheory.GoodDepth.PkgInertSt}
\formallabel{GlobalTheory.GoodDepth.PkgSt}
\end{definition}

The gauges used below are the \emph{gaps} of the weights: for $r'\in\N$ and a sequence of natural numbers $a_0\ge a_1\ge\dots\ge
a_{r'-1}$ put $a_{r'}=0$ and $\mathrm{gap}_{k'}(a)=a_{k'-1}-a_{k'}$ for $1\le k'\le r'$, and $\mathrm{gap}_{k'}(a)=0$ otherwise.

\begin{lemma}[Gap identities]\label{pf:onto2:gaplem}
For $a$ as above and every $K\ge r'$,
\begin{equation}\label{eq:pf:onto2:gaps}
\mathrm{gap}_{k'}(a)\ne0\Rightarrow1\le k'\le r',\qquad \sum_{k'=i+1}^{K}\mathrm{gap}_{k'}(a)=a_i\ (i<r'),\qquad
\sum_{k'=1}^{K}k'\,\mathrm{gap}_{k'}(a)=\sum_{i<r'}a_i .
\end{equation}
\formallabel{GlobalTheory.RowAt.gap_sums}
\end{lemma}

\begin{proof}
The first assertion is the definition. For $i<r'$ the sum telescopes: $\sum_{k'=i+1}^{r'}(a_{k'-1}-a_{k'})=a_i-a_{r'}=a_i$, and the terms
with $k'>r'$ vanish. For the last identity write $k'=\#\{i:0\le i<k'\}$ and exchange the summations:
$\sum_{k'=1}^K k'\,\mathrm{gap}_{k'}(a)=\sum_{i=0}^{K-1}\sum_{k'=i+1}^{K}\mathrm{gap}_{k'}(a)=\sum_{i<r'}a_i$, since the inner sum is $a_i$
for $i<r'$ and $0$ for $i\ge r'$.
\end{proof}

\begin{theorem}[Depth packages at good primes]\label{pf:onto2:packages}
Assume $(\mathrm U_p)$ for all primes $p$. Then for every CM field $E$ with a host structure, every datum, level, moduli space $X$ and
coefficient datum $c$, the package statement for $(X,c)$ holds.
\formallabel{GlobalTheory.goodPrime_packages}
\formallabel{GlobalTheory.T.allGoodPrimeDepthDataExist_proof}
\end{theorem}

\begin{proof}
The package statement is the conjunction of (i)--(vii); they are proved in
\lref{GlobalTheory.T.inertPrimeDepthDataExistsAtPoints_proof} (i),
\lref{GlobalTheory.T.packages_exist_frobeniusComplexConjugation} (ii) and (v),
\lref{GlobalTheory.T.splitPrimeDepthDataExistsAtPoints_proof} (iii),
\lref{GlobalTheory.T.largePrimeDepthDataExistsAtPoints_proof} (iv) and
\lref{UpperTheory.exists_pkgCs_of_degree_bound_at_inert_primes} (vi) and (vii), all along the following steps. Fix a good situation.

\emph{Step 1: a valued field.} There are a field $C$ of characteristic $0$ with a valuation $V$ of residue characteristic $\ell$ and
an embedding of the fraction field of $T'$ into $C$ compatible with $\lambda$, in which the integers prime to $\ell$ (in particular
$M$) are units and to which the chart depth transfers: $\nu\le\nu_{\mathrm{loc}}(C;e_V,M,a,n_W)$
(\lref{GlobalTheory.PlaceClosedField.exists_closed_field_over_place}). Since $\ell\ne2$, the trace of the host idempotent is $1$
(\lref{GlobalTheory.QuadTrace.cmHost_idempotent_trace_eq_one}) and $M$ is a unit, the chart $C$ admits a base change $Y$ to a
formal chart over $C$ (\lref{GlobalTheory.ChartBaseChange.exists_baseChange}). By the place tie, $\ell$ is unramified in $E$.

\emph{Step 2: labelling the CM embeddings.} The Frobenius at $\ell$ permutes the embeddings $E\to\C$; the orbits are labelled by
representatives and positions: for inert $\ell$ by $Z\times\{0,1\}\times\Z/d$ with $d=\mathrm{ord}_N(\ell)$ even
(\lref{GlobalTheory.InertLabelling.main}), for split $\ell$ by a labelling attached to the Frobenius orbits on the CM type
(\lref{GlobalTheory.T.cmType_frobenius_orbits_label_at_split_prime}), for type $\Phi$ by $\mathrm{Pl}\times\Z/2$ via $z\mapsto
\mathrm{pl}(z)$ and its conjugate, and for odd order by $Z\times\Z/2d$. In each case the labelling $\tau(z,g,s)$ is a bijection
onto the embeddings, compatible with complex conjugation and with the Frobenius congruences $\tau(x)^\ell\equiv\tau'(x)$ modulo the
maximal ideal of $C$; distinct CM embeddings stay distinct modulo the maximal ideal
(\lref{GlobalTheory.EmbSep.cmEmbeddings_integral_and_distinct_mod_maximal'}), and every label has an idempotent at the
corresponding unramified place (\lref{GlobalTheory.exists_idempotent_of_unramified_place}). Moreover the \emph{rank
conditions} hold: $\mathrm{mult}(\tau)=p=k-1+j$ for embeddings of group $\mathrm{bg}(s)$ (which lie in the CM type) and
$\mathrm{mult}(\tau)=q=m-j$ for the others, since $(p,q)=\sigma_W(j)$.

\emph{Step 3: row data from jets.} For labellings by pairs of places (cases (i), (iii), (iv)), the jets of $Y$ supply row data
$X_z$ for the representatives, with flags, of sizes $(k-1,m)$, with $\mathrm{rk}_z(g,s)=\mathrm{mult}(\tau(z,g,s))$, alternating (inert
case) or big (split case), and with the valuation-block property of $(Z,n_W,\mathrm{ga}_a,\nu)$ for the gauge
\[
\mathrm{ga}_a(z,g,s,k')=\mathrm{gap}_{k'}\bigl(a(\tau(z,g,s),\cdot)\bigr)
\]
(\lref{GlobalTheory.T.rowData_supplied_at_pairs_of_places}); the row data extend from the representatives to all of
$\mathrm{Pl}$ (\lref{GlobalTheory.RowAt.supply_extend}). The proof of this supply uses the nondegeneracy of the formal
biextension of the chart modulo the maximal ideal (\lref{GlobalTheory.T.biextension_nondegenerate_mod_maximalIdeal_integral_chart}),
the tangent condition of the raw chart (\lref{GlobalTheory.T.rawChart_tangent_maximalIdeal_of_nonDegenerate}), the Cartier
hypothesis for integral gauged formal group laws (\lref{GlobalTheory.T.cartHyp_of_integralGaugedFormalGroupLaw}) and the
character bound at pairs of places, which rests on $(\mathrm U_p)$. For type $\Phi$ (cases (ii), (v)) the row family over
$\mathrm{Pl}\times\Z/2$ is supplied at self-conjugate places (\lref{GlobalTheory.PkgConj.exists_rowFam}, from
\lref{GlobalTheory.T.rowData_supplied_at_selfConjugate_places}), and for odd order (cases (vi), (vii)) over $Z\times\Z/2d$ by the
orbit-family restriction (\lref{UpperTheory.OrbitFamilyRestriction.exists_rowFamS}, whose main ingredient is the geometric
label set for even orbits, \lref{UpperTheory.labelSetGeo_of_even_of_rankOK_of_not_dvd_discr}); in both cases the ranks are the
multiplicities and the valuation-block property holds for the gauge of gaps.

\emph{Step 4: the gauge.} By \cref{pf:onto2:gaplem}, applied to the non-increasing sequence $a(\tau,\cdot)$ of length
$\mathrm{mult}(\tau)\le r$ (dominance of $a$): the gauge is supported in $1\le k'\le\mathrm{rk}$; for an embedding $\tau=\mathrm{pl}(z')$ of
the CM type, $\mathrm{mult}(\tau)=k-1+j>k-1$ and the main term is $\mathrm{gap}_{k-1}=a(\tau,k-2)-a(\tau,k-1)=W_a(z')$; and, since the
labelling is a bijection onto the embeddings,
\[
\sum_{z,g,s}\sum_{k'}k'\,\mathrm{ga}_a(z,g,s,k')=\sum_\tau\sum_{i<\mathrm{mult}(\tau)}a(\tau,i)=\deg(a),
\]
so that the sum with the main terms omitted is $\deg(a)-(k-1)\sum_zW_a(z)=\rho_a$ (\lref{GlobalTheory.RowAt.wtHyp_gap}). For the
zero packages, $\sum_{k'}\mathrm{ga}_a(z,g,s,k')=a(\tau,0)\le s_0(\ell-1)$ by the second identity of \cref{pf:onto2:gaplem}.

\emph{Step 5: assembly.} The degree hypothesis and the size bound for $n_W$ are hypotheses of the package statement, so Steps 1--4
give the package of each kind with time $T$ and depth $\nu$ (for zero packages $K=r$ and time $1$). For (iv), the labelling is the
inert labelling of Step 2 if $\ell$ is inert and not of conjugate type
(\lref{GlobalTheory.T.realPlaceLabelling_inert_frobenius_not_conj}) and the split labelling otherwise, and the package is
assembled from the labels by \lref{GlobalTheory.LargePkg.pkg_of_labels}, which performs Steps 3 and 4.
\end{proof}

\begin{theorem}[The depth bound at good places]\label{pf:onto2:nugood}
For every $C_1\ge0$ there is $C_c>0$ with the following property. Let $(E,H,D_w,M,X,c)$ be as above, admitting integral charts at
all integral algebraic frame points and satisfying the package statement. Let $N,j,\mathrm{pl},a$ be as in a good situation, $n_W,T\in
\N$ and $b,L,w\in\R$ such that the good-page conditions hold (\lref{GlobalTheory.GoodDepth.GoodPage}: the degree hypothesis for
$P(N,j,1)$ and $(W_a,\rho_a)$ with $C_1$, $Ta(\tau,0)\le3b$ for $\tau$ of positive multiplicity,
$10r(f+f^2(k-1)^3L^2/m)\le n_W$, $3\le L$, $\log k\le L$, $\ell\ge5$ for inert primes, and $w_j(z')\le X''w_j(z)$), and
$\sum_zw_j(z)\le w$, $b\le w$. Then for every place $\lambda$ of $c.F$ whose residue characteristic $\ell$ is not in
$s_m\cup\{q,d,g_1,g_2\}$ and does not divide $M$: $\nu_W(\lambda)<\infty$,
\[
T\,\frac{\nu_W(\lambda)}{e_\lambda}\le F_j(\ell)+C_c\Bigl(\frac{\log k\cdot m\,w}{\ell}+\frac{fk(S_c+\eta)}{\ell}+k\,\mathrm{th}_k(\ell)\,w
+[X''<\ell]\frac{kw}{\ell^2}\Bigr),
\]
and $\nu_W(\lambda)=0$ if $\deg(a)<\ell$. Here $\nu_W(\lambda)$ is the local depth of $X$ at $\lambda$ for $(M,a,n_W)$ and
$\mathrm{th}_k$ the level tail of the numerical datum.
\formallabel{GlobalTheory.T.goodPlaceDepthBound_proof}
\formallabel{GlobalTheory.GoodDepth.nuGood}
\end{theorem}

\begin{proof}
\emph{Content bounds for packages.} For packages with parameters satisfying the parameter conditions ($|\mathrm{Pl}|=f$, $m\ge1$,
$m+3\le k$, $r=k-1+m$, $S_E\ge0$) and $\ell\ge5$ prime, there are constants $C_I,C_P,C_S,C_C>0$ (depending on $C_1$) with
\begin{gather*}
T\nu\le\sum_{z\in Z}F^{\mathrm{num}}_\ell\bigl((k-1)w\ \text{along the orbit of}\ z\bigr)+C_I\,\mathrm{err}(T)\quad\text{(inert)},\\
T\nu\le\frac{(k-1)\sum_zw_z}{\ell^2-1}+C_P\,\mathrm{err}(T)\quad(\Phi),\qquad
T\nu\le C_S\,\frac{m\sum_zw_z+fkS_E}{\ell}\quad\text{(split)},\\ \nu\le0\quad\text{(zero packages)},
\end{gather*}
and the inert bound with $C_C$ for packages of conjugate type, where $F^{\mathrm{num}}_\ell(u)$ is the explicit orbit content of a
function $u$ on $\Z/d$ (\lref{Formalism.IncomingBounds.Fnum}) and $\mathrm{err}(T)=\frac{\log k\cdot m\sum w}{\ell}+\frac{fkS_E}{\ell}
+k\,\mathrm{th}_k(\ell)\sum w+Tf(k-1)^3(\log k)^2$ (\lref{GlobalTheory.T.inertPrimeDepthContentBound_proof},
\lref{GlobalTheory.T.complexConjFrobeniusDepthContentBound_proof}, \lref{GlobalTheory.T.splitPrimeDepthContentBound_proof},
\lref{GlobalTheory.T.largePrime_packages_depth_nonpos}, \lref{UpperTheory.rowBound_le_Fnum_sum_add_error}; these bounds are
obtained by inserting explicit block decompositions into the valuation-block property). Put $C_c=2C_I+2C_P+2C_C+C_S+2$.

\emph{Content values.} For $\ell$ as in the statement (\lref{GlobalTheory.ContentValue.main}): $F_j(\ell)\ge0$; $\ell$ is inert,
signal or split; $P(N,j,\ell)$ satisfies the parameter conditions; if $\ell$ is of type $\Phi$ then
$\frac{(k-1)\sum w}{\ell^2-1}\le F_j(\ell)+2A_4$; and if $\ell$ is inert or signal but not of type $\Phi$, then for every set $Z$ of
representatives $\sum_{z\in Z}F^{\mathrm{num}}_\ell(\dots)\le F_j(\ell)+2A_4$, where
\[
A_1=\frac{\log k\cdot m\,w}{\ell},\quad A_2=\frac{fk(S_c+\eta)}{\ell},\quad A_3=k\,\mathrm{th}_k(\ell)\,w,\quad A_4=[X''<\ell]\frac{kw}{\ell^2}.
\]

\emph{Row budgets.} From $10r(f+f^2(k-1)^3L^2/m)\le n_W$, $3\le L$ and $\log k\le L$ one obtains by elementary estimates
(\lref{GlobalTheory.RowBudgets.main}, \lref{UpperTheory.f_mul_leastAdmissibleLevel_rows_le}, the latter using
$i_0(k,\ell)\le3+2L$): $8rf\le n_W$; $fi_0(k,\ell)(k-1)+2rf\le n_W$ and $f((i_0(k,\ell)+1)(k-1)+4r)\le n_W$ for $\ell\ge5$; if
$\ell\ge mb/(Tf(k-1)^3L)$ and $b\ge0$, there is $s_0$ with $4rs_0f\le n_W$ such that $Tx\le3b$ implies $x\le s_0(\ell-1)$; and if
$\ell<mb/(Tf(k-1)^3L)$ and $b\le w$, then $Tf(k-1)^3(\log k)^2\le A_1$.

\emph{Bound for chart depths.} Let $B=F_j(\ell)+C_c(A_1+A_2+A_3+A_4)\ge0$. We show that every $\nu$ occurring in a good situation at
$\lambda$ satisfies $T\nu\le B$. If $b<0$, no $\tau$ has positive multiplicity (as $0\le Ta(\tau,0)\le3b$ would fail), and a zero
package with $s_0=0$ gives $\nu\le0$. If $b\ge0$ and $\ell\ge mb/(Tf(k-1)^3L)$, take $s_0$ from the row budgets; then
$a(\tau,0)\le s_0(\ell-1)$ by the first-row condition, and a zero package (of conjugate type if $\ell$ is of conjugate type) gives
$\nu\le0$. Otherwise $\ell<mb/(Tf(k-1)^3L)$, and the row budgets give $\mathrm{err}(T)\le2(A_1+A_2+A_3)$ (with $\sum w\le w$). If $\ell$ is of
type $\Phi$, the package of type $\Phi$ gives $T\nu\le F_j(\ell)+2A_4+2C_P(A_1+A_2+A_3)\le B$; if $\ell$ is of conjugate type, the
package of conjugate type and the content values give $T\nu\le F_j(\ell)+2A_4+2C_C(A_1+A_2+A_3)\le B$; if $\ell$ is otherwise inert or
signal, the same with the package of inert type and $C_I$; and if $\ell$ is split, the split package gives
$T\nu\le C_S(A_1+A_2)\le B$, using $m\sum w\le\log k\cdot m\,w$ ($\log k\ge1$ as $k\ge5$). In each case the required package exists by
the package statement: its degree hypothesis is the good-page condition (for $P(N,j,\ell)$ it follows from that for $P(N,j,1)$) and
its size bound is one of the row budgets.

\emph{Conclusion.} If every chart depth $\nu$ at $\lambda$ satisfies $\nu\le B/T$ with $B/T\ge0$, then the local depth, a supremum over
points of an infimum over charts, is finite with $\nu_W(\lambda)/e_\lambda\le B/T$, and it vanishes if $B=0$
(\lref{GlobalTheory.GoodDepth.nuW_bound_of_charts}; integral charts exist at all points). This gives the bound. If
$\deg(a)<\ell$, then $a(\tau,0)\le\deg(a)\le\ell-1$ for $\tau$ of positive multiplicity, so the zero packages with $s_0=1$ (whose size
bound $4rf\le n_W$ is a row budget) give $\nu\le0$ for every chart depth, hence $\nu_W(\lambda)=0$.
\end{proof}

\begin{theorem}[The good-prime leaf from depth bounds]\label{pf:onto2:leafgood}
For every $C_1\ge0$ there is $C_0$ with the following property. Let $(E,H,M)$, data $D_w(j)$, moduli spaces $X_j$ and a coefficient
datum $c$ be given, together with depth data, a number field $K$ over $c.F$, an indefinite datum $S$ over $K$, a record object $O$, an
index map $\iota$, numbers $N_S$ for $S$, a numerical datum $N$ and a place map $\mathrm{pl}$, such that: $C_0\le C$; $E(v,j)$,
$\sigma(v,j)$, $T$ and the primes $g_1,g_2,q_1,d,s_m$ of $O$ are those of the depth data at $\iota(j)$; the depth data split the
level primes and have level $g_1g_2$; the involution of $\OO_E$ is the host conjugation; integral charts exist and the package
statement holds for every $X_{\iota(j)}$; $\mathrm{pl}$ satisfies the place tie; for every $j$, $1\le\iota(j)\le m$, the signature is
$\sigma_W(\iota j)$, the weight $a_L(\iota j)$ is dominant and the good-page conditions hold with $(b,L)=(O.b,O.L)$; $O$ is tied to $N$
numerically ($k,m,f$ agree, $S_c+\eta\le S_e$, $X_2=X''$, $\theta=\mathrm{th}_k$, $\alpha=0$, $\sum_zw_{\iota j}(z)\le\mathrm{wt}(j)$,
$b\le\mathrm{wt}(j)$, $\deg a_L(\iota j)\le D_F$); the page leaf \ref{pf:onto2:P1} holds; the content tie $F_0(w,j)=F_{\iota j}(p_w)$ holds;
and the primes of $O$ are those of $N$. Then $O$ satisfies \ref{pf:onto2:L6} for $N_S$.
\formallabel{GlobalTheory.T.depthTable_goodPrimeLeaf}
\formallabel{GlobalTheory.GoodDepth.leafGood}
\end{theorem}

\begin{proof}
Let $C_c$ be given by \cref{pf:onto2:nugood} for $C_1$ and put $C_0=\max(C_c,20)$. By the numerical tie, $\theta(\ell)=\mathrm{th}_k(\ell)\ge0$
and $\alpha(\ell)=0$; hence $T\sum_{p\in Q}\alpha(p)\log p=0\le Cmb$, and
$\sum_{p\in Q}\theta(p)\log p\le20/k\le C/k$ for every finite set $Q$ of primes
(\lref{GlobalTheory.T.sum_level_tail_mul_log_primes_le}: $\sum_{p\in Q}\mathrm{th}_k(p)\log p\le20/k$).

Let $w\in\ontoIITw(S)$ with $\ell=p_w\notin s_m\cup\{q_1,d,g_1,g_2\}$ and $j$ an index. Then $\ell$ does not divide the level $g_1g_2$. For
$v\in\{w,\bar w\}$ the place $\lambda_v$ of $c.F$ below $v$ has residue characteristic $\ell$, and at such good places
(\lref{GlobalTheory.T.onObj_exponent_eq_rounded_depth})
\[
e_0(v,j)\le T\,\frac{\nu_W(\lambda_v)}{e_{\lambda_v}}+1,\qquad \nu_W(\lambda_v)=0\Rightarrow e_0(v,j)=0,
\]
with $\nu_W$ the local depth of $X_{\iota j}$ for $(M,a_L(\iota j),n_W(\iota j))$; this is the formula of \cref{pf:onto2:depdat} at a
good prime, where $c_P=0$ and the prime is not bad. \Cref{pf:onto2:nugood}, applied with $w=\mathrm{wt}(j)$ and $(b,L)=(O.b,O.L)$, gives
$T\nu_W(\lambda_v)/e_{\lambda_v}\le F_{\iota j}(\ell)+C_c(\cdots)$, and the error term is at most $a_E(j,\ell)$, since $C_c\le C$, $S_c+\eta\le S_e$,
$k,m\ge1$, $f\ge8$ and all terms are non-negative (\lref{GlobalTheory.GoodDepth.err_mono_leaf}). With the content tie,
$e_0(v,j)\le F_0(w,j)+a_E(j,\ell)+1$. If $Y<\ell$, then $\deg a_L(\iota j)\le D_F\le CD_F=Y<\ell$ (as $C\ge20$), so
$\nu_W(\lambda_v)=0$ by \cref{pf:onto2:nugood} and $e_0(w,j)=e_0(\bar w,j)=0$; finally $F_0(w,j)=F_{\iota j}(\ell)\ge0$ by the content values
(\lref{GlobalTheory.T.forcedContent_bounds_programmeValue_goodPrime}). This is \ref{pf:onto2:L6}.
\end{proof}

\begin{proposition}[The good-prime leaf for record objects]\label{pf:onto2:leafgoodrun}
There is $C_0$ such that, for a core with its laws and all data as before \cref{pf:onto2:dlevel}, if the depth tie holds, the package
statement holds for the moduli spaces of the core at all indices $\iota(j)$, $C_0\le C$, \ref{pf:onto2:P1}, \ref{pf:onto2:P2},
$10\le C_{nW}$, $3\le x^{e_z/5}$, $3\le x$, $kb'\le b$ and the content tie with $\mathrm{num}(D,L)$ hold, then $O$ satisfies
\ref{pf:onto2:L6}.
\formallabel{GlobalTheory.GoodDepth.leafGood_run}
\formallabel{GlobalTheory.T.depthLeaf_goodPrimes_of_point_data}
\end{proposition}

\begin{proof}
Apply \cref{pf:onto2:leafgood} with $C_1=100$, the depth datum of the core, $N=\mathrm{num}(D,L)$ and $\mathrm{pl}$ the composite of the
place map of $D$ with the embeddings of its types, and check its hypotheses. By the depth tie, $O.x=x$ and $O.\varepsilon=\varepsilon$,
so $O.b=b(L)$ and $O.L=L$; the depth tie also gives the identities for $E$, $\sigma$, $T$ and the primes, the splitting of the level
primes and the level, and the star structure. Integral charts exist at all integral framed points
(\lref{GlobalTheory.T.exists_integral_chart_integral_framed_point}), and the real places match the CM type through Frobenius
(\lref{GlobalTheory.T.realPlaces_match_cmTypeFrobenius}), which is the place tie. At each $\iota(j)$, admissibility and the
signature come from the pins of the laws, dominance from that of the run weight, and the good-page conditions with $C_1=100$ from the
depth parameters (\lref{GlobalTheory.T.goodPrimeWeightHypotheses_of_depthParameters}, using $10\le C_{nW}$,
$3\le x^{e_z/5}$, $3\le x$, $kb'\le b$ and $\log k\le L$ from \ref{pf:onto2:P2}). For the numerical tie: $k,m,f,X_2,\theta,\alpha$ agree by
the depth tie, $S_c+\eta\le S_e$ (\lref{GlobalTheory.GoodDepth.numOf_sizes_le}), $w_{\iota j}(z)\le2b$ for all $z$
(\lref{GlobalTheory.T.firstRows_bound_and_sizes_bounds}) so that $\sum_zw_{\iota j}(z)\le2fb=\mathrm{wt}(j)$, $b\le\mathrm{wt}(j)$ as
$f\ge8$, and $\deg a_L(\iota j)\le D_F$ by the degree identity for weight exponents
(\lref{GlobalTheory.WeightExponents.weight_exponents}) and the depth tie.
\end{proof}

\begin{proposition}[Depth at good primes]\label{pf:onto2:dgood}
Assume the formal-group statement $(\mathrm U_p)$ for all primes $p$. There is $C_0$ such that, under these assumptions and
$C_0\le C$, $10\le C_{nW}$ and $3\le x^{e_z/5}$, $O$ satisfies \ref{pf:onto2:L6} for $N$.
\formallabel{GlobalTheory.leafDepthGood_of_chartForm}
\formallabel{GlobalTheory.T.depthTable_goodPrimes_of_chartForm}
\end{proposition}

\begin{proof}
By \cref{pf:onto2:packages}, $(\mathrm U_p)$ for all $p$ implies the package statement for every moduli space, in particular for the
moduli spaces of the core at the indices $\iota(j)$. Take $C_0$ from \cref{pf:onto2:leafgoodrun}; its remaining hypotheses are among
the assumptions fixed before \cref{pf:onto2:dlevel} and those of the statement.
\end{proof}

\begin{corollary}[Depth leaves for a core]\label{pf:onto2:depthrun}
Assume $(\mathrm A_p)$, $(\mathrm H_p)$ and $(\mathrm U_p)$ for all primes $p$. There is $C_0$ such that, under the assumptions fixed
before \cref{pf:onto2:dlevel} and $C_0\le C$, $10\le C_{nW}$ and $3\le x^{e_z/5}$, $O$ satisfies \ref{pf:onto2:L6}--\ref{pf:onto2:L9}
for $N$.
\formallabel{GlobalTheory.RunDepth.leafDepth_all_run}
\end{corollary}

\begin{proof}
Take $C_0$ from \cref{pf:onto2:dgood}; then \cref{pf:onto2:dgood,pf:onto2:dram,pf:onto2:d2,pf:onto2:dlevel} give
\ref{pf:onto2:L6}, \ref{pf:onto2:L7}, \ref{pf:onto2:L8} and \ref{pf:onto2:L9}.
\end{proof}

\begin{proposition}[The depth leaves]\label{pf:onto2:depthall}
There is $C_0$ such that for all data as at the beginning of \cref{sec:pf:onto2:geo}, all numbers $N$ for the false side, and every
record object $O$: if the record relation, \ref{pf:onto2:P1}, \ref{pf:onto2:P2}, $3\le x$, $kb'\le b$, the content tie of $N$ with
$\mathrm{num}(D,L)$ for $\iota(j)=j+1$, $C_0\le C$, $10\le C_{nW}$ and $3\le x^{e_z/5}$ hold, then $O$ satisfies
\ref{pf:onto2:L6}--\ref{pf:onto2:L9} for $N$.
\formallabel{GlobalTheory.BigOn.leafDepth_all_record}
\end{proposition}

\begin{proof}
By \cref{pf:onto2:depthrun}, whose hypotheses $(\mathrm A_p)$, $(\mathrm H_p)$ and $(\mathrm U_p)$ hold for every prime $p$
(\lref{GlobalTheory.ChF.colength_bound_split_place_chart}, \lref{GlobalTheory.ChF.hodgeTate_inclusion_chart},
\lref{GlobalTheory.ChF.unramified_companion_chart_characters}), it suffices to provide the data fixed before
\cref{pf:onto2:dlevel}. For a
record object as in the statement, component (iv) of the record relation provides a star structure on $\OO_E$ and the depth tie for
the core $J$ of the runs and $\iota(j)=j+1$, and component (iii) provides the laws of $J$; we take $K=\mathrm{LL}$ with the algebra
structure of the host.
\end{proof}

\subsubsection{Proof of \texorpdfstring{\cref{pf:onto2:thmA}}{Theorem A}}\label{sec:pf:onto2:geoproof}

Let $C_K,x_a$ be the constants of \cref{pf:onto2:runshold}, $x_b$ the constant given by \cref{pf:onto2:latleaves} for $C_K$, and $C_d$
the constant of \cref{pf:onto2:depthall}. By \lref{GlobalTheory.T.runNumericalConditions_hold_along_family}, for every $x_1$ the
following statement $\mathrm{Num}(x_1)$ holds closedly over $\sfname{FamHyp}$: there is $y$ such that for all $x>x_1(\mathcal F)$ with
$x\ge y$ and all $i$, the parameters of the family at $(x,i)$ satisfy $\ontoIIPN(x_1)$, $3\le x$, $kb'\le b$, $3\le x^{e_z/5}$ and
$8fr+4\le n_V$. Closed statements are closed under conjunction and implication (\lref{Formalism.MainNotions.Closed.and},
\lref{Formalism.MainNotions.Closed.imp}), so it suffices to show, for a fixed family $\mathcal F$ of $\sfname{FamHyp}$ and $C_H$,
that the conjunction of the inner statements of the page statement, $\mathrm{Num}(x_a)$ and $\mathrm{Num}(x_b)$ implies the inner
statement of the geometric statement.

Let $c_0$ be the constant of the page statement and put $c_0'=\max(c_0,C_d)$. Let $c$ be constants with $c_0'\le C$, $14\le C_{b1}$
and $66C+4\le C_{b2}$, let $x_0$ be the threshold of the page statement for $c$, and $y_a,y_b$ those of $\mathrm{Num}(x_a)$,
$\mathrm{Num}(x_b)$. Let $x>x_1(\mathcal F)$ with $x\ge\max(x_0,y_a,y_b)$, $i$ an index, $p$ a point over $(x,i)$ satisfying the point
conditions and $\mathcal M(p)$, $r$ a pair of runs and $O$ a record object of $p$ along $r$. The page statement gives \ref{pf:onto2:P1},
\ref{pf:onto2:P2} and $O.b=b(L)$; $\mathrm{Num}(x_a)$ gives $\ontoIIPN(x_a)$, $3\le x$, $kb'\le b$ and $3\le x^{e_z/5}$;
$\mathrm{Num}(x_b)$ gives $\ontoIIPN(x_b)$. From $\mathcal M(p)$ we get the core ties of the core of $p$ with $R_2(p)$, and $O$ satisfies
the record relation and the own-prime pin. By \lref{GlobalTheory.T.links2_of_single_index_twisted} (the link conditions of the
two-sided reduction hold at every point satisfying the point conditions at which every single-index place is twisted, the first
component of $\mathcal M(p)$), $R_2(p)$ satisfies the group condition and the three tuple conditions.

Now \cref{pf:onto2:runshold} gives the run bounds $\ontoIIRB(C_K,r)$ (with the weight definition of the auxiliary record of $\mathcal F$
at $(x,i)$); \cref{pf:onto2:latleaves} gives \ref{pf:onto2:L1}--\ref{pf:onto2:L4}; \cref{pf:onto2:index} gives \ref{pf:onto2:L5}, using
\ref{pf:onto2:L1}, \ref{pf:onto2:L2} and $O.b=b(L)$. Finally, for $C,C_H'\in\R$ the numbers $N_p(C,C_H')$ satisfy the content
tie with $\mathrm{num}(D,L)$ for $\iota(j)=j+1$: both contents are, by definition, the content function of the numerical datum of the
record of $p$ (\lref{GlobalTheory.BigOn.indefinitePoint_splitPartContent_eq_forcedContent}, which rewrites the content of
$\mathrm{num}(D,L)$ with \lref{GlobalTheory.BigOn.numOf_Fj_eq_ofRecord} using the fields $\mathrm{Del}$, $h_0$ of the point).
Since $C_d\le C$ and $C_{nW}=10$ by $\ontoIIPN(x_a)$, \cref{pf:onto2:depthall} gives \ref{pf:onto2:L6}--\ref{pf:onto2:L9} for
$N_p(C,C_H')$. \qed

%======================================================================
\subsection{The exponent statement}\label{sec:pf:onto2:exp}

\subsubsection{Weight dictionaries}\label{sec:pf:onto2:dict}

\begin{definition}[Weyl exponents]\label{pf:onto2:weylexp}
Let $N\ge1$ and let $\nu=(\nu_0\ge\dots\ge\nu_{N-1})$ be a dominant weight. In the polynomial ring $\Z[x_{ab}]_{a,b<N}$ let
$\nabla_\Z(\nu)$ be the module of polynomials of column weight $\nu$ which are invariant under the column operations
$x_{\cdot j}\mapsto x_{\cdot j}+t\,x_{\cdot l}$ ($l<j$), and $\Delta_\Z(\nu)\subseteq\nabla_\Z(\nu)$ the span of the iterated divided row
operations applied to $\prod_{s<N}\det_{s+1}^{\nu_s-\nu_{s+1}}$, where $\det_s$ is the leading principal $s\times s$ minor
(\lref{Formalism.GradedWeylDivisibility.dualWeyl}, \lref{Formalism.GradedWeylDivisibility.weyl}; these are integral models
of the induced and Weyl modules of $\GL_N$ \cite{Jantzen2003}). For a prime $\ell$, the \emph{Weyl exponent} $\ontoIIexp_\ell(\nu)$ is
the least $e\in\N$ such that $m\ell^e\nabla_\Z(\nu)\subseteq\Delta_\Z(\nu)$ for some $m$ prime to $\ell$, and $0$ if there is no such
$e$ (\lref{Formalism.CohomologicalDescent.LA.exZ}). For $\lambda:\N\to\N$ we write $\lambda|_N$ for $(\lambda(0),\dots,\lambda(N-1))$
and $U_N(\lambda,a)=\lambda(a)+N-1-a$.
\end{definition}

The \emph{cut inequality} (\lref{Formalism.CohomologicalDescent.E.ExCut}) is the statement: for all $N$, all $\lambda$ with
$\lambda|_N$ antitone, every prime $\ell$, all $\mathrm{lo}\in\Z$ and $w\in\N$ with $\mathrm{lo}\le U_N(\lambda,a)\le\mathrm{lo}+w$ for
$a<N$, and every $a\le N$,
\begin{equation}\label{eq:pf:onto2:excut}
\ontoIIexp_\ell(\lambda|_N)\le\ontoIIexp_\ell(\lambda|_a)+\ontoIIexp_\ell\bigl(\lambda(a+\cdot)|_{N-a}\bigr)+a(N-a)\lfloor\log_\ell w\rfloor .
\end{equation}
It is proved in \lref{GlobalTheory.LeviCut.exCut}, and implies (\lref{Formalism.CohomologicalDescent.E.exIntervalLt_of_exCut})
the interval bound $\ontoIIexp_\ell(\lambda|_N)\le\binom N2\lfloor\log_\ell w\rfloor$, with $\ontoIIexp_\ell(\lambda|_N)=0$ if $w<\ell$.

\begin{definition}[Parameters and weights]\label{pf:onto2:parwts}
A \emph{parameter} $P$ (\lref{Formalism.CohomologicalDescent.LogSmallConstantBounds.Par}) consists of natural numbers
$k,m,t,f,n,N_r$ and reals $L,x,b',\eta,S$; put $b=e^L$, $h=b+S+\eta+t+1$, $l^*=mb/\bigl(nt^2(k+m(t+1)+t)L^2+kn^2(S+t+L)\bigr)$,
$l_J=mbL/\bigl(k(m+t)\log(3h)\bigr)$, $k_{\max}=\lfloor2b/l^*\rfloor$, and let $\mathrm{Excl}(\ell)$ mean that $\ell$ divides an integer
in $[b-\tfrac{t-1}2-k_{\max},\,h+\tfrac{t-1}2+k_{\max}]$. The conditions $\mathrm{Hyp}(P)$ and $\mathrm{Large}(P)$ are explicit lists of
inequalities between these numbers (for instance $N_r+1=k+m+t$, $t\ge k+m+4$, $3m<t$, $1\le m\le k$, $L\ge4$, $48k^2t^2<b$), and
$\mathrm{Rows}(P,X',A)$, $\mathrm{Big}(L,x,X',A)$ are further explicit inequalities.

\emph{Weights} $W$ on a finite set $\iota$ (\lref{Formalism.CohomologicalDescent.ShortWeightBounds.Wts}) consist of functions
$\lambda_\tau:\N\to\N$ ($\tau\in\iota$), a subset $T_2\subseteq\iota$, integers $j_\tau,q_\tau,x^\pm_\tau$, reals $c_\tau,e_\tau$, maps
$\mathrm{fr}:\N\times\iota\to\iota$, $\mathrm{cj}:\N\times\iota\to\{0,1\}$ and a finite set $S_p$ of primes. Put
$U_\tau(a)=U_{N_r}(\lambda_\tau,a)$ and $h_0=\lfloor h\rfloor$. The condition $\mathrm{Ok}_P(W)$ requires $\lambda_\tau|_{N_r}$ antitone, a
constant block of length $t$ starting at $k-1+j_\tau$ with $k-1+j_\tau+t+q_\tau=N_r$, and explicit position conditions on the values
$U_\tau(a)$ relative to $c_\tau,e_\tau,x^\pm_\tau$ (\lref{Formalism.CohomologicalDescent.ShortWeightBounds.Wts.Ok}). The
\emph{exponents} are
\[
\ontoIIexp(\tau,\ell)=\ontoIIexp_\ell(\lambda_\tau|_{N_r}),\qquad \ontoIIJA(\ell)=\sum_\tau\ontoIIexp(\tau,\ell);
\]
the prime $\ell$ is \emph{in range} if $l_J<\ell\le h$, $\neg\mathrm{Excl}(\ell)$ and $\ell\notin S_p$; the \emph{coarse bound}
$J^{\mathrm{c}}(\ell)$ is an explicit integer built from the exponents $\ontoIIexp(\tau,\ell)$, the Weyl exponents of the first $k-1$
entries of $\lambda_\tau$ and counts of the values $U_\tau(a)$ in residue classes modulo $\ell$; and the \emph{recursive exponent} is
\[
\ontoIIJr(\ell)=\begin{cases}\ontoIIJA(\ell)&\text{if }\ell\le l_J,\ \mathrm{Excl}(\ell)\text{ or }\ell\in S_p,\\
J^{\mathrm{c}}(\ell)&\text{otherwise}\end{cases}
\]
The pair $(P,W)$ is \emph{admissible} if $\mathrm{Hyp}(P)$,
$\mathrm{Large}(P)$, $\mathrm{Ok}_P(W)$ hold and $\mathrm{Rows}(P,X',A)$, $\mathrm{Big}(P.L,P.x,X',A)$ hold for some $X',A$.
\formallabel{Formalism.CohomologicalDescent.LargeParameterEstimates.Rows}
\formallabel{Formalism.CohomologicalDescent.LargeParameterEstimates.Big}
\formallabel{Formalism.CohomologicalDescent.ShortWeightBounds.Wts.coarseWeightBound}
\formallabel{Formalism.CohomologicalDescent.ShortWeightBounds.Wts.Jrec7}
\formallabel{GlobalTheory.ExpTie.NumOk}
\end{definition}

\begin{definition}[Blocks and lattice inclusions]\label{pf:onto2:blocks}
A \emph{block datum} $B$ over a number field $K$ with involution and a finite group $G$ (\lref{Formalism.L2.DIV.BlockDat}) consists
of a space of class-sum forms with a representation of $G$, Hecke operators $B.\mathrm{op}(P,i)$ indexed by places $P$ and integers
$i$, and block data; it carries two $\OO_K$-lattices of forms, $\nabla(B)$ and $\Delta(B)$. For a level $\ell_v$ (a finite set of
functions on $G$ with a subgroup and a character) and a $K$-subspace $Z$ of forms, the \emph{cut lattices} $\mathrm{cut}_Z\nabla$ and
$\mathrm{cut}_Z\Delta$ are the images of $\nabla(B)$, $\Delta(B)$ in coordinates for the good part of $Z$. For a place $w$ of $K$ with residue
characteristic $\ell=p_w$ and $J\in\N$ we write
\[
\ontoIILI_w(J):\iff \ell^J\cdot\mathrm{cut}_Z\nabla\subseteq\ontoIIloc{w}{(\mathrm{cut}_Z\Delta)},
\]
and the \emph{normalised exponent} is
\[
\ontoIIJu(B,\ell_v,Z,w)=\frac{J_W}{e(w|\ell)},\qquad J_W=\min\{J:\ a\cdot\mathrm{cut}_Z\nabla\subseteq\ontoIIloc
w{(\mathrm{cut}_Z\Delta)}\ \text{for all }a\in w^J\},
\]
with $\min\emptyset=0$. Clearly $\ontoIILI_w(J)\Rightarrow\ontoIILI_w(J+1)$. For a point $p$ of a family, an index $j$ and a type
$t$, the point carries a block datum $\mathrm{blk}(p,j,t)$ (\lref{GlobalTheory.ExpTie.blk}), a level $\ell_v(p)$ and a subspace
$Z(p,j,t)$, those of its two-sided reduction.
\formallabel{Formalism.L2.DIV.BlockDat.cutZ}
\formallabel{Formalism.L2.DIV.BlockDat.delZ}
\formallabel{GlobalTheory.ExpTie.LIncl}
\formallabel{Formalism.L2.DIV.BlockDat.JexpU}
\formallabel{GlobalTheory.ExpTie.lvlL}
\formallabel{GlobalTheory.ExpTie.subZ}
\end{definition}

By definition of the numbers of the two-sided reduction, their exponents are these normalised exponents:
\begin{equation}\label{eq:pf:onto2:J0}
J_0\bigl(N_p(C,C_H)\bigr)(t,w,j)=\ontoIIJu\bigl(\mathrm{blk}(p,j,t),\ell_v(p),Z(p,j,t),w\bigr)\qquad(C,C_H\in\R).
\end{equation}

\begin{definition}[Dictionaries]\label{pf:onto2:dictdef}
For a datum $D$ with size data $L$, an index $j<m$, a parameter $P$ and weights $W$ on the type set $\mathrm{Tau}(D)$, $(P,W)$ is a
\emph{dictionary at $j$} if: $(k,m,t,f)$ of $P$ are the ranks of $D$, $|\mathrm{Tau}|=f$,
$N_r=\rank V_j$, $x=L.x$, $L=L.x^{L.\varepsilon}$ and $b=L.b$; each $\lambda_\tau$ is the shift by a constant of the forced weight
$\lambda^{\mathrm{def}}_{j,\tau}$ of $D$ (\lref{Formalism.IndefR.Bridge.lamLdef}), which is a genuine weight of $D$; $0\le
U_\tau(a)\le h_0$ for $a<N_r$; $|T_2|\le|T_2(D)|$; $S_p=s_m\cup\{q,d,p_{\mathrm{de}},p'_{\mathrm{de}}\}$; and $h<g_i$ for both level
primes. The four statements used below are of the following form. We say that we are \emph{in range at $(p,j,P,W,w)$} if $p$ is a
point satisfying the extra point conditions, $(P,W)$ is an admissible dictionary at $j$, $w$ is a twisted place of the reduction of $p$
and $p_w$ is in range for $(P,W)$; a property $Q$ \emph{holds in range} if, closedly over $\sfname{FamHyp}$ and eventually at points,
$Q(p,j,P,W,w)$ holds whenever we are in range at $(p,j,P,W,w)$.
\begin{itemize}
\item The \emph{numerical statement}: $\mathrm{AtPts}$ of ``for every $j$ there is an admissible dictionary at $j$''.
\item The \emph{exponent-sum statement}: $\mathrm{AtPts}$ of ``for every admissible dictionary $(P,W)$ at $j$, every type $t$ and every
  place $w$, $\ontoIILI_w(\ontoIIJA(p_w))$ holds for $(\mathrm{blk}(p,j,t),\ell_v(p),Z(p,j,t))$''.
\item The \emph{smallness}, \emph{frame}, \emph{Galois} and \emph{model statements}: the conditions of
  \cref{pf:onto2:small,pf:onto2:frame,pf:onto2:galois,pf:onto2:model} hold in range.
\end{itemize}
\formallabel{GlobalTheory.ExpTie.Dict}
\formallabel{GlobalTheory.ExpTie.InRangeAlt}
\end{definition}

%----------------------------------------------------------------------
\subsubsection{Reduction to four statements}\label{sec:pf:onto2:red}

\begin{theorem}[The exponent statement]\label{pf:onto2:thmB}
Assume the following published results, in the form in which they are fields of the package of Part~I (\cref{def:hyps}): (H1)
Clozel's purity theorem for regular algebraic cuspidal representations \cite{Clozel1990}; (H2) the $\ell$-adic characters attached to
algebraic Hecke characters \cite{Serre1968}; (H3) Chebotarev's density theorem with degree-one primes \cite{Chebotarev1926}; (H4)
Galois representations for polarised regular algebraic cuspidal representations of $\GL_a$ over CM fields
\cite{ClozelHarrisLabesse2011,Shin2011,ChenevierHarris2013,Caraiani2012,BarnetLambEtAl2014b}; (H5) base change from definite unitary groups
\cite{Labesse2011}; (H6) the Bhatt--Scholze theorem on crystalline lattices \cite{BhattScholze2023}; (H7) the crystallinity of the
$\ell$-adic characters above $\ell$ with the expected Hodge--Tate weights; and (H8) the existence of local Langlands parameters
($\sfname{RecTExists}$, \cref{def:rectexists}). Then the exponent statement holds: closedly over $\sfname{FamHyp}$, for large constants
and eventually at points, every record object $O$ at a point $p$ satisfies \ref{pf:onto2:L10} for $N_p(C,C_H)$, for all $C,C_H$.
\formallabel{GlobalTheory.T.expSt_of_base_change_and_crystalline_lattices}
\formallabel{GlobalTheory.BigOn.exponent_bound_of_unitaryThetaStandardTheorems_alt}
\end{theorem}

See \cref{sec:literature:aut,sec:literature:padic} for the published results formalised by (H1)--(H8).

\begin{proof}
By \cref{pf:onto2:small} the smallness statement holds unconditionally, so by \cref{pf:onto2:frame} the frame statement holds and by
\cref{pf:onto2:galois} (which assumes (H1)--(H8); (H1) is not used, \cref{pf:onto2:rem:unused}) the Galois statement holds. By \cref{pf:onto2:model} (which uses (H3) and (H6)) the
smallness, frame and Galois statements imply the model statement. The cut inequality holds
(\lref{GlobalTheory.LeviCut.exCut}), the numerical
statement holds by \cref{pf:onto2:numst} and the exponent-sum statement by \cref{pf:onto2:sumexp}. \Cref{pf:onto2:expst} now gives the
exponent statement.
\end{proof}

\begin{lemma}[Exponents from inclusions]\label{pf:onto2:jexpU}
If $\ontoIILI_w(J)$ holds for $(B,\ell_v,Z)$, then $0\le\ontoIIJu(B,\ell_v,Z,w)\le J$.
\formallabel{GlobalTheory.ExpTie.LIncl.jexpU_bounds}
\formallabel{GlobalTheory.T.exponent_datum_le_of_primePowerLatticeInclusion}
\end{lemma}

\begin{proof}
Non-negativity is clear. Let $\ell=p_w$ and $e=e(w|\ell)$. Let $a\in w^{Je}$ and $z\in\mathrm{cut}_Z\nabla$. There are $t\in\OO_K\setminus w$
and $b\in\OO_K$ with $ta=\ell^Jb$ (\lref{GlobalTheory.ExpTie.exists_mul_eq_prW_pow}: $a$ has $w$-adic valuation at least that of
$\ell^J$). By hypothesis there is $s\in\OO_K\setminus w$ with $s\ell^Jz\in\mathrm{cut}_Z\Delta$; then $(st)(az)=b(s\ell^Jz)\in\mathrm{cut}_Z\Delta$
and $st\notin w$, so $az$ lies in the localisation of $\mathrm{cut}_Z\Delta$ at $w$. Hence $J_W\le Je$ and $\ontoIIJu=J_W/e\le J$.
\end{proof}

\begin{lemma}[Recursive exponents versus formulas]\label{pf:onto2:recform}
Assume the cut inequality, $\mathrm{Hyp}(P)$, $\mathrm{Large}(P)$ and $\mathrm{Ok}_P(W)$, and let $\ell$ be prime. If $\ell$ is not in range
then $\ontoIIJA(\ell)\le\ontoIIJr(\ell)$; if $\ell$ is in range then $J^{\mathrm c}(\ell)=\ontoIIJr(\ell)$.
\formallabel{GlobalTheory.exponent_of_record_vs_formulas}
\end{lemma}

\begin{proof}
If $\ell$ is in range then none of $\ell\le l_J$, $\mathrm{Excl}(\ell)$, $\ell\in S_p$ holds, so $\ontoIIJr(\ell)=J^{\mathrm c}(\ell)$ by
definition. Suppose $\ell$ is not in range. If $\ell\le l_J$, $\mathrm{Excl}(\ell)$ or $\ell\in S_p$, then $\ontoIIJr(\ell)=\ontoIIJA(\ell)$.
Otherwise $l_J<\ell$, $\neg\mathrm{Excl}(\ell)$ and $\ell\notin S_p$, so not being in range means $h<\ell$. In that case
$\ontoIIexp(\tau,\ell)=0$ for every $\tau$ by \lref{Formalism.CohomologicalDescent.E.beyond_height_exceptional} (for primes
$\ell>\max(l_J,h)$ all the Weyl exponents $\ontoIIexp(\tau,\ell)$ vanish; this is proved from the cut inequality and its interval
consequence), so $\ontoIIJA(\ell)=0\le\ontoIIJr(\ell)$.
\end{proof}

\begin{lemma}[The exponent at a place]\label{pf:onto2:jrec}
Assume the cut inequality and that $(P,W)$ is admissible; let $B,\ell_v,Z$ be as in \cref{pf:onto2:blocks} and $w$ a place with
$\ell=p_w$. Assume: if $\ell$ is not in range then $\ontoIILI_w(\ontoIIJA(\ell))$; if $\ell$ is in range then the local-model condition
of \cref{pf:onto2:hida} holds at $(B,\ell_v,Z,w)$. Then $\ontoIIJu(B,\ell_v,Z,w)\le\ontoIIJr(\ell)$.
\formallabel{GlobalTheory.ExpTie.Asm.jexpU_le_jrec7}
\end{lemma}

\begin{proof}
If $\ell$ is in range, \cref{pf:onto2:lincl} gives $\ontoIILI_w(J^{\mathrm c}(\ell))$ from the local-model condition.
Hence $\ontoIIJu\le J^{\mathrm c}(\ell)=\ontoIIJr(\ell)$ by \cref{pf:onto2:jexpU,pf:onto2:recform}. If $\ell$ is not in range, the
hypothesis and \cref{pf:onto2:jexpU} give $\ontoIIJu\le\ontoIIJA(\ell)\le\ontoIIJr(\ell)$ by \cref{pf:onto2:recform}.
\end{proof}

\begin{lemma}[Summation with few exceptional primes]\label{pf:onto2:sumcut}
Assume the cut inequality, $\mathrm{Hyp}(P)$, $\mathrm{Large}(P)$, $\mathrm{Rows}(P,X',A)$ and $\mathrm{Big}(P.L,P.x,X',A)$, and let $W$ be
weights on a finite set $\iota$ with $\mathrm{Ok}_P(W)$, $|\iota|=P.f$ and $|S_p|\le5$. Then for every $T\ge|T_2|$ and every finite set
$Q$ of primes
\[
\sum_{\ell\in Q}\ontoIIJr(\ell)\log\ell\le30\,T\,m\,b\,L^2+2mb .
\]
\formallabel{GlobalTheory.sum_exponent_log_le_of_cut}
\end{lemma}

\begin{proof}
By \lref{Formalism.CohomologicalDescent.E.recursion_sum_exceptional}, the sum is at most an explicit quantity
$B_{\mathrm{tot}}(P,|T_2|,30\,mbL^2)$ (\lref{Formalism.CohomologicalDescent.LogSmallConstantBounds.totalBound}); that result verifies
that the recursive exponents satisfy the inputs of an elementary summation over primes, using the cut inequality, its interval
consequence (\lref{Formalism.CohomologicalDescent.E.exIntervalLt_of_exCut}) and the prime number bound $\pi(y)<1.26\,y/\log y$ for
$y>1$ (\cite[Cor.~1]{RosserSchoenfeld1962}; proved in the development as
\lref{Formalism.CohomologicalDescent.Top.prime_count_bound_rs}), and it needs $|\iota|=P.f$ and $|S_p|\le5$. By the tail inequality
\lref{Formalism.CohomologicalDescent.LargeParameterEstimates.tail_bound_holds}, which needs $|T_2|\le P.f$ (true since
$T_2\subseteq\iota$) together with $\mathrm{Hyp}$, $\mathrm{Rows}$ and $\mathrm{Big}$, this quantity is at most $30|T_2|mbL^2+2mb$. Since
$b>0$ under these hypotheses (\lref{Formalism.CohomologicalDescent.LargeParameterEstimates.large_parameter_facts_of_hyp}),
$mbL^2\ge0$, and $|T_2|\le T$ gives the claim.
\end{proof}

\begin{proposition}[Summation of recursive exponents]\label{pf:onto2:sumlog}
Assume the cut inequality, $\mathrm{Hyp}(P)$, $\mathrm{Large}(P)$, $\mathrm{Rows}(P,X',A)$, $\mathrm{Big}(P.L,P.x,X',A)$, $\mathrm{Ok}_P(W)$,
$|\iota|=P.f=D.R.f$ and $S_p=s_m\cup\{q,d,p_{\mathrm{de}},p'_{\mathrm{de}}\}$ for a datum $D$. Then for every $T\ge|T_2|$ and every finite
set $Q$ of primes
\[
\sum_{\ell\in Q}\ontoIIJr(\ell)\log\ell\le30\,T\,m\,b\,L^2+2mb .
\]
\formallabel{UpperTheory.ExcSet.sum_exponent_log_le}
\end{proposition}

\begin{proof}
Let $W'$ be $W$ with $S_p$ replaced by $S'=\{q,d,p_{\mathrm{de}},p'_{\mathrm{de}}\}$; $\mathrm{Ok}_P(W')$ still holds (the set $S_p$ does
not enter $\mathrm{Ok}_P$) and $|S'|\le5$. The recursive exponents of $W$ and $W'$ agree: the small primes of the host satisfy $\ell\le l_J$
(\lref{UpperTheory.ExcSet.small_le_lJ}), so for every $\ell$ the disjunction ``$\ell\le l_J$ or $\mathrm{Excl}(\ell)$ or $\ell\in S_p$''
is equivalent to the same disjunction with $S'$. Since $T_2$ is unchanged, \cref{pf:onto2:sumcut} applied to $W'$ gives the claim.
\end{proof}

\begin{lemma}[The exponent bound at a record object]\label{pf:onto2:hasexp}
Assume the cut inequality. Let $p$ be a point and $O$ a record object for its two-sided reduction with $m=D.R.m$, $x=L.x$,
$\varepsilon=L.\varepsilon$ and $T_2=|T_2(D)|$. Let $(P_j,W_j)_j$ be admissible dictionaries at the indices $j$, such that
$\ontoIILI_w(\ontoIIJA_j(p_w))$ holds for $(\mathrm{blk}(p,j,t),\ell_v(p),Z(p,j,t))$ for all $j,t,w$, and such that the local-model
condition holds for $(\mathrm{blk}(p,j,\mu(w)),\ell_v(p),Z(p,j,\mu(w)),w)$ at every twisted place $w$ with $p_w$ in range for $(P_j,W_j)$.
Then $O$ satisfies \ref{pf:onto2:L10} for $N_p(C,C_H)$, for all $C,C_H$, with $J_p(j,\ell)=\ontoIIJr_j(\ell)$.
\formallabel{GlobalTheory.ExpTie.Asm.hasExpBound_at}
\end{lemma}

\begin{proof}
\emph{Non-negativity.} By \eqref{eq:pf:onto2:J0} and \cref{pf:onto2:jexpU} applied to the inclusion with $J=\ontoIIJA_j(p_w)$,
$J_0(\mu(w),w,j)\ge0$. \emph{Summed bound.} $J_p(j,\ell)=\ontoIIJr_j(\ell)\ge0$. By the dictionary, $P_j.x=L.x$ and
$P_j.L=L.x^{L.\varepsilon}=x^\varepsilon=x_e$, hence $P_j.b=e^{x_e}=O.b$, and $P_j.m=m=O.m$. \Cref{pf:onto2:sumlog} (with $T=O.T_2\ge|T_2(W_j)|$,
by the dictionary) gives $\sum_{\ell\in Q}J_p(j,\ell)\log\ell\le30\,O.T_2\,O.m\,O.b\,x_e^2+2\,O.m\,O.b=E_J$. \emph{Domination.} For a twisted
place $w$, \eqref{eq:pf:onto2:J0} and \cref{pf:onto2:jrec} give $J_0(\mu(w),w,j)=\ontoIIJu\le\ontoIIJr_j(p_w)=J_p(j,p_w)$.
\end{proof}

\begin{proposition}[Exponent statement from four statements]\label{pf:onto2:expst}
The cut inequality, the numerical statement, the exponent-sum statement and the model statement imply the exponent statement.
\formallabel{GlobalTheory.ExpTie.Asm.expSt_of_alt}
\formallabel{GlobalTheory.T.expSt_of_num_sumExp_model}
\end{proposition}

\begin{proof}
The three statements of the form $\mathrm{AtPts}$ combine to one (conjunction of closed statements). Given a family, take $c_0=0$, and
for constants $c$ take a threshold $x\ge\max(x_1,x_2,x_3)$ for the three statements. Let $p$ be a point with $\mathcal M(p)$, $r$ runs and
$O$ a record object of $p$ along $r$. By the numerical statement choose admissible dictionaries $(P_j,W_j)$ at every $j$. The depth tie
of the record relation gives $O.m=D.R.m$, $O.x=L.x$, $O.\varepsilon=L.\varepsilon$, and the record relation gives $O.T_2=|T_2(D)|$. The
exponent-sum statement gives the inclusions with $\ontoIIJA_j$, and the model statement gives the local-model condition at the twisted
places in range. \Cref{pf:onto2:hasexp} concludes.
\end{proof}

%----------------------------------------------------------------------
\subsubsection{The numerical and exponent-sum statements}\label{sec:pf:onto2:sum}

\begin{definition}[The parameter of size data]\label{pf:onto2:parof}
For size data $L$ and an index $j$, the \emph{parameter of $L$ at $j$} is $P(L,j)=(k,m,t,f,n,\rank V_j,L.L,L.x,L.b',L.\eta_{\mathrm{sz}},
L.S_c)$, with $k,m,t,f,n$ the ranks of $D$.
\formallabel{GlobalTheory.ExpTie.parOf}
\end{definition}

\begin{proposition}[The numerical statement]\label{pf:onto2:numst}
Closedly over $\sfname{FamHyp}$, eventually at points $p$ satisfying the extra point conditions, for every index $j$ there is an
admissible dictionary at $j$.
\formallabel{GlobalTheory.ExpTie.numSt_alt}
\end{proposition}

\begin{proof}
The following six statements of the form $\mathrm{AtPts}$ hold, for every $j$: $P(L,j)$ satisfies $\mathrm{Hyp}$
(\lref{GlobalTheory.ExpTie.parOf_hyp_alt}) and $\mathrm{Large}$ (\lref{GlobalTheory.T.parameter_large_at_all_points}); there are
$X',A$ with $\mathrm{Rows}(P(L,j),X',A)$ and $\mathrm{Big}(L,x,X',A)$ (\lref{GlobalTheory.T.exists_large_parameter_estimates_at_pts});
$h<g_i$ for both level primes (\lref{GlobalTheory.T.parOf_height_lt_genus_at_pts}); every forced weight $\lambda^{\mathrm{def}}_{j,\tau}$
is a genuine weight of $D$ (\lref{GlobalTheory.T.isLamPrime_at_all_points}); and there are weights $W$ with $\mathrm{Ok}_{P(L,j)}(W)$, each
$\lambda_\tau$ a shift of $\lambda^{\mathrm{def}}_{j,\tau}$ by a constant, $0\le U_\tau(a)\le h_0$, $|T_2(W)|\le|T_2(D)|$ and
$S_p(W)=s_m\cup\{q,d,p_{\mathrm{de}},p'_{\mathrm{de}}\}$ (\lref{GlobalTheory.T.weight_lambda_shift_nonneg'}). Statements of the form
$\mathrm{AtPts}$ are closed under conjunction (take the larger of the two thresholds) and under pointwise implication (keep the
threshold), because closed statements are closed under conjunction and implication (\lref{Formalism.MainNotions.Closed.and},
\lref{Formalism.MainNotions.Closed.imp}). At a point where all six hold, $P=P(L,j)$ and this $W$ form a dictionary: the rank, size and
scale conditions hold by definition of $P(L,j)$, $|\mathrm{Tau}(D)|=f$ by the datum, and the remaining conditions are the six statements.
\end{proof}

\begin{proposition}[The exponent-sum statement]\label{pf:onto2:sumexp}
For every point $p$ of every family, every index $j$, every dictionary $(P,W)$ at $j$, every type $t$ and every place $w$,
$\ontoIILI_w(\ontoIIJA(p_w))$ holds for $(\mathrm{blk}(p,j,t),\ell_v(p),Z(p,j,t))$:
\[
\ell^{\ontoIIJA(\ell)}\cdot\mathrm{cut}_Z\nabla\subseteq\ontoIIloc w{(\mathrm{cut}_Z\Delta)},\qquad \ell=p_w .
\]
In particular the exponent-sum statement holds.
\formallabel{GlobalTheory.ExpTie.lIncl_sumExp}
\formallabel{GlobalTheory.ExpTie.sumExpSt_alt}
\formallabel{GlobalTheory.T.sum_expansion'}
\end{proposition}

\begin{proof}
The block of a point is tied to its weight data: there are an isomorphism $e$ between the coefficient space of the block and the
space of the relevant Hermitian data and, for every class $c$, identifications of the lattices $\nabla(B)_c$ and $\Delta(B)_c$ with the
images of the $\nabla$- and $\Delta$-lattices attached to a Hermitian lattice of $V_j$, the forced weights
$\lambda^{\mathrm{def}}_{j,\tau}$ and the embeddings $\tau$ (the tie of the character family of $p$). By
\lref{GlobalTheory.OntoExponentsLat.exZ_sum_smul_mem_locP_delAtr}, for every place $w$ the element
$\ell^{\sum_\tau\ontoIIexp_\ell(\lambda^{\mathrm{def}}_{j,\tau})}$ maps the $\nabla$-lattice of a Hermitian lattice into the localisation at
$w$ of its $\Delta$-lattice. Transport along $e$ (\lref{GlobalTheory.SumExpAt.locIncl_of_equiv}) and passage to the cut lattices
(\lref{GlobalTheory.SumExpAt.cutLat_locIncl}) give the inclusion with exponent $\sum_\tau\ontoIIexp_\ell(\lambda^{\mathrm{def}}_{j,\tau})$.
Finally the Weyl exponent is invariant under shifting a weight by a constant (\lref{GlobalTheory.SumExpAt.exZ_res_eq_natW}), each
$\lambda_\tau$ is such a shift of $\lambda^{\mathrm{def}}_{j,\tau}$ and $N_r=\rank V_j$ by the dictionary; hence the exponent equals
$\ontoIIJA(\ell)$. The statement holds at every point, so the closed statement holds with arbitrary thresholds.
\end{proof}

%----------------------------------------------------------------------
\subsubsection{Smallness, frames, Galois representations and local models}\label{sec:pf:onto2:four}

In this subsection we are in range at $(p,j,P,W,w)$ (\cref{pf:onto2:dictdef}), we write $\ell=p_w$, $B=\mathrm{blk}(p,j,\mu(w))$,
$Z=Z(p,j,\mu(w))$, $E$ for the host CM field and $\mathrm{Tau}=\mathrm{Tau}(D)$.

\begin{proposition}[Smallness]\label{pf:onto2:small}
In range, $\ell$ is a prime with $N_r+1<\ell$ and $\ell\nmid\mathrm{disc}(E)$, the subspace $Z$ is stable under all Hecke operators of
$B$, and $Z$ lies in the block part of $B$.
\formallabel{GlobalTheory.ExpTie.smallSt_alt}
\formallabel{GlobalTheory.T.small_estimate}
\end{proposition}

\begin{proof}
$\ell$ is prime as a residue characteristic. Since $\ell$ is in range, $l_J<\ell$, and for parameters with $\mathrm{Hyp}$ and $\mathrm{Large}$
every prime $\ell>l_J$ satisfies $N_r+1<\ell$ (\lref{Formalism.CohomologicalDescent.LogSmallConstantBounds.paramsLargePrime_holds}, an
explicit comparison of $l_J$ with $N_r$). Since $\ell\notin S_p$, which contains the primes ramified in the host
(\lref{GlobalTheory.ExpTie.not_dvd_discr_host}), $\ell\nmid\mathrm{disc}(E)$. The point conditions contain the primitive row conditions
and the separation of types, from which the complex data of the point satisfy the stability and block conditions at the base level
(\lref{Formalism.IndefR.Links.complex_data_conditions_base_level_of_primitive_links}); these pass to the subspaces $Z(p,j,t)$ of the
reduction and to all indices $j$, directly from the definitions of these subspaces as restrictions of the subspaces of the
point. The conditions do not depend on the
family or its constants, so the closed statement holds trivially.
\end{proof}

\begin{definition}[Frames, readings and ties]\label{pf:onto2:framedef}
Let $\ell$ be a prime, $\iota$ a finite set and $N\ge1$. An \emph{$\ell$-adic frame} $F$ (\lref{GlobalTheory.ExpFrame.Frame}) consists
of a discrete valuation ring $\OO_F$ of characteristic $0$ with a ring map $\OO_F\to\Qlbar$ (completed), a group $G$ with subgroups $\Gamma$
and $K$, homomorphisms $\kappa:K\to\GL_N(\OO_F)^\iota$ and $\chi:K\to\OO_F^\times$, and an index set $I$ of Hecke operators with finite
sets $\mathrm{tr}(i)\subseteq G$ and constants. It is \emph{admissible} if $\OO_F$ maps isomorphically onto the integers of a finite
extension of $\Q_\ell$, $G=\Gamma S K$ for a finite set $S$, $\ell$ does not divide the (finite) orders of the groups
$\Gamma\cap gKg^{-1}$, and the Hecke operators satisfy permutation and commutation conditions
(\lref{GlobalTheory.ExpFrame.Frame.Ok}). For a weight $\mu$ on $\iota$, \emph{$F$-forms of weight $\mu$} are functions
$f:G\to\nabla_{\Qlbar}(\mu)=\bigotimes_\tau\nabla(\mu_\tau)$, left $\Gamma$-invariant and right $K$-equivariant through $\chi$ and $\kappa$;
an \emph{eigensystem} $\theta:I\to\Qlbar$ of weight $\mu$ is the system of Hecke eigenvalues of a non-zero $F$-form. A \emph{reading}
$R$ of a block $B$ at a place $w$ through $F$ is an embedding $e_0:K\to\Qlbar$ inducing $w$
with $e_0(\OO_K)\subseteq\OO_F$, a map from $I$ to the Hecke operators of $B$, and an $e_0$-semilinear map from the forms of $B$ to the
$F$-forms of weight $\lambda$, preserving linear independence, compatible with the Hecke operators, mapping $\nabla(B)$ into the
$\nabla$-lattice, and such that forms mapping into the $\Delta$-lattice lie in $\Delta(B)$ up to a unit at $w$. A \emph{tie} of $(F,R)$ to
the point $p$ (\lref{GlobalTheory.ExpTie.Tied}) identifies $G$ with the finite-adelic unitary group of $V_j$, $\Gamma$ with its
rational points, $K$ with the level group of $p$, $\kappa$ with the $\ell$-adic images of the level matrices under embeddings
$\mathrm{inc}(\tau):E\to\Qlbar$ forming a \emph{half system} (one of each pair of conjugate embeddings), $\chi$ with the level character
twisted by determinants, and the Hecke operators of $F$ with those of the block at the places prime to $\ell$.
\formallabel{GlobalTheory.ExpFrame.Reads}
\end{definition}

\begin{proposition}[The frame statement]\label{pf:onto2:frame}
In range, there are an admissible $\ell$-adic frame $F$ with $N=N_r$, a reading $R$ of $B$ at $w$ through $F$ and a half system
$\mathrm{inc}$ such that $(F,R)$ is tied to $p$.
\formallabel{GlobalTheory.FrameOf.frameSt_of_smallSt'}
\formallabel{GlobalTheory.T.frame_statement_of_small_statement}
\end{proposition}

\begin{proof}
By \lref{GlobalTheory.FrameOf.frameAt_of}, the conclusion follows at $(p,j,P,W,w)$ from $N_r=\rank V_j$ and seven local statements:
the embeddings $e_0\circ\mathrm{lab}$ and their conjugates form half systems; the $\ell$-adic level matrices are integral over the field
generated by $e_0$; the level character has finite order; tameness ($\ell\nmid|\Gamma\cap gKg^{-1}|$); the permutation and commutation
properties of the Hecke double cosets; and the existence of a reading. Each of the seven statements holds in range whenever the
conditions of \cref{pf:onto2:small} hold (\lref{GlobalTheory.T.inRange_halfAt_of_smallAt},
\lref{GlobalTheory.T.inRange_kapAt_of_smallAt}, \lref{GlobalTheory.T.chiAt_of_smallAt},
\lref{GlobalTheory.T.tame_at_of_small_at}, \lref{GlobalTheory.T.permAt_of_smallAt},
\lref{GlobalTheory.T.commAt_of_smallAt_in_range}, \lref{GlobalTheory.T.readsAt_of_smallAt}); combining them with
\cref{pf:onto2:small} (conditions holding in range are closed under conjunction, by taking the larger threshold) and applying the
combination pointwise to every dictionary gives the claim.
\end{proof}

\begin{definition}[Crystalline Galois sides]\label{pf:onto2:galside}
Let $E$ be a CM field, $\ell$ a prime, $\mathrm{inc}$ a half system of embeddings $E\to\Qlbar$, $(P,W)$ a dictionary with
$h_0=\lfloor h\rfloor$, $F$ an $\ell$-adic frame with $N=N_r$, and $\mathcal B$ a property of eigensystems. A Galois representation
$r:G_E\to\GL_n(\Qlbar)$ is \emph{crystalline with jumps} $J(\tau)$ ($\tau:E\to\Qlbar$) if it is continuous and at every place above
$\ell$ its restriction to the decomposition group is crystalline with Hodge--Tate multiset $J(\tau)$ at every embedding $\tau$.
The \emph{Galois side} $\mathrm{Gal}(\mathrm{inc},P,W,F,\mathcal B)$ asserts the existence of labels $\mathrm{lab}(Q)$ (places of $E$), positions of the Hecke
operators, a continuous character $\mathrm{tw}$ and $\epsilon\in\{\mathrm{id},c\}$ such that: (i) for every dominant $\mu$ and every
eigensystem $\theta$ of weight $\mu$ there is $r:G_E\to\GL_{N_r}(\Qlbar)$, crystalline with jumps $J$ satisfying
$J(\epsilon\,\mathrm{inc}\,\tau,a)=U(\mu_\tau,a)$ and $J(\epsilon\,c\,\mathrm{inc}\,\tau,a)=h_0-U(\mu_\tau,N_r-1-a)$, whose characteristic
polynomials at Frobenius elements of the places $\mathrm{lab}(Q)$, twisted by $\mathrm{tw}$, are the Hecke polynomials of $\theta$; (ii)
for every eigensystem $\theta$ of weight $\lambda$ with $\mathcal B(\theta)$ there are $a+t=N_r$, a continuous $\rho:G_E\to\GL_a(\Qlbar)$
and characters $\chi_0,\dots,\chi_{t-1}$, $\chi_s$ crystalline with jump $\mathrm{bot}(\tau)+s$, where
$\mathrm{bot}(\epsilon\,\mathrm{inc}\,\tau)=\mathrm{bot}_W(\tau)$ and $\mathrm{bot}(\epsilon\,c\,\mathrm{inc}\,\tau)=h_0-\mathrm{bot}_W(\tau)-t+1$
($\mathrm{bot}_W(\tau)=U_\tau(k-2+j_\tau+t)$ the bottom of the block of $\lambda_\tau$), such that the Hecke polynomials of $\theta$ are, in
the same sense, the characteristic polynomials of $\rho\oplus\bigoplus_s\chi_s$.
\formallabel{GlobalTheory.ExpFrame.CrysWith}
\formallabel{GlobalTheory.ExpFrame.GalSide}
\end{definition}

\begin{proposition}[Galois representations for definite unitary groups]\label{pf:onto2:galrep}
Assume (H1)--(H8) of \cref{pf:onto2:thmB}. Let $E$ be a CM field with $E/E^+$ unramified at all finite places, $n\ge1$, $\pi$ an
automorphic representation of the definite unitary group $\Uni(n)$ of $E/E^+$ whose archimedean weights are dominant weights
$\lambda_\tau$ ($\tau:E\to\C$), $\ell$ a prime and $\iota:\Qlbar\cong\C$, and assume that $\pi$ is unramified at the places above $\ell$.
Then:
\begin{enumerate}[label=(\roman*)]
\item there is $r:G_E\to\GL_n(\Qlbar)$, crystalline with jumps $J_{\iota,\lambda}(\tau)=\{k-\lambda_{\iota\tau}(k):k<n\}$, such that at every
  split place $w\nmid\ell$ where $\pi$ is unramified, with Satake multiset $T_w$, and every Frobenius element $g$ at $w$,
  $\det(X-r(g^{-1}))=\prod_{z\in T_w}\bigl(X-\iota^{-1}(z\,\Nm w^{(n-1)/2})\bigr)$;
\item if $n<2t$, every $\lambda_\tau$ has a run of length $t$ (a maximal block of $t$ equal entries) and $\pi$ has a \emph{string} of
  length $t$ (at some split place $w\nmid\ell$ where $\pi$ is unramified, the Satake multiset contains $\{\nu\Nm w^{(t-1)/2-c}:c<t\}$
  with $|\nu|=1$), then there are a continuous $\rho:G_E\to\GL_{n-t}(\Qlbar)$ and characters $\chi_0,\dots,\chi_{t-1}$, $\chi_s$ crystalline
  with jump $\mathrm{bot}(\tau)+s$ where $\mathrm{bot}(\tau)=p-\lambda_{\iota\tau}(p)$ for any run position $p$ of $\lambda_{\iota\tau}$, such
  that the identity of (i) holds with $r$ replaced by $\rho\oplus\bigoplus_s\chi_s$.
\end{enumerate}
\formallabel{UpperTheory.GalN.galoisRep_definiteUnitary_alt}
\end{proposition}

\begin{proof}
\emph{Base change.} By (H5) in the form \lref{UpperTheory.BaseChangeIdentities.base_change_speh_narrow_swap}, there are integers
$a_i,d_i\ge1$ ($i<c$) with $\sum_id_ia_i=n$ and conjugate self-dual cuspidal representations $\Pi_i$ of $\GL_{a_i}(\A_E)$ with infinity
types $m_i(\tau)$ (multisets of distinct numbers), such that the local components of $\pi$ at split places are the products of the
Speh components of the $\Pi_i$ of lengths $d_i$, and $\sum_i\mathrm{Speh}_{d_i}(m_i(\tau))$ is the infinitesimal character of $\lambda_\tau$.

\emph{Galois representations of the pieces.} For every $i$ and $k<d_i$, \lref{GlobalTheory.CsdTwistGal.exists_galRep_csd_twist_alt_alt}
(under (H2)--(H4) and (H6)--(H8): Galois representations attached to the conjugate self-dual cuspidal representation $\Pi_i$,
twisted by a power of the norm) gives a continuous $r_{i,k}:G_E\to\GL_{a_i}(\Qlbar)$
which is crystalline with jumps $J_{i,k}(\tau)=\{f_{i,k}-z:z\in m_i(\iota\tau)\}$, $f_{i,k}=\tfrac{n-1}2+\tfrac{d_i-1}2-k$, as soon as
$\Pi_i$ is spherical above $\ell$, and whose Frobenius polynomials at split places prime to $\ell$ are given by the Satake parameters
of $\Pi_i$ scaled by $\Nm w^{f_{i,k}}$. The sphericity holds because $\pi$ is unramified above $\ell$
(\lref{GlobalTheory.SphBC.isSpherical_of_isUnramU}).

\emph{Assembly.} The direct sum $R=\bigoplus_{i,k}r_{i,k}$ is continuous, crystalline with jumps $\sum_{i,k}J_{i,k}$, and has
$\det(X-R(g))=\prod_{i,k}\det(X-r_{i,k}(g))$ (\lref{GlobalTheory.BlockSum.main}: a block-diagonal sum of continuous
representations is continuous, its characteristic polynomials are the products, and it is crystalline with the summed jumps when the
summands are). The sum of the jumps is $J_{\iota,\lambda}$ by the identity of infinitesimal characters, and the Satake
parameters of $\pi$ at a split place are the regrouped scaled Satake parameters of the $\Pi_i$
(\lref{GlobalTheory.satake_linear_factors_regroup}); this gives (i).

\emph{Strings.} For (ii), let $w$ be a place carrying a string of length $t$. The Satake parameters of the conjugate self-dual cuspidal
pieces at $w$ have absolute value $<\Nm w^{1/2}$ (\lref{UpperTheory.GalN.satake_window}), so by
\lref{UpperTheory.GalN.string_forces_long_part} (a string of length $t$ in a multiset made of segments of lengths $d_i$ centred at
parameters of absolute value $<\Nm w^{1/2}$ forces some $d_{i_0}\ge t$) there is $i_0$ with $d_{i_0}\ge t$. Since $n<2t$ and the weights
have runs of length $t$, \lref{GlobalTheory.GalDef.spehPart_is_run} gives $d_{i_0}=t$, $d_i<t$ for $i\ne i_0$, and $m_{i_0}(\tau)$ a
single number determined by the run; hence $a_{i_0}=1$. Put $\rho=\bigoplus_{i\ne i_0,k}r_{i,k}$ (\lref{GlobalTheory.BlockSum.main})
and $\chi_s=\det r_{i_0,t-1-s}$. These characters are crystalline with jumps $\mathrm{bot}(\tau)+s$
(as $a_{i_0}=1$, the determinant of $r_{i_0,t-1-s}$ is this representation itself, whose single jump
$f_{i_0,t-1-s}-m_{i_0}(\iota\tau)$ equals $\mathrm{bot}(\tau)+s$ by the value of $m_{i_0}$ given by the run), $\mathrm{bot}$ does not depend on the chosen run position, $(n-t)+t=n$, and the product formula for $\det(X-R(g))$ gives the
identity in (ii).
\end{proof}

\begin{proposition}[The Galois statement]\label{pf:onto2:galois}
Assume (H1)--(H8) of \cref{pf:onto2:thmB} and the smallness statement. In range, for every admissible $\ell$-adic frame $F$ with $N=N_r$,
every reading $R$ of $B$ at $w$ and every half system $\mathrm{inc}$ such that $(F,R)$ is tied to $p$, the Galois side
$\mathrm{Gal}(\mathrm{inc},P,W,F,\mathcal B_R)$ holds, where $\mathcal B_R(\theta)$ means that $\theta$ is read from a system of eigenvalues
of $B$ of block shape.
\formallabel{GlobalTheory.FrameGal.galois_steinberg_local_alt}
\formallabel{GlobalTheory.T.galois_frame_compatibility_alt}
\end{proposition}

\begin{proof}
\emph{Step 1 (automorphy).} By \lref{GlobalTheory.T.frame_aut_of_small_estimate}, the smallness statement implies, in range: the
places $Q$ of the block are labelled by split places $\mathrm{lab}(Q)$ of $E$ with $\Nm\,\mathrm{lab}(Q)=q_Q$, and the labels of the
places prime to $\ell$ include all primes of $E$ of prime norm outside a finite set, the block parameters $\nu$ satisfy
$|\nu|^2=q_Q^{N_r-t}$; and for every tied frame and reading, every dominant $\mu$ and every eigensystem $\theta$ of weight $\mu$, there is
an automorphic representation $\pi$ of the definite unitary group of $V_j$ with weights $\mu_\tau+\mathrm{sh}(\tau)$ at the embeddings
$\iota\circ\mathrm{inc}(\tau)$ and the dual weights at their conjugates, unramified above $\ell$ and at the places $\mathrm{lab}(Q)$, whose
Satake parameters at $\mathrm{lab}(Q)$ give the eigenvalues $\theta$.

\emph{Step 2 (a crystalline twist).} Put $c(\mathrm{inc}\,\tau)=h_0-N_r+1+\mathrm{sh}(\tau)$ and $c(c\,\mathrm{inc}\,\tau)=-\mathrm{sh}(\tau)$;
then $c(\sigma)+c(c\sigma)=h_0-N_r+1$ for every embedding $\sigma$. There is an algebraic Hecke character of $E$ with infinity type $-c$,
unramified above $\ell$ (\lref{GlobalTheory.AlgCharUnr.exists_algChar_unr}), and by (H2) and (H7) its $\ell$-adic realisation
$\mathrm{tw}$ is continuous and crystalline with jump $c(\sigma)$ at $\sigma$ (\lref{GlobalTheory.FrameGal.exists_crys_char}).

\emph{Step 3 (the Galois side).} Apply \cref{pf:onto2:galrep} to $E$ (unramified over $E^+$ at finite places,
\lref{GlobalTheory.DefiniteBC.host_unramified}), $n=N_r$ and the representations $\pi$ of Step~1. For (i): given $\mu,\theta$, the
twist $r\otimes\mathrm{tw}$ is crystalline with jumps $J$ as required (the jumps of $r$ are $k-\lambda(k)$, those of $\mathrm{tw}$ are $c$,
and the identities $J(c\,\mathrm{inc}\,\tau,a)=U(\mu_\tau,a)$, $J(\mathrm{inc}\,\tau,a)=h_0-U(\mu_\tau,N_r-1-a)$ for $\epsilon=c$ are a direct
computation), and its Frobenius characteristic polynomials read $\theta$ (\lref{GlobalTheory.FrameGal.SystemOf.reads}). For (ii): let
$\theta$ be an eigensystem of weight $\lambda$ of block shape. The weights of $\pi$ have runs of length $t$, coming from the constant
block of $\lambda_\tau$ (\lref{GlobalTheory.FrameGal.isRunAt_block}). By (H3) the set of primes of $E$ of degree one outside a
finite set is non-empty, so some label $\mathrm{lab}(Q)$ is such a prime; at this place the Hecke polynomial of $\theta$ is divisible
by a block polynomial with a root $\nu$ of absolute value $\Nm(\mathrm{lab}\,Q)^{(N_r-t)/2}$ (\lref{GlobalTheory.FrameGal.block_dvd}),
which gives a string of length $t$ for $\pi$ (\lref{GlobalTheory.FrameGal.string_of_dvd}). Since $N_r<2t$ by $\mathrm{Hyp}(P)$, part
(ii) of \cref{pf:onto2:galrep} gives $\rho$ and $\chi_s$, and twisting by $\mathrm{tw}$ yields the required data
(\lref{GlobalTheory.FrameGal.galois_side_of_compatible_systems_alt},
\lref{GlobalTheory.FrameGal.galois_at_place_of_automorphic_at_place_alt}).
\end{proof}

\begin{definition}[Hida settings and local models]\label{pf:onto2:hida}
A \emph{Hida setting} $X$ for $(\iota,N,\lambda)$ consists of
a noetherian local domain $\OO$ of characteristic $0$, an $\OO$-algebra $A$ acting on $\nabla_\OO(\lambda)=\bigotimes_\tau\nabla_\OO(\lambda_\tau)$,
a group $G$ with subgroups $\Gamma,K$ and a homomorphism $j_K:K\to A$, a commutative $\OO$-algebra $T$ of Hecke operators acting on the
$\nabla_\OO(\lambda)$-valued forms on $\Gamma\backslash G$ of type $j_K$, and a multiplicative set $S\subseteq T$. Its \emph{local
inclusion with exponent $J$} says that after localisation at $S$,
$\ell^J$ maps the $\nabla$-valued forms into the forms with values in $\bigotimes_\tau\Delta_\OO(\lambda_\tau)$. A family of Hida settings
\emph{covers} $(B,\ell_v,Z,w)$ if $\ontoIILI_w(J)$ holds for every $J$ for which all members have
the local inclusion with exponent $J$. The \emph{local-model condition} at $(B,\ell_v,Z,w)$ for $(P,W)$
asserts the existence of $(\mathrm{fr},\mathrm{cj})$ and a family of Hida settings for
$(\iota,N_r,\lambda)$ covering $(B,\ell_v,Z,w)$ such that in each of them $\ell$ is not a unit and the \emph{link inputs}
hold for the weights obtained from $W$ by replacing
$(\mathrm{fr},\mathrm{cj})$: $\OO/\ell$ is finite, $\ell$ divides a power of an element $\varpi$, there are weight projectors in $A$ and a
family of $A$-modules $L_\mu$ (containing the simple subquotients of $\nabla_\OO(\lambda)/\varpi$) with Hecke-equivariant forms, and every
$L_\mu$ which is a subquotient of $\nabla_\OO(\lambda)/\varpi$ with non-zero localised forms has, at every $\tau\in T_2$, a highest weight
satisfying the \emph{lattice link} (the existence of Frobenius period-change data over a power series ring for that weight).
\formallabel{Formalism.CohomologicalDescent.ShortWeightBounds.HidaSetting}
\formallabel{Formalism.CohomologicalDescent.DepthNonnegative.TopIncl}
\formallabel{GlobalTheory.ExpTie.Covers}
\formallabel{GlobalTheory.ExpTie.ModelAt}
\formallabel{Formalism.CohomologicalDescent.ShortWeightBounds.WInputsLink}
\formallabel{Formalism.CohomologicalDescent.ShortWeightBounds.LatticeLink}
\end{definition}

\begin{proposition}[Lattice inclusions from local models]\label{pf:onto2:lincl}
Let $(B,\ell_v,Z)$ and a place $w$ be as in \cref{pf:onto2:blocks}, with $\ell=p_w$ prime, let $(P,W)$ be admissible with $\ell$ in
range, let $(\mathrm{fr},\mathrm{cj})$ be re-indexing maps and let $(X_i)_{i\in I}$ be a family of Hida settings for $(\iota,N_r,\lambda)$
covering $(B,\ell_v,Z,w)$ such that in every $X_i$ the prime $\ell$ is not a unit and the link inputs hold for the weights $W'$
obtained from $W$ by replacing $(\mathrm{fr},\mathrm{cj})$. Then $\ontoIILI_w(J^{\mathrm c}(\ell))$ holds.
\formallabel{GlobalTheory.ExpTie.lIncl_of_models}
\formallabel{GlobalTheory.T.primePowerLatticeInclusion_of_localModels}
\end{proposition}

\begin{proof}
\emph{Re-indexing.} $\mathrm{Ok}_P(W')$ holds and, $\ell$ being in range, the coarse bounds of $W'$ and $W$ at $\ell$ coincide
(\lref{GlobalTheory.expo_frob_indep}: the conditions $\mathrm{Ok}_P$ do not involve $(\mathrm{fr},\mathrm{cj})$, and for primes in range
the coarse bound is unchanged by the re-indexing).

\emph{Weight inputs.} Fix $i$. Since $l_J<\ell$, the link inputs of $X_i$ imply its \emph{weight inputs}
(\lref{Formalism.CohomologicalDescent.ShortWeightBounds.weight_inputs_of_linkage_holds}): the frame conditions on the weight
projectors and the Hecke operators and the description of the simple subquotients of $\nabla_\OO(\lambda)/\varpi$ are part of the link
inputs, and the remaining \emph{first input of weight elimination} asks that, at every $\tau\in T_2$, the depth of the highest weight
of every simple subquotient with non-zero localised forms be at most the integer $i_W(\ell,\tau)$ attached to the size data of
$(P,W')$; this is deduced from the lattice link at $\tau$ (\lref{Formalism.CohomologicalDescent.ShortWeightBounds.depth_of_linkage_holds},
through the orientation lattices of the Frobenius period-change data,
\lref{Formalism.CohomologicalDescent.WorkedExample.orientation_lattice_holds}).

\emph{The local inclusion.} For a Hida setting with its weight inputs, in which the prime $\ell>l_J$ is not a unit, the local inclusion
holds with exponent $J^{\mathrm W}(\ell)=\sum_\tau c_\tau$, where $c_\tau=\ontoIIexp(\tau,\ell)$ for $\tau\notin T_2$ and $c_\tau$ is, for
$\tau\in T_2$, an explicit floor-type quantity built from the Weyl exponent of the first $k_c$ entries of $\lambda_\tau$, the depth
bound $i_W(\ell,\tau)$ and the block data of $W'$ at $\tau$ (\lref{Formalism.CohomologicalDescent.ShortWeightBounds.Wts.JW}). This is
\lref{Formalism.CohomologicalDescent.ShortWeightBounds.inclusion_w_real_holds}, which derives the inclusion from a shape bound and
the graded form of a Pieri-type lemma on products of column minors
(\lref{Formalism.CohomologicalDescent.Down.hG}, proved from the Pieri rule). Moreover $J^{\mathrm W}(\ell)\le J^{\mathrm c}(\ell)$
(\lref{Formalism.CohomologicalDescent.Down7.JW_le_coarseWeightBound}), and a local inclusion with exponent $J$ implies the one with
exponent $J'\ge J$ (multiply by $\ell^{J'-J}$, which preserves the $\Delta$-valued forms). Hence every $X_i$ has the local inclusion with
exponent $J^{\mathrm c}_{W'}(\ell)=J^{\mathrm c}(\ell)$ (\lref{Formalism.CohomologicalDescent.Top.inclusion_w_seven_real_holds}).

\emph{Covering.} Since the family covers $(B,\ell_v,Z,w)$, the last statement gives $\ontoIILI_w(J^{\mathrm c}(\ell))$.
\end{proof}

\begin{proposition}[Crystalline block lattices over a power series ring]\label{pf:onto2:wach}
Assume (H6). Let $E$ be a CM field, $\ell\ne2$ a prime unramified in $E$, and $\sigma$ a Frobenius permutation of the embeddings
$E\to\Qlbar$ ($|\sigma(\tau)(x)-\tau(x)^\ell|<1$ for $x\in\OO_E$). Let $n,t,h\in\N$, let $J(\tau,a)\in[0,h]$ ($a<n$) and $\mathrm{bot}(\tau)$
with $\mathrm{bot}(\tau)+s\in[0,h]$ for $s<t$, and assume that there is a \emph{crystalline block pair} for $(n,t,J,\mathrm{bot})$:
continuous $r:G_E\to\GL_n(\Qlbar)$, crystalline with jumps $J$, a continuous $\rho:G_E\to\GL_{n-t}(\Qlbar)$ and characters $\chi_s$
crystalline with jumps $\mathrm{bot}+s$ ($s<t$), such that the coefficients of $\det(X-r(g))-\det(X-\rho(g))\prod_s(X-\chi_s(g))$ have
absolute value $<1$ for all $g$ (\lref{GlobalTheory.ExpFrame.BlockPair}). Then there are a domain $\OO'$ over $\Z_\ell$, complete
for an ideal $I$, a field $F$ of characteristic $\ell$, a field $L'$ with $\OO'\to L'$, Frobenius period-change lattice data of
type $(t,n)$ over $(\OO',F)$, a unit $u\in\Z_\ell^\times$, $c\in\Z_\ell$, integers $\kappa(\tau,s)$, $\zeta\in\OO'$ and maps of power series
rings, which satisfy the inputs of the Frobenius period change for $(\sigma,\ell,h,u,c,J,\mathrm{bot}+s,\kappa)$
(\lref{Formalism.CohomologicalDescent.FrobeniusPeriodChange.NInputs}).
\formallabel{GlobalTheory.T.crystalline_block_lattices_over_power_series_ring}
\formallabel{GlobalTheory.WachBlock.wach_inputs}
\end{proposition}

\begin{proof}
The block pair yields integral crystalline block data whose reductions agree
(\lref{GlobalTheory.T.integralCrystallineBlockData_of_reductionPair}). From these and (H6) one obtains Wach-type lattices over a power series ring over $\OO'$, rank-one models for the characters
$\chi_s$ and a reduced \'etale filtration (\lref{GlobalTheory.T.wachLattices_rankOne_models_reducedFiltration}). Adapted bases put the Frobenius and the filtration in block-triangular form
with the prescribed diagonal blocks (\lref{GlobalTheory.T.adaptedBases_blockTriangular_of_reducedEtaleFiltration}). The bounds
$0\le J\le h$ and $0\le\mathrm{bot}+s\le h$, together with these bases and the data of the previous step, are exactly the inputs of the
Frobenius period change.
\end{proof}

\begin{proposition}[Lattice links from crystalline block pairs]\label{pf:onto2:latlink}
Assume (H6). Let $E$ be a CM field, $\ell\ne2$ a prime unramified in $E$, $\iota$ a finite set and $\mathrm{inc}:\iota\to\{E\to\Qlbar\}$
a half system. Then there are re-indexing maps $(\mathrm{fr},\mathrm{cj})$ such that for every parameter $P$ with $\mathrm{Hyp}(P)$, all
weights $W$ with $\mathrm{Ok}_P(W)$ and $0\le U_\tau(a)\le h_0$ for $a<N_r$, every $\mu$ \emph{linked} to $\lambda$ modulo $\ell$ (for
every $\tau$: $\mu_\tau|_{N_r}$ is non-increasing, $\sum_a\mu_\tau(a)=\sum_a\lambda_\tau(a)$, $\lambda_\tau(N_r-1)\le\mu_\tau(a)\le
\lambda_\tau(0)$, and the multisets $\{\lambda_\tau(a)-a\}$ and $\{\mu_\tau(a)-a\}$ coincide modulo $\ell$;
\lref{GlobalTheory.ExpFrame.Linked}) for which there is a \emph{block datum} (\lref{GlobalTheory.ExpFrame.BlockFor}: a
crystalline block pair with $\epsilon\in\{\mathrm{id},c\}$ and jumps $J(\epsilon\,\mathrm{inc}\,\tau,a)=U(\mu_\tau,a)$,
$J(\epsilon\,c\,\mathrm{inc}\,\tau,a)=h_0-U(\mu_\tau,N_r-1-a)$, $\mathrm{bot}(\epsilon\,\mathrm{inc}\,\tau)=\mathrm{bot}_W(\tau)$,
$\mathrm{bot}(\epsilon\,c\,\mathrm{inc}\,\tau)=h_0-\mathrm{bot}_W(\tau)-t+1$), and every $\tau$, the lattice link holds at $\tau$ for $\mu_\tau$
and the weights $W'$ obtained from $W$ by replacing $(\mathrm{fr},\mathrm{cj})$.
\formallabel{GlobalTheory.T.crystallineLatticeLink_of_blockData}
\formallabel{GlobalTheory.ExpLink.latticeLink_of_blockFor}
\end{proposition}

\begin{proof}
Since $\ell$ is unramified in $E$, there is a Frobenius permutation $\sigma$ of the embeddings $E\to\Qlbar$ commuting with complex
conjugation (\lref{GlobalTheory.T.exists_arithFrobenius_embedding_perm_cmConj}). Every embedding is $\mathrm{inc}(\tau)$ or
$c\,\mathrm{inc}(\tau)$ for a unique $\tau$; let $\mathrm{fr}(\tau)$ be the index with $\sigma(\mathrm{inc}\,\tau)\in\{\mathrm{inc}\,\mathrm{fr}(\tau),
c\,\mathrm{inc}\,\mathrm{fr}(\tau)\}$, and $\mathrm{cj}(\tau)$ true if and only if $\sigma(\mathrm{inc}\,\tau)\ne\mathrm{inc}\,\mathrm{fr}(\tau)$.
Then $\sigma(\mathrm{inc}\,\tau)$ is $c\,\mathrm{inc}\,\mathrm{fr}(\tau)$ or $\mathrm{inc}\,\mathrm{fr}(\tau)$ according to $\mathrm{cj}(\tau)$.

Let $P,W,\mu$ and a block datum $(\epsilon,J,\mathrm{bot})$ be given. \emph{Bounds.} By the linkage, $\lambda_\tau(N_r-1)\le\mu_\tau(a)\le
\lambda_\tau(0)$, so $0\le U(\mu_\tau,a)\le h_0$ follows from $0\le U_\tau(a)\le h_0$ at $a=0$ and $a=N_r-1$; hence $0\le J\le h_0$ at all
embeddings (also at the conjugates, where $J=h_0-U(\mu_\tau,\cdot)$). The block of $\lambda_\tau$ is constant
(\lref{GlobalTheory.ExpLink.U_block}), so $\mathrm{bot}_W(\tau)+s$ is a value of $U_\tau$, and $0\le\mathrm{bot}+s\le h_0$ at all
embeddings. \emph{Lattices.} \Cref{pf:onto2:wach} (with $n=N_r$, $h=h_0$ and the block pair of the block datum) gives Frobenius
period-change data with their inputs. \emph{The lattice link.} It remains to provide the labels of the lattice link: on the embeddings
$\mathrm{inc}(\tau')$ and $c\,\mathrm{inc}(\tau')$ the labels take the values $U_{\tau'}$, $\mathrm{bot}_W(\tau')$ and $U(\mu_{\tau'},\cdot)$, and
their conjugates $h_0-(\cdot)$ reversed; with these labels the reflection condition holds because $\sigma$ commutes with complex
conjugation, the linkage condition modulo $\ell$ is the residue condition of the linkage of $\mu$ and $\lambda$ (rewritten with
$U(\nu,a)+1-N_r=\nu(a)-a$), the carrying condition is the description of $J$ and $\mathrm{bot}$ in the block datum, the
$\sigma$-compatibility of $\mathrm{inc}(\tau)$ is the definition of $(\mathrm{fr},\mathrm{cj})$, and the equality of the sums of $U(\mu_\tau,\cdot)$
and $U_\tau$ follows from $\sum_a\mu_\tau(a)=\sum_a\lambda_\tau(a)$.
\end{proof}

\begin{proposition}[Link inputs from the Galois side]\label{pf:onto2:winputs}
Assume (H3). Let $E$ be a CM field, $\ell$ a prime, $\mathrm{inc}:\iota\to\{E\to\Qlbar\}$, re-indexing maps $(\mathrm{fr},\mathrm{cj})$,
a parameter $P$ and weights $W$ with $\mathrm{Ok}_P(W)$ such that the lattice links hold for all $\mu$ linked to $\lambda$ modulo
$\ell$ with a block datum (as in the conclusion of \cref{pf:onto2:latlink}). Let $F$ be an admissible $\ell$-adic frame with $N=N_r$,
$\mathcal B$ a property of eigensystems with the Galois side $\mathrm{Gal}(\mathrm{inc},P,W,F,\mathcal B)$, $j_K$ a level map for $F$
and $\lambda$, and a Hecke module structure on the $\nabla$-valued $F$-forms with the Hecke property; assume $\OO_F/\ell$ finite and $\ell$
not a unit in $\OO_F$. Let $\mathfrak m$ be a maximal ideal of the Hecke ring and $\theta_0$ an eigensystem of weight $\lambda$ with
$\mathcal B(\theta_0)$ lying over $\mathfrak m$. Then the link inputs hold for the Hida setting of $F$ localised at $\mathfrak m$ and the
weights $W'$.
\formallabel{GlobalTheory.ExpModel.winputs_frame}
\end{proposition}

\begin{proof}
Since $\OO_F$ is a discrete valuation ring in which $\ell$ is not a unit, there are an irreducible $\varpi$ and $n_0$ with $\varpi\mid\ell$
and $\ell\mid\varpi^{n_0}$ (\lref{GlobalTheory.ExpModel.exists_unif_dvd}). The algebra generated by the group action and the weight
projectors contains, for each $\tau$ and weight $w$, a unique projector $e(\tau,w)$ onto the weight-$w$ component at $\tau$
(\lref{GlobalTheory.ExpFrame.exists_unique_isProj}); these projectors preserve the chain lattices
(\lref{GlobalTheory.ExpFrame.alg_chainB_stable}), and the Hecke operators preserve the $\nabla$-valued forms
(\lref{GlobalTheory.T.heckeStable_of_heckeFrame_polynomial_action}). By \lref{GlobalTheory.ExpModel.winputs_of_simple} it
remains to show, for every simple subquotient $G$ of $\nabla_{\OO_F}(\lambda)/\varpi$: $G$ carries a Hecke module structure on its forms
compatible with that of $\nabla$; $G$ has a highest weight $\mu$ with $e(\tau,\mu_\tau)G\ne0$ for all $\tau$; and if the forms with
values in $G$ have non-zero localisation at $\mathfrak m$, then the lattice link holds at every $\tau$ for $\mu_\tau$.

The Hecke structure comes from \lref{GlobalTheory.ExpFrame.HeckeMod.hecke_module}. A simple subquotient $G$ embeds into the reduction of
the dual Weyl module of a highest weight $\mu$ with $\lambda_\tau(N_r-1)\le\mu_\tau(a)\le\lambda_\tau(0)$, $\sum_a\mu_\tau(a)=\sum_a\lambda_\tau(a)$
and $\mu_\tau$ non-increasing, and the projector $e(\tau,\mu_\tau)$ does not kill $G$
(\lref{GlobalTheory.T.simpleSubquotient_dualWeyl_embeds_highestWeight}). Assume that the localised forms of $G$ are non-zero
(\lref{GlobalTheory.T.latticeLink_of_surviving_simple_subquotient}). Then: there is an eigensystem $\theta$ of weight $\mu$ in
characteristic $0$ lying over $\mathfrak m$ (\lref{GlobalTheory.T.exists_charZero_eigenform_surviving_compositionFactor}); a Galois
conjugate $\theta'$ of $\theta$ is congruent to $\theta_0$ modulo the maximal ideal
(\lref{GlobalTheory.T.exists_galoisConj_eigenSystem_congruent}); by the Galois side for $\theta'$ and $\theta_0$ and Chebotarev's theorem
(H3), the congruence of Hecke polynomials becomes a congruence of characteristic polynomials of the attached Galois representations,
which is a block datum for $\mu$ (\lref{GlobalTheory.T.blockData_of_congruent_heckeSystem}); and since $G$ is a simple subquotient of
the reduction of $\nabla(\lambda)$ with highest weight $\mu$, the residues of $\mu_\tau(a)-a$ and $\lambda_\tau(a)-a$ modulo $\ell$ agree as
multisets (\lref{GlobalTheory.T.simpleModule_linkage_dualWeylLattice_reduction}), so $\mu$ is linked to $\lambda$ modulo $\ell$. The
hypothesis on lattice links now gives the lattice link at every $\tau$.
\end{proof}

\begin{proposition}[The model statement]\label{pf:onto2:model}
Assume (H3) and (H6) of \cref{pf:onto2:thmB}. If the smallness, frame and Galois statements hold, then the local-model condition holds
in range (for $(B,\ell_v(p),Z,w)$ and $(P,W)$).
\formallabel{GlobalTheory.ExpModel.modelSt_of_small_frame_gal}
\formallabel{GlobalTheory.T.model_statement_of_crystalline_chebotarev_frame_gal}
\end{proposition}

\begin{proof}
Fix $(p,j,P,W,w)$ in range. By \cref{pf:onto2:small} $\ell$ is a prime with $N_r+1<\ell$ (so $\ell\ne2$, as $N_r+1=k+m+t\ge3$),
$\ell\nmid\mathrm{disc}(E)$, and $Z$ is Hecke-stable and contained in the block part. By the frame statement choose an admissible frame $F$,
a reading $R$ and a half system $\mathrm{inc}$ tied to $p$; by the Galois statement the Galois side $\mathrm{Gal}(\mathrm{inc},P,W,F,\mathcal
B_R)$ holds.

\emph{Lattice links.} By \cref{pf:onto2:latlink}, applied to $E$, $\ell\ne2$ (unramified in $E$) and the half system
$\mathrm{inc}$, there are $(\mathrm{fr},\mathrm{cj})$ such that the lattice links hold for all parameters with $\mathrm{Hyp}$, all
$W$ with $\mathrm{Ok}$ and $0\le U_\tau(a)\le h_0$, and all $\mu$ linked to $\lambda$ modulo $\ell$ with a block datum.

\emph{Hida settings.} Choose a level map $j_K$ realising $\chi\cdot\kappa$ in the algebra generated by the group action and the weight
projectors (\lref{GlobalTheory.ExpFrame.levelMap_exists}) and the corresponding Hecke module structure on $F$-forms
(\lref{GlobalTheory.ExpFrame.HeckeMod.hecke_module}). Let $I$ be the set of maximal ideals $\mathfrak m$ of the Hecke polynomial
ring $T=\OO_F[x_i]_{i\in I_F}$ for which some eigensystem $\theta$ occurring in the reading of $Z$ lies over $\mathfrak m$, and let
$X_{\mathfrak m}$ be the Hida setting of $F$ localised at $\mathfrak m$. These settings cover $(B,\ell_v(p),Z,w)$
(\lref{GlobalTheory.T.maximalIdeals_meeting_subspace_cover}: an inclusion of $\nabla$- into $\Delta$-valued $F$-forms after localising at
every such $\mathfrak m$ transports through the reading to the cut lattices of $Z$).

\emph{Link inputs.} For $\mathfrak m\in I$ with an eigensystem $\theta_0$ over $\mathfrak m$ occurring in $Z$: $\ell$ is not a unit in
$\OO_F$ (\lref{GlobalTheory.T.frameCoeffRing_prime_not_isUnit}) and $\OO_F/\ell$ is finite
(\lref{GlobalTheory.T.finite_frameCoeffRing_quotient_prime}); since $Z$ lies in the block part and the Hecke operators of degree $0$ of
$B$ are trivial, $\theta_0$ has block shape and is an eigensystem of weight $\lambda$
(\lref{GlobalTheory.T.hasBlocks_of_system_occurs_in_block_part}). Then \cref{pf:onto2:winputs} (which uses (H3), the Galois side and the lattice links; the
hypotheses $\mathrm{Ok}_P(W)$ and $0\le U_\tau(a)\le h_0$ are part of the dictionary) gives the link inputs for $X_{\mathfrak m}$.
\end{proof}

%======================================================================
\subsection{Remarks on the formal statements}\label{sec:pf:onto2:remarks}

\begin{remark}[Lengths, valuations and minima]\label{pf:onto2:rem:len}
In the formal development the length of a module is a natural number, a module of infinite length being given length $0$, and
the valuations are extended by $v_w(0)=0$. Thus \eqref{eq:pf:onto2:ix} is defined for arbitrary submodules; for lattices in a
common subspace all lengths in it are finite and $\ontoIIix_w$ is the usual local index. In \cref{pf:onto2:l5} the module $T_h$ is
arbitrary, and the inequality $\ontoIIix_v(A,B)\le\len(\ontoIIloc vA/\ontoIIloc v{(A\cap B)})$ used in its proof rests on the
non-negativity of lengths under this convention. Likewise the Weyl exponent of \cref{pf:onto2:weylexp} and the exponent $J_W$ of
\cref{pf:onto2:blocks} are defined to be $0$ when no exponent with the required property exists; the upper bounds of
\cref{pf:onto2:jexpU,pf:onto2:jrec} are deduced from inclusions which are proved, so they are bounds for genuine exponents.
\end{remark}

\begin{remark}[Hypotheses which are not used]\label{pf:onto2:rem:unused}
(i) Clozel's purity theorem (H1) is a hypothesis of \cref{pf:onto2:galrep,pf:onto2:galois,pf:onto2:thmB}, but it is not used:
the formal proof of \cref{pf:onto2:galrep} does not use it, and it enters the proofs of \cref{pf:onto2:galois,pf:onto2:thmB} only
as a hypothesis passed to \cref{pf:onto2:galrep}. Hence the exponent statement follows from (H2)--(H8).
(ii) The lemma on direct sums \lref{GlobalTheory.BlockSum.main} used in the proof of \cref{pf:onto2:galrep} carries the
Bhatt--Scholze theorem (H6) as a hypothesis which is not used in its proof; the direct sum of crystalline representations is
treated directly. The theorem (H6) enters \cref{pf:onto2:model} through the lattice links
(\cref{pf:onto2:latlink}).
(iii) For the hypotheses of \cref{pf:onto2:latleaves} which are not used, see \cref{pf:onto2:rem:latunused}.
\end{remark}

\begin{remark}[Existentially quantified data]\label{pf:onto2:rem:exist}
The leaves \ref{pf:onto2:L5} and \ref{pf:onto2:L10} assert the existence of data $T_{h1},a,r_1,r_2,\delta_{e2}$, respectively of a
function $J_p$, with the stated properties; they do not involve the corresponding components of the record object $O$. In the
proofs above these data are $T_{h1}=T_h$, the remainders $r_1,r_2$ constructed in \cref{pf:onto2:index}, the thick density of $L$,
and $J_p(j,\ell)=\ontoIIJr_j(\ell)$ (\cref{pf:onto2:hasexp}).
\end{remark}

\begin{remark}[Inputs proved in the development]\label{pf:onto2:rem:proved}
\Cref{pf:onto2:thmA} has no hypotheses besides the page statement: the formal-group statements $(\mathrm A_p)$, $(\mathrm H_p)$ and
$(\mathrm U_p)$ used for the depth leaves are proved in the development for every prime $p$
(\lref{GlobalTheory.ChF.colength_bound_split_place_chart}, \lref{GlobalTheory.ChF.hodgeTate_inclusion_chart},
\lref{GlobalTheory.ChF.unramified_companion_chart_characters}). Similarly, the prime number bound $\pi(y)<1.26\,y/\log y$ used in
\cref{pf:onto2:sumlog} is proved in the development (\lref{Formalism.CohomologicalDescent.Top.prime_count_bound_rs}) and is not
a hypothesis.
\end{remark}

\begin{remark}[The functions $\theta$ and $\alpha$]\label{pf:onto2:rem:alpha}
For the record objects considered here the depth tie gives $\alpha=0$ and $\theta=\mathrm{th}_k$ (see the proofs of
\cref{pf:onto2:leafgood,pf:onto2:leafgoodrun}). Hence the conditions on $\alpha$ in \ref{pf:onto2:L6} hold trivially, and the
condition on $\theta$ is the bound $\sum_{p\in Q}\mathrm{th}_k(p)\log p\le20/k$.
\end{remark}

\begin{remark}[Closed statements]\label{pf:onto2:rem:closed}
\Cref{pf:onto2:thmA,pf:onto2:thmB} are implications between closed statements over $\sfname{FamHyp}$ (\cref{def:families:closed}):
they concern only families of the class $\sfname{FamHyp}$ and, for each such family, only sufficiently large constants and scales.
In \cref{pf:onto2:latleaves} the threshold is $x_1=0$, and the statements of \cref{pf:onto2:sumexp,pf:onto2:small} hold at every
point, so that the thresholds play no role there.
\end{remark}
